% 本文件由 LaTeX 排版 Agent 根据有效 translation.json 自动生成。
% author=Augustine of Hippo; work_id=1601; title=论山中圣训

\chapter{论山中圣训}

% structure: p0001 source=h1 target=omitted reason=page-title-already-rendered-by-parent

\Emph{卷一} 论山中圣训第一部分（\BibleRef{玛 5}）\newline{} \Emph{卷二} 论山中圣训第二部分（\BibleRef{玛 6-7}）\par

\SourceRecord{Translated by William Findlay.；Nicene and Post-Nicene Fathers, First Series；Vol. 6.；Edited by Philip Schaff.；Buffalo, NY: Christian Literature Publishing Co.,；1888.；<http://www.newadvent.org/fathers/1601.htm>.}

% structure: p0001 source=h1 target=section reason=page-title

\section{论山中圣训 第一册}

\BibleBook{玛窦}{福音} 5\par

% structure: p0003 source=h2 target=subsection reason=relative-heading-rank; base=h2

\subsection{第一章}

1. 若有人虔敬而审慎地思考我们的主耶稣基督在山上所讲的那篇训言——正如我们在 \BibleBook{玛窦}{福音} 中所读到的——我认为，按最高道德标准而言，他将在其中发现基督徒生活的完美规范。这并非我们轻率地妄加断言，而是从主自己的话中得出的结论。因为这篇训言本身以这样的方式结束，使人清楚看到其中包含塑造生命的所有诫命。主如此说：\BibleQuote{所以，凡听见我这些话而实行的，我就把他比作一个聪明人，他把房子盖在磐石上；雨淋、水冲、风吹，袭击那座房子，它却不坍塌，因为基础建在磐石上。凡听见我这些话而不实行的，就像个愚昧人，他把房子盖在沙土上；雨淋、水冲、风吹，袭击那座房子，它就坍塌了，且倒塌得很厉害。} 因此，主并非简单地说\BibleQuote{凡听见我话的}，而是加上\BibleQuote{凡听见我这些话的}，我认为这充分表明，祂在山上所说的这些话如此完美地指导那些愿意按之生活的人，使他们可以正当地被比作在磐石上盖房子的人。我这样说，只是要表明我们面前的这篇训言在塑造基督徒生活的一切诫命上是完美的；至于这一部分的具体内容，将在适当的地方作更详细的论述。\par

2. 这篇训言的开头是这样写的：\BibleQuote{耶稣看见这许多群众，就上了山；坐下后，祂的门徒来到祂跟前，祂就开口教训他们说：} 若问这座\Quote{山}有何含义，完全可以理解为更大的义德诫命，因为曾有一些较小的诫命赐给了犹太人。然而，同一天主，通过祂的圣先知和仆人，按照完全安排好的时期分赐，将较小的诫命赐给一个尚需被恐惧约束的民族，又通过祂的圣子，将更大的诫命赐给一个如今已适合被爱释放的民族。此外，当小的赐给小的，大的赐给大的时，都是出于那唯一知道如何适时为人类提供合适药物的天主。并不奇怪，更大的诫命是为天国而赐，较小的诫命是为地上的国而赐，两者都出自那创造了天地的同一天主。因此，关于那更大的义德，先知说：\BibleQuote{你的义德如同天主的山岳。} 这可以理解为，只有那位唯一适合教导如此重大事情的导师，在山上施教。祂坐着施教，这符合导师职位的尊严；门徒来到祂跟前，不仅在身体上更近以便听祂的话，也在精神上接近以履行祂的诫命。\BibleQuote{祂就开口教训他们说：} 这个迂回的说法\Quote{祂就开口}，或许优雅地通过这一停顿暗示这篇训言将比平时稍长；除非，也许这并非没有意义：如今祂被说成\Quote{开了祂自己的口}，而在旧约之下，祂通常开启先知的口。\par

3. 那么，祂说了什么？\BibleQuote{神贫的人是有福的，因为天国是他们的。} 我们在圣经中读到，关于追求世俗之事：\BibleQuote{一切都是虚幻，是精神的夸耀；} 但精神的夸耀指的是傲慢和骄傲；通常骄傲的人也被说成有\Quote{大精神}，这正是因为风也被称为\Quote{精神}。因此也记载着：\BibleQuote{火、冰雹、雪、冰、暴风的精神。} 然而，有谁不知道骄傲的人被说成是\Quote{膨胀的}，仿佛被风吹胀？所以宗徒也说：\BibleQuote{知识使人膨胀，爱德却建造人。} 这里的\Quote{神贫的人}正确地被理解为谦卑和敬畏天主的人，即没有那种使人膨胀的精神的人。有福的开端不能从任何其他点开始，如果它确实要达到最高智慧的话；\BibleQuote{敬畏主是智慧的开端；} 另一方面，\Quote{骄傲}也被称为\BibleQuote{一切罪恶的开端。} 因此，让骄傲的人追求和喜爱地上的国；但\BibleQuote{神贫的人是有福的，因为天国是他们的。}\par

% structure: p0007 source=h2 target=subsection reason=relative-heading-rank; base=h2

\subsection{第二章}

4. \BibleQuote{温良的人是有福的，因为他们要继承土地：} 我猜想，这土地就是在 \BibleBook{圣}{咏} 中所说的：\BibleQuote{你是我的避难所，我在活人之地的福分。} 因为它象征着永恒产业的一种稳固和坚定，在那里，灵魂因着良好的性情，仿佛安息在自己的处所，正如身体安息在土地上，并从土地中以其食物得滋养，如同身体从土地得滋养。这就是圣徒的安息与生命。温良的人是那些向恶行屈服，不抵抗邪恶，而以善胜恶的人。因此，让不温良的人为地上的、暂时的事物争吵争斗；但\BibleQuote{温良的人是有福的，因为他们要继承土地，} 他们不能被赶出这块土地。\par

5. \BibleQuote{哀恸的人是有福的，因为他们要受安慰。}哀恸是因失去所爱之物而产生的忧伤；但那些归向天主的人，失去了他们从前在世上所习惯视如珍宝的东西：因为他们不再为从前所喜乐的事物而喜乐；直到永恒之爱在他们内形成之前，他们总被某种程度的忧伤所刺伤。因此，他们必由圣神获得安慰，圣神为此特别被称为\OriginalTerm{护慰者}{Paraclete}，\Emph{即}安慰者，为使他们失去暂时之乐时，能尽情享受永恒之乐。\par

6. \BibleQuote{饥渴慕义的人是有福的，因为他们要得饱饫。}祂称这些人为真实不朽之善的爱好者。因此，他们将以那食物——主亲自说：\BibleQuote{我的食物就是承行派遣我者的旨意}，这旨意即是义——以及那水——正如祂也说：凡喝这水的，\BibleQuote{在他内要成为涌到永生的水泉}——而得饱饫。\par

7. \BibleQuote{怜悯人的人是有福的，因为他们要受怜悯。}祂说，那些同情受苦者的人是有福的，因为这样他们所得的回报，就是他们自己摆脱了苦难。\par

8. \BibleQuote{心里洁净的人是有福的，因为他们要看见天主。}因此，那些用外在肉眼寻求天主的人是多么愚昧，因为祂是用心来看见的！正如别处所载：\BibleQuote{怀着正直诚朴的心寻求他。}因为纯洁的心就是单一的心；正如这光只能以清洁的眼睛看见，同样，除非那能看见天主的官能是纯洁的，否则就看不见天主。\par

9. \BibleQuote{缔造和平的人是有福的，因为他们要称为天主的子女。}和平的完满在于没有任何对立；而天主的子女是缔造和平的，因为没有什么能抵抗天主，子女当然应有他们父亲的相似。现在，那些在自己内将灵魂的一切情感纳入秩序、使之服从理性——\Emph{即}心智和神性——并彻底制服其肉体私欲的人，他们自身就成了天主的国：在这国中，一切都安排得如此，使人身上那首要和卓越的部分毫无抵抗地统治着那与我们同野兽共有的其他元素；而那人身上卓越的部分，\Emph{即}心智和理性，又被置于更优越者即真理本身、天主的独生子之下。因为人除非先服从那更高的，就不能统治那较低的。这就是赐给地上善意之人的和平；这是成全而完美智者的生命。在这种完全和平与秩序的国度中，这世界的首领被驱逐出去，他统治着乖谬无序之处。当这和平在内里建立并巩固之后，那被逐者从外部挑起的任何迫害，只会增加那属于天主的荣耀；他无法动摇那建筑中的任何部分，反而因他阴谋的失败，显明那建筑从内到外是用多大的力量建造的。因此接着便说：\BibleQuote{为义而受迫害的人是有福的，因为天国是他们的。}\par

% structure: p0014 source=h2 target=subsection reason=relative-heading-rank; base=h2

\subsection{第三章}

10. 因此，总共有这八句话。因为在余下的部分，祂以直接称呼在场者的方式说话，说：\BibleQuote{几时人辱骂你们，迫害你们，你们是有福的。}但先前的话祂是以一般方式说的：因为祂没有说：你们神贫的人是有福的，因为天国是你们的；而是说：\BibleQuote{神贫的人是有福的，因为天国是他们的。}也不说：你们温良的人是有福的，因为你们要承受土地；而是说：\BibleQuote{温良的人是有福的，因为他们要承受土地。}直到第八句话都是如此，在那里祂说：\BibleQuote{为义而受迫害的人是有福的，因为天国是他们的。}此后祂开始以直接称呼在场者的方式说话，尽管先前所说的话也涉及在场的听众；而后面的、似乎特别对在场者说的话，也涉及那些不在场或以后将出现的人。\par

因此，我们面前的这些真福句的数目应仔细思量。因为真福以谦卑开始：\BibleQuote{神贫的人是有福的\BibleRef{玛 5:3}}，即那些不自高自大的人，而灵魂顺服于天主的权威，唯恐在今生之后走向惩罚，尽管在今生它可能自以为幸福。然后，它（灵魂）开始认识圣经，在此它应表现出温良的虔敬，以免冒然谴责那些对未受教者似乎荒谬的事，并因固执的争辩而自身变得不可教导。此后，它开始知道自己因肉体的习惯和罪恶而被此世的何种纠缠所束缚：于是在这第三阶段——即认知的阶段——为丧失至善而哀恸，因为它固着于至低之物。接着，在第四阶段是劳苦，在此需要竭力奋斗，以使心灵摆脱那些因有害的甜蜜而纠缠的事物：因此此处饥渴慕义，并且非常需要刚毅；因为以快乐保持的东西，不经痛苦是不会放弃的。然后，在第五阶段，对于那些在劳苦中坚持的人，给予了摆脱它的忠告；因为除非每个人得到更高者的帮助，他绝不可能靠自己从如此巨大的苦难纠缠中解脱出来。但这是一个公正的忠告：希望得到更强者的帮助的人，应在自己更强的地方帮助更弱者：因此\BibleQuote{怜悯人的人是有福的，因为他们要受怜悯\BibleRef{玛 5:7}}。在第六阶段是心地纯洁，能凭良善工作的良心默观那至善，只有纯洁而平静的理智才能辨识至善。最后是第七，即智慧本身——亦即对真理的默观，使整个人安宁，并肖似天主，总结如下：\BibleQuote{缔造和平的人是有福的，因为他们要称为天主的子女\BibleRef{玛 5:9}}。第八，仿佛回到起点，因为它展示并推荐了圆满和成全：因此在第一和第八中，都提到了天国，\BibleQuote{神贫的人是有福的，因为天国是他们的\BibleRef{玛 5:3}}；以及\BibleQuote{为义而受迫害的人是有福的，因为天国是他们的\BibleRef{玛 5:10}}：正如现在所说，\BibleQuote{谁能使我们与基督的爱隔绝？是困苦吗？是窘迫吗？是迫害吗？是饥饿吗？是赤身露体吗？是危险吗？是刀剑吗？\BibleRef{罗 8:35}}因此，带来成全的事物共有七件：因为第八件使那成全的显明并展示出来，仿佛重新开始，其余各件也藉着这些阶段得以成全。\par

% structure: p0017 source=h2 target=subsection reason=relative-heading-rank; base=h2

\subsection{第四章}

11. 因此，依撒意亚所说的圣神七恩，在我看来也与这些阶段和真福句相对应。但顺序有差异：因为那里从较卓越的开始，而这里从较卑微的开始。因为那里以智慧开始，以敬畏天主结束；但\BibleQuote{敬畏上主是智慧的开始\BibleRef{箴 1:7}}。因此，如果我们按逐渐上升的系列来算，那里敬畏天主是第一，孝爱第二，知识第三，刚毅第四，超见第五，明达第六，智慧第七。敬畏天主对应于谦卑的人，即此处所说的\BibleQuote{神贫的人是有福的\BibleRef{玛 5:3}}，即不自高自大、不傲慢的人：宗徒对他们说：\BibleQuote{不可心高气傲，但要恐惧\BibleRef{罗 11:20}}；即不可骄傲。孝爱对应于温良的人：因为虔敬地探求圣经的人，尊敬圣经，不谴责他尚未理解的事，因此不反抗；这就是温良：因此此处说\BibleQuote{温良的人是有福的\BibleRef{玛 5:5}}。知识对应于哀恸的人，他们已在圣经中发现，自己因哪些邪恶而被束缚，而这些邪恶他们曾无知地当作美好有益而贪恋。刚毅对应于饥渴的人：因为他们努力急切地渴望从真正美善的事物中得到喜乐，并急切地寻求将爱从世俗和肉体的事物中转离：因此此处说\BibleQuote{饥渴慕义的人是有福的\BibleRef{玛 5:6}}。超见对应于怜悯的人：因为这是摆脱如此巨大邪恶的唯一补救，即我们宽恕，如同我们愿意自己被宽恕；并尽己所能帮助他人，如同我们在自己无能为力之处渴望得到帮助：因此此处说\BibleQuote{怜悯人的人是有福的\BibleRef{玛 5:7}}。明达对应于心地纯洁的人，眼睛仿佛被洗净，借此可以看见眼睛未曾看见、耳朵未曾听见、人心也未曾想到的事：因此此处说\BibleQuote{心地纯洁的人是有福的\BibleRef{玛 5:8}}。智慧对应于缔造和平的人，在他们身上一切都已井然有序，没有激情反抗理性，而是万物都服从于人的精神，而人自身也服从天主：因此此处说\BibleQuote{缔造和平的人是有福的\BibleRef{玛 5:9}}。\par

12. 此外，那唯一的赏报，即天国，按这些阶段有不同的名称。在第一阶段，正如应当的那样，是天国，即理性灵魂的完美与最高智慧。因此，经文如此说：\BibleQuote{神贫的人是有福的，因为天国是他们的。}仿佛是说：\BibleQuote{敬畏上主是智慧的开端。}温良的人获得产业，犹如孝子寻求父的遗嘱：\BibleQuote{温良的人是有福的，因为他们要承受土地。}哀恸的人得到安慰，如同那些知道自己失去什么、沉沦于何种邪恶的人：\BibleQuote{哀恸的人是有福的，因为他们要受安慰。}饥渴慕义的人得到饱足，犹如那些为救恩劳苦奋战的人获得滋养：\BibleQuote{饥渴慕义的人是有福的，因为他们要得饱饫。}怜悯人的人得到怜悯，如同那些遵循真实而卓越劝告的人，以至于那更强有力的对待他们，正如他们对待那较弱者：\BibleQuote{怜悯人的人是有福的，因为他们要受怜悯。}心地纯洁的人获得看见天主的权能，如同那些带着纯洁之眼辨别永恒事物的人：\BibleQuote{心地纯洁的人是有福的，因为他们要看见天主。}缔造和平的人获得天主的肖像，如同那因更新之人的重生而完全明智、按照天主的肖像被塑造的人：\BibleQuote{缔造和平的人是有福的，因为他们要称为天主的子女。}这些许诺确实可以在今生实现，正如我们相信它们在宗徒们身上实现了一样。因为那在此生之后许诺的、全面改变为天使形态的事，无法用任何言语解释。因此，\BibleQuote{为义而受迫害的人是有福的，因为天国是他们的。}这第八端真福，回到起点，彰显完美的人，或许其意义在旧约中第八天的割礼，以及主在安息日后、那第八天同时也是第一天的复活，以及我们在新人的重生时所庆祝的八天庆节，以及五旬节本身的数字中有所表明。因为七乘以七，我们算得四十九，仿佛加上了第八，使五十得以完成，我们似乎又回到起点；在那一天，圣神被派遣，藉着祂，我们被引入天国，承受产业，得到安慰，得到饱足，获得怜悯，得以洁净，成为缔造和平的人；如此完美，我们为真理和正义的缘故，承受从外而来的所有患难。\par

% structure: p0020 source=h2 target=subsection reason=relative-heading-rank; base=h2

\subsection{第五章}

13. 祂说：\BibleQuote{几时人辱骂你们，迫害你们，捏造一切坏话毁谤你们，为了我，你们是有福的。你们欢喜踊跃罢！因为你们在天上的赏报是丰厚的。}凡以基督徒之名追求此世之乐和暂时之财的人，当思我们的福乐是在内里；就如先知的口论及教会的灵魂所说：\BibleQuote{君王女儿的一切荣耀，是在内里；}因为外面应许的是辱骂、迫害和毁谤；然而，从这些事中，在天上有丰厚的赏报，这在忍受者心中可感受到，他们如今能说：\BibleQuote{我们在磨难中也欢跃，因为我们知道磨难生坚忍，坚忍生老练，老练生望德，望德不叫人蒙羞，因为天主的爱，藉着所赐给我们的圣神，已倾注在我们心中了。}因为忍受这些事并非单凭其本身有益，而是为了基督的名，以平静的心甚至欢欣的心忍受这些事。因为许多异端者以基督徒之名欺骗灵魂，忍受许多此类事；但他们被排除在那赏报之外，因为经文没有仅仅说：\BibleQuote{忍受迫害的人是有福的；}而是加上：\BibleQuote{为义。}凡没有真实信德的地方，就没有义，因为义人因信德而生活。分裂教会者也不可指望那赏报的任何部分；因为同样，凡没有爱德的地方，就不能有义，因为\BibleQuote{爱德不加害于人；}如果他们有爱德，就不会撕碎基督的身体，即教会。\par

14. 但有人可能会问，当祂说\BibleQuote{人若辱骂你们}和\BibleQuote{若有人说一切坏话毁谤你们}有什么区别？因为辱骂正是说坏话毁谤。但有一件事：当辱骂的话当着被辱骂者的面带着轻蔑抛出时，就如有人对我的主说：\BibleQuote{我们说你是撒玛黎雅人，并附有魔鬼，岂不对吗？}；另一件事：当我们的名声在我们不在场时受到损害，就如关于祂的记载：\BibleQuote{有人说：祂是先知；也有人说：不是的，祂是欺骗百姓的。}此外，迫害是施加暴力或以计谋攻击，正如出卖祂的人和钉死祂的人所做的。诚然，关于这一点也并非简单表述为\BibleQuote{若有人说一切坏话毁谤你们}，但加上了\BibleQuote{说谎}这个词，以及\BibleQuote{为了我}的说法；我认为加上这些是为了那些希望在迫害和名誉卑贱中夸耀，并因此说基督属于他们的人，因为有许多坏事是针对他们的；但一方面，所说的话若是关于他们的错误，则是真实的；另一方面，如果有时也抛出一些虚假的指控（这常因人的鲁莽而发生），他们却不是为了基督的缘故而遭受这些。因为凡不是按照真信德和公教纪律被称为基督徒的，就不是基督的门徒。\par

15. 祂说：\BibleQuote{你们欢喜踊跃吧！因为你们在天上的赏报是丰厚的。}我不认为这里所谓的天是指这可见世界的较高部分。因为我们的赏报，应当是稳固而永恒的，不应置于易逝和暂时的事物中。但我认为\BibleQuote{在天上}这句话的意思是在精神的穹苍中，那里存有永恒的公义。与之相比，邪恶的灵魂被称为地，当它犯罪时，就对它说：\BibleQuote{你是土，还要归于土。}关于这个天，宗徒说：\BibleQuote{我们的家乡是在天上。}因此，那些在精神之善中喜乐的人，现在就已经意识到那赏报；但那时，当这必死的也穿上了不朽，它将在每一部分得以完善。祂说：\BibleQuote{因为人们也这样迫害了你们以前的先知。}在此，祂以一般的含义使用了\BibleQuote{迫害}一词，同样适用于辱骂的言语和破坏名誉的行径；并且以榜样很好地鼓励了他们，因为说真话的人通常遭受迫害：然而，古代的先知们并没有因此，因惧怕迫害而放弃宣讲真理。\par

% structure: p0024 source=h2 target=subsection reason=relative-heading-rank; base=h2

\subsection{第六章}

16. 因此，下面这句话是完全正确的：\BibleQuote{你们是地上的盐。}这表明那些或因热切追求丰厚的地上祝福，或因惧怕匮乏而失去那既不能由人给与也不能被夺走的永恒事物的人，将被视为无味。\BibleQuote{盐若失了味，可用什么使它再咸呢？}即如果你们——藉着你们，各族在某种程度上得以保存[免于败坏]——因惧怕暂时的迫害而失去天国，那么从哪里能找到人来除去你们中的错误呢？因为天主拣选了你们，为的是藉着你们除去其他人的错误。因此，失了味的盐\BibleQuote{毫无用处，只好被丢在外面，任人践踏。}所以，不是那受迫害的人，而是那因恐惧迫害而变得无味的人，被人践踏。因为只有最底层的人才会被践踏；但他并非最底层的人，无论他在世上身体遭受多少痛苦，他的心却固守在天上。\par

17. \BibleQuote{你们是世上的光}。正如祂前面所说的\BibleQuote{地上的盐}，如今祂说\BibleQuote{世上的光}。因为在前者中，\Emph{地}不是指我们用肉身之脚所踏的地，而是指居住在地上的人，甚至是指罪人——主派遣宗徒的盐就是为了保存他们并消除他们的败坏。而在这里，\Emph{世界}必须理解不是指天和地，而是指在世上或爱世界的人——宗徒们被派遣正是为了光照他们。\BibleQuote{建在山上的城是不能隐藏的}，\Emph{即}［一座城］建立在伟大而显著的正义之上，这也是我们的主讲道的那座山本身的意义。\BibleQuote{人点灯，不放在斗底下}。我们当如何理解？即\BibleQuote{放在斗底下}这一说法仅表示灯的隐藏，好像祂在说：没有人点灯而将它隐藏起来？还是\OriginalTerm{斗}{bushel measure}也有含义，以致将灯放在斗底下就是将身体的舒适置于真理的宣讲之上，因而因害怕在肉身和暂时的事物中遭受任何烦恼而不宣讲真理？\Quote{斗}的称呼很恰当，无论是因为量器的报应——每个人都按自己身体所行的接受——宗徒说：\BibleQuote{好使每一个人都得着他在身体上所行的}；另一处似乎论及身体的这一\Quote{斗}时说：\BibleQuote{你们用什么量器量给人，也要用什么量器量给你们}——还是因为暂时的好事，在身体中实现，在一定数目的日子中开始并结束，这或许就是\BibleQuote{斗}的含义；而永恒和属灵的事不受此限制，\BibleQuote{因为天主赐圣神没有限量}。因此，凡用暂时的舒适遮蔽和掩盖美好教义之光的人，就是把他的灯放在斗底下。\BibleQuote{而是放在灯台上}。现在，人若将身体服从于天主的服务，就把它放在灯台上，这样真理的宣讲是更高的，身体的服侍是更低的；然而，藉着甚至身体的服侍——即藉着声音、舌头和善行中身体的其他动作——教义更加明亮地照耀出来。因此，宗徒把灯放在灯台上，他说：\BibleQuote{所以我斗拳，不像打空气的；我克制我的身体，使它服从，免得我给别人传了福音，自己反被弃绝}。然而当祂说\BibleQuote{照亮屋中所有的人}时，我认为\Emph{屋}是指人的居所，即世界本身，因为祂之前说过\BibleQuote{你们是世上的光}；或者若有人愿意将\Emph{屋}理解为教会，这也并非不妥。\par

% structure: p0027 source=h2 target=subsection reason=relative-heading-rank; base=h2

\subsection{第七章}

18. \BibleQuote{你们的光，}祂说，\BibleQuote{当这样在人前照耀，使他们看见你们的善行，并光荣你们在天上的圣父。}如果祂仅仅说：\BibleQuote{你们的光当这样在人前照耀，使他们看见你们的善行，}祂似乎将目标定在人的赞美上——这是假善人和那些追求荣誉、贪图最虚幻荣耀的人所寻求的。针对这类人，经上说：\BibleQuote{如果我还讨人的欢喜，我就不是基督的仆人了}；先知说：\BibleQuote{讨人喜欢的人蒙受羞辱，因为天主轻视了他们}；又说：\BibleQuote{天主打碎了讨人喜欢的人的骨头}；宗徒又说：\BibleQuote{不要贪图虚荣}；另一次说：\BibleQuote{各人该省察自己的行为，这样，只在自身有可夸耀的，而不在别人身上}。因此，我们的主不仅说了\BibleQuote{使他们看见你们的善行}，还加上\BibleQuote{并光荣你们在天上的圣父}：这样，一个人藉善行讨人喜欢，并不以此为目的——讨人喜欢；而应将此归于对天主的赞美，并为此讨人喜欢，好使天主在他内受到光荣。因为这对那些赞美的人有益：他们应尊崇的，不是人，而是天主；我们的主在抬瘫痪者的例子中表明，当瘫子被治愈时，群众惊奇祂的能力，如福音所载：\BibleQuote{害怕起来，光荣了那赐给人们这样权能的天主}。祂的效仿者宗徒保禄说：\BibleQuote{他们只是听说：那从前迫害我们的，如今传布他曾经要使消灭的信仰；他们就为我光荣了天主。}\par

19. 因此，在祂劝勉听众预备为真理和正义忍受一切，并不要隐藏他们将要领受的美善，而应以如此的慈爱学习以便教导他人，在善行中不追求自己的赞美，只求天主的荣耀之后，祂开始指示和教导他们应当教导什么；好像他们问祂说：看，我们愿意为祢的名忍受一切，也不隐藏祢的教义；但祢禁止我们隐藏、祢命令我们忍受的一切，究竟是什么？祢是否要提及与法律和先知所记相反的事？\BibleQuote{不，}祂说，\BibleQuote{你们不要以为我来是废除法律或先知；我来不是为废除，而是为成全。}\par

% structure: p0030 source=h2 target=subsection reason=relative-heading-rank; base=h2

\subsection{第八章}

20. 这句话的含义是双重的。我们必须从两个方面来处理它。因为那说\BibleQuote{我来不是为废除法律，而是为成全}（\BibleRef{玛 5:17}）的那位，要么是指添加缺少的部分，要么是指实行其中的内容。那么，让我们先考虑我首先提出的那种含义：因为添加缺少的部分的人，当然不是废除他所找到的，而是通过使之完善来巩固它；因此，他接着陈述说：\BibleQuote{我实在告诉你们：即使天地过去了，一撇或一画也决不会从法律中过去，直到一切全都实现。}（\BibleRef{玛 5:18}）因为，如果连那些为完成而添加的东西都已实现，那么那些作为开端而预先送出的东西就更得以实现了。然后，至于他所说的\BibleQuote{一撇或一画也决不会从法律中过去}（\BibleRef{玛 5:18}），除了理解为一种强烈的完美表述之外，别无他意，因为这是通过单个字母指出的，而在这些字母中，\Quote{\OriginalTerm{一撇}{\GreekText{ἰῶτα}}}比其他字母更小，因为它由一笔写成；而\Quote{\OriginalTerm{一画}{\GreekText{κεραία}}}甚至只是该字母顶端的一种微小部分。通过这些话语，他表明法律中所有最微小的细节都必须付诸实行。之后他补充说：\BibleQuote{所以，谁若废除这些诫命中最小的一条，并这样教训人，他在天国里将称为最小的。}（\BibleRef{玛 5:19}）因此，最小的诫命就是指\Quote{一撇}和\Quote{一画}。所以，\BibleQuote{谁若废除并这样教训人}——即按照他所废除的，而不是按照他所发现和阅读的——\BibleQuote{他在天国里将称为最小的}（\BibleRef{玛 5:19}）；因此，也许他根本就不在天国里，因为那里只有伟大的才能进入。\BibleQuote{但谁若实行并这样教训人}——即他不废除，并按照他不废除的这样教训人——\BibleQuote{他在天国里将称为大的。}（\BibleRef{玛 5:19}）至于那在天国里称为大的人，他自然也在天国里，因为大的得以进入；因为随后的话正指此而言。\par

% structure: p0032 source=h2 target=subsection reason=relative-heading-rank; base=h2

\subsection{第九章}

21. \BibleQuote{我告诉你们：除非你们的义德超过经师和法利塞人的义德，你们决不能进入天国。}（\BibleRef{玛 5:20}）即，除非你们不仅履行那些开始成人的法律中最小的诫命，而且也履行由我添加的那些诫命——我来不是为废除法律，而是为成全——你们就不能进入天国。但你会对我说：如果以上他谈论那些最小的诫命时，他说谁废除其中一条并按照他的违抗教训人，他在天国里称为最小的；但谁实行它们并这样教训人，则称为大的，因而已经是在天国里，因为他是大的；那么，为什么还需要给法律的最小诫命添加东西呢？如果因为谁实行它们并这样教训人就是大的，他已经能进入天国了吗？因此，那句话应这样理解：\BibleQuote{但谁若实行并这样教训人，他在天国里将称为大的}（\BibleRef{玛 5:19}）——即不是按照那些最小的诫命，而是按照我即将提到的那些。那么这些是什么呢？他说：\BibleQuote{使你们的义德超过经师和法利塞人的义德}（\BibleRef{玛 5:20}）；因为除非你们的义德超过他们的，你们决不能进入天国。所以，谁废除那些最小的诫命并这样教训人，将称为最小的；但谁实行那些最小的诫命并这样教训人，不一定就被视为大的和适合天国的；但他还不至于像废除它们的人那么小。然而，为了使他成为大的并适合那个国度，他应当照基督现在所教导的去实行和教训，即为了使他的义德超过经师和法利塞人的义德。法利塞人的义德是：不可杀人；那些注定进入天主之国的人的义德是：不可无端发怒。因此，最小的诫命就是不可杀人；谁废除这条，将在天国里称为最小的；但谁履行这条不可杀人的诫命，未必就自然而然地成为大的和适合天国，不过他却登上了一级台阶。然而，如果他无端不发怒，他将得以成全；如果他做到这点，他将离谋杀更远。因此，那教导我们不应发怒的人，并不废除不可杀人的法律，反而成全了它；这样我们既在外表上不杀人，也在心里不发怒，从而保持我们的无辜。\par

22. 因此，他说：\BibleQuote{你们曾听见对古人说：不可杀人；凡杀人的，应受审判。我却对你们说：凡向自己的弟兄发怒的，应受审判；谁若向自己的弟兄说\InnerQuote{\OriginalTerm{拉卡}{\GreekText{ῥακά}}}，应受公议会制裁；谁若说\InnerQuote{\OriginalTerm{愚妄的人}{\GreekText{μωρέ}}}，应受火狱的罚。}（\BibleRef{玛 5:21-22}）那么，应受审判、应受公议会制裁和应受火狱的罚之间有什么区别？因为最后一项听起来最严厉，并提醒我们某些阶段从较轻到较重地过渡，直到达到火狱。因此，如果应受审判比应受公议会制裁更轻，应受公议会制裁也比应受火狱的罚更轻，那么我们就必须理解，无端向弟兄发怒比说\Quote{拉卡}更轻；同样，说\Quote{拉卡}比说\Quote{愚妄的人}更轻。因为如果罪过不是按等级提到，危险也不会有等级。\par

23. 但这里有一个含混的词出现了，因为\Quote{Raca}既不是拉丁文也不是希腊文。然而，其他词在我们的语言中通用。有些人想从希腊文推导这个词的解释，认为一个衣衫褴褛的人被称为\Quote{Raca}，因为一块破布在希腊文中叫\OriginalTerm{破布}{\GreekText{ῥάκος}}；但当人们问他们一个衣衫褴褛的人在希腊文中叫什么时，他们并不回答\Quote{Raca}；而且，拉丁文译者在原本应该用\Quote{衣衫褴褛}的地方用了\Quote{Raca}，而没有使用一个既不存在于拉丁语、在希腊语中也罕见的词。因此，我更倾向于我从一位询问过的希伯来人那里听到的看法；他说这个词不代表任何东西，而只是表达愤怒情绪。语法学家将那些表达激动情绪的情感称为\OriginalTerm{感叹词}{interjections}；如同悲伤的人说\Quote{唉}，或愤怒的人说\Quote{哈}。这些词在所有语言中都是专有名词，不易翻译成另一种语言；这个原因确实迫使希腊文和拉丁文译者都保留了这个词本身，因为他们找不到翻译它的方法。\par

因此，所提及的罪有一个层级，首先是发怒，并将这种情绪作为意念藏在心中；但如果这种情绪发出一声愤怒的表达，没有任何特定含义，只是通过爆发的行为显示出那被愤怒所攻击者的情绪状态，这当然比将升起的愤怒在沉默中抑制更严重；但如果听到的不仅是一声愤怒的表达，而且还有一个词，使用者通过这个词指向并表示对受指责者的一种明确谴责，谁会怀疑这比仅仅发出一声愤怒的感叹更严重呢？因此，第一等有一样东西，即仅仅是愤怒；第二等有两样东西，即愤怒和表达愤怒的词；第三等有三样东西，即愤怒、表达愤怒的词，以及这个词中吐露的明确谴责。现在也看审判、公议会、火狱这三个责任的等级。因为在审判中，仍然有机会辩护；然而在公议会中，虽然通常也有审判，但由于这种区分本身迫使我们在这一点上承认存在某种差异，判决的宣示似乎属于公议会，因为现在不是考虑被告本人是否应被定罪的问题，而是审判官们彼此商议应定他什么刑罚，他已显然要被定罪；但火狱既不似审判那样对待定罪作为悬疑之事，也不似公议会那样对待被定罪者的刑罚；因为在火狱中，被定罪者的定罪和刑罚都是确定的。这样，在罪和罪的惩罚中看到一定的等级；但谁能说出它们在灵魂的惩罚中如何无形地显示出来呢？因此，我们要学习法利塞人的义和那更大的引人进入天国的义之间的差别有多大：因为杀人比用言语侮辱更严重，在前一情形，杀人使人受审判，而在后一情形，愤怒使人受审判，愤怒是这三种罪中最轻的；因为在前一情形，他们讨论的是人间杀人问题，而在后一情形，一切都由天主审判处置，被定罪者的结局是火狱。但若有人说在更大的义之下，杀人受更严厉的惩罚，因为侮辱已受火狱之罚，这就迫使我们理解火狱也有差别。\par

确实，在我们面前的这三句话中，我们必须注意有些词是隐含的。因为第一句话包含了所有必需的词。祂说：\Quote{凡无缘无故向弟兄发怒的，应受审判。}但在第二句，祂说：\Quote{凡向弟兄说Raca的}，其中隐含了\Quote{无缘无故}，因此接着是\Quote{应受公议会的裁判。}在第三句，祂说：\Quote{但凡说\InnerQuote{愚人}的}，其中隐含了两件事，即\Quote{向弟兄}和\Quote{无缘无故}。这样，我们为宗徒辩护，他称迦拉达人为愚人，也称呼他们为弟兄；因为他不是无缘无故这样做的。这里理解\Quote{弟兄}一词是因为后面谈到了仇敌的情况，以及在更大的义之下应如何对待他。\par

% structure: p0038 source=h2 target=subsection reason=relative-heading-rank; base=h2

\subsection{第十章}

26. 接下来是：\BibleQuote{所以，你若将你的礼物带到祭台前，在那里想起你的弟兄有什么怨你的事，就将你的礼物留在祭台前，先去与你的弟兄和好，然后再来献上你的礼物。}由此显然，上面所说的关于弟兄的话，是同下面这句话相连的，因为接下来的句子以这样的连词连接，证实了前面的话；因为祂不是说：\BibleQuote{但是，你若将你的礼物带到祭台前}，而是说：\BibleQuote{所以，你若将你的礼物带到祭台前。}因为如果向弟兄无缘无故发怒，或说\BibleQuote{辣加}，或说\BibleQuote{愚人}都是不合法的，那么，在心里保留任何东西以致愤恨变成仇恨就更不合法了。这也属于另一段经文所说的：\BibleQuote{不可让太阳在你们的愤怒上落下。}因此，我们被命令，当我们要将礼物带到祭台前时，如果我们想起弟兄有什么怨我们的事，就要将礼物留在祭台前，先去与弟兄和好，然后再来献上礼物。但如果按字面理解，也许有人会认为，如果弟兄在场，就应该这样做；因为命令你留下礼物在祭台前，所以不能拖延太久。因此，如果你想起某个不在场、甚至可能远在海外的人，那么认为要留下礼物在祭台前，直到你跨过陆地和海洋之后才能献给天主，这是荒谬的。因此，我们不得不求助于完全内在和属灵的解释，以便所说的话毫无荒谬之处。\par

27. 因此，我们可以从属灵的角度解释祭台，即天主内在圣殿中信德本身，其标记是可见的祭台。因为我们献给天主的任何奉献，无论是预言、教导、祈祷、圣咏、赞美诗，或其他任何类似的属灵礼物出现在心中，除非有真诚的信德支撑，并如同稳固地、不动地放在其上，使我们所言说的保持完整无损，否则都不能蒙天主悦纳。因为许多异端者没有祭台，即真信德，却将亵渎当作赞美说出；他们被世俗的见解压制，因而如同将他们的奉献扔在地上。但献祭者也应当有纯洁的意向。因此，当我们心中要献上任何这样的奉献时（即在天主的内在圣殿中，因为经上说：\BibleQuote{天主的圣殿是圣的，这圣殿就是你们。}以及\BibleQuote{使基督因信德住在你们心里，住在你们的内人里。}），如果我们心中想起弟兄有什么怨我们的事，即我们若在什么事上伤害了他（因为那时他有什么怨我们的事；相反，他若伤害了我们，我们就有怨他的事；在这种情况下，无需进行和好：因为你不会向伤害你的人请求宽恕，而只是宽恕他，如同你希望主宽恕你得罪祂的事一样），因此，当我们心中想起我们或许在什么事上伤害了弟兄时，就应去和好；但这不是用身体的脚，而是用心中的情感，你要以谦卑的态度俯伏在弟兄面前，以亲切的思念投奔他，在天主面前——你将要献上奉献的那一位——如此，即使他就在面前，你也能以不矫饰的心软化他，通过请求宽恕使他恢复善意，如果你首先在天主面前这样做，不是以身体的缓慢行动，而是以爱的急速冲动走向他；然后前来，即重新将你的心思转向你开始要做的事，你就会献上你的礼物。\par

28. 但谁能做到既不无缘无故向弟兄发怒，也不无缘无故说\BibleQuote{辣加}，也无缘无故称他为愚人（这些事都是最骄傲地犯下的），或者，如果偶然陷入其中任何一种，能以恳求的心请求宽恕（这是唯一的补救）呢？除了那不被虚荣夸耀的精神所膨胀的人，还有谁呢？因此，\BibleQuote{神贫的人是有福的，因为天国是他们的。}现在让我们看看下文。\par

% structure: p0042 source=h2 target=subsection reason=relative-heading-rank; base=h2

\subsection{第十一章}

29. 祂说：\BibleQuote{你要赶快与你的对头和解，当你还在路上同他一起时；免得对头把你交给审判官，审判官把你交给差役，你就被投入监里。我实在告诉你，除非你还清最后一文小钱，你决不能从那里出来。}我明白审判官是谁：\BibleQuote{因为圣父不审判任何人，但把一切审判都交给了圣子。}我明白差役是谁：经上说：\BibleQuote{天使前来伺候祂。}我们相信祂将率领祂的天使审判活人死人。我明白监牢的意思：显然是指惩罚的黑暗，祂在另一处称之为\Quote{外界的黑暗}；我相信，这是因为神圣奖赏的喜乐是内心本身的某种内在喜乐，或者甚至更隐秘的，正如对那配受的仆人所说的：\BibleQuote{进入你主人的喜乐吧！}同样，在这个共和政体下，一个人被投入监牢，是从议事厅或审判官的宫殿被赶出去的。\par

30. 但如今，关于偿还最后一文钱，这可以无荒谬地理解为：要么是指没有留下任何未受惩罚之事；正如我们日常说话中也说\Quote{到渣滓为止}，当我们想表达某物被倒空以致毫无剩余时；要么，藉着\Quote{最后一文钱}可能指属地之罪。因为，作为此世界各组成部分的第四部分，且实际上作为最后一部分，地是被发现的；因此，你以天为始，以空气为第二，水为第三，地为第四。因此，说\Quote{直到你还清了最后的第四部分}，似乎可恰当地理解为\Quote{直到你赔补了你的属地之罪}：因为罪人也曾听闻：\BibleQuote{你是土，还要归于土。}那么，关于\Quote{直到你还清了}这一表达，我怀疑它是否指称为永罚的那种惩罚。因为，当不再有悔改和度更正确生活的机会时，那债如何偿还呢？或许，此处\Quote{直到你还清了}的意思，与那段经文中所说的相同：\BibleQuote{坐在我右边，直到我使你的仇敌成为你的脚凳}；因为即使仇敌已被置于祂脚下，祂也不会停止坐在右边；或者宗徒的声明：\BibleQuote{因为祂必须为王，直到把一切仇敌放在祂脚下}；因为即使它们已被放在祂脚下，祂也不会停止为王。因此，正如那里对祂的理解——关于祂说\BibleQuote{祂必须为王，直到把祂的仇敌放在祂脚下}——即祂将永远为王，因为它们将永远在祂脚下：同样，这里也可以理解为关于他的——关于他说\BibleQuote{你决不可能从那里出来，直到你还清了最后一文钱}——即他将永远不会出来；因为他永远在偿还最后一文钱，只要他正在承受他属地之罪的永罚。我这样说，并非要阻止对罪的惩罚进行更仔细的讨论，即圣经中如何称之为永；尽管无论如何，它更应被避免而非被认识。\par

31. 但让我们现在看看，这对头本人是谁，我们被命令要与他快速达成和解，当我们与他同行在路上时。因为他要么是魔鬼，要么是人，要么是肉体，要么是天主，要么是祂的诫命。但我不明白，我们如何被命令要与魔鬼处于善意之中，即与他同心或同意。因为有些人将此处出现的希腊文词译为\Quote{同心}，另一些人译为\Quote{同意}：但我们既未被命令向魔鬼表示善意（因为哪里有善意，哪里就有友谊；没有人会说我们要与魔鬼交朋友）；与他达成和解也无益，因为我们曾一次而永远地弃绝了他，向他宣战，并将在战胜他后获得冠冕；我们现在也不应向他屈服，因为如果我们从未向他屈服，我们绝不会陷入如此悲惨的境地。再者，关于对头是一个人的观点，尽管我们被命令，只要力所能及，要与众人和平相处，在那里当然可以理解为善意、和谐与同意；然而我不明白，我如何能接受这种观点，即我们被一个人交给审判官，在此情况下我理解基督是审判官，宗徒说我们\BibleQuote{都必须在祂的审判台前出现}：那么，他怎能将我交给审判官，而他自己也将与我一同出现在审判官面前呢？或者，如果有人因伤害了他人而被交给审判官，尽管受伤害的一方并未交出他，更合适的观点是，犯罪者因违反了伤害人时与之相悖的法律而被交给审判官。这还有一个原因：如果有人因杀人而伤害了他人，那么现在就没有时间与他达成和解了；因为他不再与他同行在路上，即在此生中：然而，补救并不会因此被排除，如果一个人悔改，并以痛悔的心作为祭献，投奔于那赦免归向祂者之罪、并\BibleQuote{为一个悔改者比九十九个义人更欢喜}的那位的仁慈。但更少的是，我明白如何被命令要向肉体怀有善意、与之达成一致或向它屈服。因为，反而是罪人更爱自己的肉体，与之达成一致并屈服于它；但那迫使肉体服从的人，并不是向它屈服的人，而是强迫它向他们屈服。\par

32. 因此，或许我们奉命向天主让步，善待祂，以便我们能和祂和好，我们因犯罪而远离了祂，以致祂可被称为我们的对头。因为祂理当被称为祂所抵抗者的对头，因为\BibleQuote{天主拒绝骄傲人，却赐恩宠于谦卑人。}并且\BibleQuote{骄傲是万恶的起源，但人的骄傲的开始是背离天主；}宗徒也说：\BibleQuote{因为假如我们还在为仇的时候，因祂圣子的死而和天主和好了，那么，在和好之后，我们更要因祂的生命而得救了。}由此可知，没有一个本性因是恶劣的而成为天主的敌人，因为正是那些曾为仇敌的人正在和好。因此，任何人在此路上，即在此生中，若未因祂圣子的死而与天主和好，将被祂交给审判者，因为\BibleQuote{父不审判任何人，但把审判的全权交给了子。}而此后所描述的其余事情便依次发生，这些我们已经讨论过了。关于这种解释，只有一个难题，即：如果我们在此段落中将祂本人理解为恶人的对头，而我们要奉命与祂尽快和好，那么如何能正确地说我们与天主同在途中呢？除非，或许因为祂无处不在，我们在此途中也必然与祂同在。因为经上说：\BibleQuote{我若上升于高天，你已在那里；我若下降于阴府，你也在那里。我若飞旋在晨曦的翅膀，定居在大海的尽头，你的手仍在那里引导我，你的右手仍在那里扶持我。}或者，如果不接受恶人被称为与天主同在的观点（尽管没有一处不是天主临在——正如我们不说瞎子和光在一起，尽管光环绕他们的眼睛），那么还有一个办法：我们应将这里的对头理解为天主的诫命。因为有什么比天主的诫命，即祂的法律和神圣圣经，更与那些想要犯罪的人作对呢？这圣经是为今生赐给我们的，使它与我们一同行在途中，我们不可违背它，免得它将我们交给审判者，而应当赶快顺从它。因为没有人知道自己何时会离开此生。那么，除了虔诚地诵读或聆听，并因它拥有最高权威而服从它的人以外，谁还会顺服圣经呢？这样他对自己所理解的，并不因为觉得它反对自己的罪恶而憎恨它，反而喜爱受它的责备，并因自己的病患不被姑息直到痊愈而欢喜；这样，纵使遇到他认为晦涩或荒谬之处，他也不会因此挑起好争辩的矛盾，而是祈求得以明白，但心中牢记要对如此伟大的权威表现出善意和敬畏。但谁这样做呢？除非那人不以严厉威胁，而是以虔诚的温良前来，为要开启并确定他父亲遗嘱的内容。因此，\BibleQuote{温良的人是有福的，因为他们要承受土地。}让我们看下面的话。\par

% structure: p0047 source=h2 target=subsection reason=relative-heading-rank; base=h2

\subsection{第十二章}

33. \BibleQuote{你们曾听见古人的话说：不可奸淫。我却对你们说：凡注视妇女，有意贪恋她的，他早已在心里奸淫了她。}因此，较小的义是不因肉体的结合而行奸淫；但天国的更大的义是不在心中行奸淫。凡不在心中行奸淫的人，更容易防备在实际行动中行奸淫。因此，颁布后来诫命的那位证实了先前的诫命；因为祂来不是要废除法律，而是要成全。很值得思考的是，祂并没有说\Quote{凡贪恋妇女的}，而是说\Quote{凡注视妇女，有意贪恋她的}，即以这个目的和意向转向她，为要贪恋她；这实际上不仅仅是肉欲的触动，而是完全同意贪欲；以至于禁止的欲望不被抑制，而是机会若允许，就得到满足。\par

34. 因为有三件事构成完整的罪：暗示、取乐和同意。暗示通过记忆或身体感官发生，当我们看见、听见、嗅到、尝到或触摸到某物时。如果从中取乐，这种快乐如果是非法的，就必须被抑制。正如我们禁食时，看见食物，味觉被激起，这不可能没有快乐；但我们不同意这种喜好，我们以理性的权利压制它。但如果同意发生，罪就完成，在天主心中是已知的，尽管可能没有通过行为被人知晓。因此，这些步骤是：暗示由蛇发出，即身体的短暂而快速、即逝的运动；因为如果在灵魂中也有任何这样的图像移动，它们是从外部从身体得来的；如果除了那五种感官之外还有任何隐藏的身体感觉触及灵魂，那也是短暂即逝的；因此它越隐秘地潜入以影响思维过程，就越适宜地被比作蛇。因此，这三个阶段，正如我开始说的，类似于\WorkTitle{创世纪}中描述的那件事，暗示和一定程度的劝诱由蛇发出；而取乐在于肉欲，如同在\Person{厄娃}中；同意在于理性，如同在男人中；这些事情完成后，人就被逐出乐园，即从最幸福的光明正义中进入死亡——这是最公义的。因为发出劝诱的并不强迫。所有本性按其等级在其秩序中是美丽的；但我们绝不能从较高的（理性心灵被置于其中的较高位置）下降到较低的。也没有人被强迫这样做；因此，如果他做了，他就因天主的公义法律而受惩罚，因为他不是不情愿地犯此罪。但是，在习惯形成之前，要么没有快乐，要么快乐如此轻微以致几乎不存在；屈服于它是一个大罪，因为这种快乐是非法的。现在，当一个人屈服时，他就在心中犯了罪。然而，如果他还付诸行动，欲望似乎得到满足并熄灭；但后来，当暗示重复时，更大的快乐被点燃，然而这仍然远小于通过持续实践转化为习惯的快乐。因为克服后者非常困难；然而即使习惯本身，如果一个人不证明自己不忠，不因害怕基督的战斗而退缩，那么在他的（即基督的）领导和帮助下，他将战胜它；这样，按照原始的和平与秩序，男人服从基督，女人服从男人。\par

35. 因此，正如我们通过三个步骤达到罪——暗示、取乐、同意——同样，罪本身有三种变化——在心中、在行为中、在习惯中——如同三种死亡：一种，如同在屋内，即当我们心中同意肉欲时；第二种，如同已带到门外，即当同意走向行动时；第三种，当心灵被坏习惯的力量压制住，如同被土堆覆盖，现在好像在坟墓中腐烂时。凡阅读\WorkTitle{福音}者，都看到我们的\Person{主}使这三种死人复活。或许他也思考在复活者的同一句话中可能发现什么区别，因为祂一次说：\BibleQuote{闺女，起来！}另一次说：\BibleQuote{青年人，我对你说，起来！}另一次，祂在心神中叹息，流泪，又叹息，然后\BibleQuote{大声喊说：\Person{拉匝禄}，出来！}\par

36. 因此，在本段提到的奸淫类别下，我们必须理解所有肉体和感官的肉欲。因为圣经如此频繁地称偶像崇拜为淫乱，而\Person{保禄}宗徒称贪婪为偶像崇拜，谁还怀疑任何邪恶的肉欲都可以恰当地称为淫乱呢？因为灵魂忽视了统治它的更高法律，并为了较低本性的卑鄙快乐而自贬（可以说是作为报酬），从而被败坏？因此，凡在自己因犯罪习惯而感觉肉体快乐反抗正确倾向、其未受制服的暴力将其拖入捆绑的人，应当尽可能回想他因犯罪而失去了何种和平，并呼喊说：\BibleQuote{我这个人真不幸呀！谁能救我脱离这该死的肉身呢？感谢天主，藉着我们的主耶稣基督。}因为这样，当他喊出自己不幸时，在哀叹中他恳求安慰者的帮助。而当他认识到自己的不幸时，这绝非小的接近幸福；因此\BibleQuote{哀恸的人是有福的，因为他们要受安慰。}\par

% structure: p0052 source=h2 target=subsection reason=relative-heading-rank; base=h2

\subsection{第十三章}

37. 接着，祂继续说道：\BibleQuote{若你的右眼使你跌倒，就把它剜出来，丢掉；因为对你来说，丧失一个肢体，而全身不投入地狱，是有益的。} 这里当然需要极大的勇气，以便割掉自己的肢体。因为无论这\Quote{眼睛}指的是什么，无疑都是被热爱的。因为那些想强烈表达情感的人常常这样说：我爱他如自己的眼睛，甚至胜过自己的眼睛。然后，当加上\Quote{右}字时，也许是为了增强爱的强度。因为虽然我们这些肉体的眼睛在看东西时方向相同，且若两者都转，能力相等，但人们更害怕失去右眼。所以在这种情况下，意思是：无论你如此珍爱、视为右眼的东西，如果它使你跌倒，\Emph{即}它阻碍你通往真福的道路，就剜出来丢掉。因为对你来说，丧失一个你所深爱、如肢体般依附的东西，比整个身体被投入地狱更有益。\par

38. 但由于祂接着对右手作了类似的声明：\BibleQuote{若你的右手使你跌倒，就把它砍下来，丢掉；因为对你来说，丧失一个肢体，而全身不投入地狱，是有益的。} 这迫使我们去更仔细地探究祂所说的眼睛是什么。就这一探究而言，我觉得最恰当的解释莫过于一位深爱的朋友：因为无疑，我们可以正当地称其为肢体，即我们热爱的对象；而且这位朋友是一位顾问，因为它就像一只眼睛，指示道路；并且是在神圣事物中，因为它是右眼：如此左眼确实是一位亲爱的顾问，但关乎世俗的、身体必需的事物；关于这些作为绊脚石的事物，提及它是多余的，因为连右眼也不能饶恕。现在，一位在神圣事物中的顾问，若他试图在宗教和教义的外表下引导人陷入任何危险的异端，就成了绊脚石。因此，也让右手被理解为在神圣工程中亲爱的帮助者和助手：因为正如默观被恰当地理解为位于眼睛中，行动则位于右手中；这样左手可以被理解为关乎今生和身体所需的工程。\par

% structure: p0055 source=h2 target=subsection reason=relative-heading-rank; base=h2

\subsection{第十四章}

39. \BibleQuote{又有话说：凡休妻的，应当给她休书。} 这是法利塞人较小的义，主所说的并不反对这一点：\BibleQuote{但我告诉你们：凡休妻的，除非为了淫乱的缘故，就是使她犯奸淫；人若娶被休的妇人，也是犯奸淫。} 因为那位命令要给休书的主，并没有命令休妻；而是说\Quote{凡休妻的}，就得\Quote{给她休书}，以便休书的念头能缓和休妻者鲁莽的怒气。因此，那位设法在休妻中插入延迟的主，已尽其所能向硬心之人表明祂不愿分离。因此，主自己在另一处，当有人就此问题询问祂时，回答说：\BibleQuote{梅瑟因你们心硬，才容许这样做。} 因为无论一个想休妻的人心有多硬，当他想到一旦给了休书，她就能无风险地再嫁，他就容易平息怒气。因此，主为巩固妻子不应轻易被休的原则，只以淫乱为例外；但命令所有其他的烦恼，如果出现的话，为了婚姻的忠贞和贞洁，要坚忍地承受；祂也称那娶被休之妇的人是奸夫。宗徒保禄显示了这一状态的限制，他说只要丈夫活着，就要遵守；但丈夫死后，他允许再婚。因为他自己也持守这个规则，并且在其中引用的不是自己的意见，如同他某些劝诫中的情况，而是主的命令，他说：\BibleQuote{至于那已经嫁娶的，我吩咐他们，其实不是我吩咐，乃是主吩咐说：妻子不可离开丈夫；若是离开了，不可再嫁，或是仍同丈夫和好。丈夫也不可离弃妻子。} 我相信，根据类似的规则，如果丈夫休妻，他应保持独身，或与妻子和好。因为可能他因淫乱的缘故休妻，这是主所允许的例外。但现在，既然妻子在离开丈夫、丈夫还活着时不能结婚，丈夫在妻子活着时也不能另娶，那么任何方面的非法淫乱更是不可行的。那些婚姻，无论是生育子女后，还是因轻视地上的后裔，双方能共同同意彼此节制，更应被视为有福的：因为这样做既不违反主禁止休妻的命令（因为他并没有休妻，而是与她过着非肉身而是属灵的生活），也遵守了宗徒所表达的原则：\BibleQuote{从此以后，那有妻子的，要像没有妻子。}\par

% structure: p0057 source=h2 target=subsection reason=relative-heading-rank; base=h2

\subsection{第十五章}

40. 但主在另一处所作的、常使小者（那些虽热切渴望现今按基督的诫命生活的人）困惑的声明，更是如此：\BibleQuote{如果谁到我这里来，而不恨自己的父亲、母亲、妻子、儿女、兄弟、姊妹，甚至自己的性命，他不能作我的门徒。} 因为在不太聪明的人看来，这可能是一个矛盾：此处主禁止休妻，除非为了奸淫的缘故；但在别处祂却说，凡不恨自己妻子的，不能作祂的门徒。但若祂是就性交而言，就不会将父亲、母亲、兄弟放在同一类别。然而\BibleQuote{天国是以猛力夺取的，以猛力夺取的人就攫取了它}是何等真实！因为一个人要爱仇敌，却恨自己的父亲、母亲、妻子、儿女、兄弟，需要何等的猛力！因为那召叫我们进入天国的，同时命令了这两件事。在祂的引导下，容易显示这些事如何并不互相矛盾；但一旦理解了，实行起来却困难，尽管在祂亲自助佑时，这也非常容易。因为在那永恒的国度里——祂恩准将祂的门徒（祂也称他们为兄弟）召叫进去——没有这类时间性的关系。因为\BibleQuote{没有犹太人，也没有希腊人，没有奴隶，也没有自由人，没有男人，也没有女人；} \BibleQuote{惟有基督是一切，并在一切之内。} 主亲自说：\BibleQuote{因为在复活时，人不娶也不嫁，而是像天上的天使一样。} 因此，任何人若愿意在此刻追求那国度的生命，就必须恨的不是人本身，而是那些维系我们生命的暂时关系——这生命是短暂的，由生和死构成；因为凡不恨这些关系的，还不爱那没有生和死条件、且在地上婚姻中结合双方的生命。\par

41. 因此，若我问一位有妻子的好基督徒，即使他可能仍与她生育儿女，他是否愿意在那国度里有她；他当然记得天主的应许，以及那生命——那时这可朽坏的将穿上不可朽坏，这可死的将穿上不死；尽管目前因着爱的伟大，或至少某种程度的爱而犹豫，他会带着咒骂回答说他非常厌恶。我若再问他，是否愿意他的妻子在复活后、经历了应许给圣人的天使般的改变后，与他一同生活在那里，他会回答说他极其渴望，正如他厌恶前者一样。因此，一位好基督徒在同一个女人身上，爱天主的受造物——他渴望她得以改变和更新；但恨那朽坏和有死的婚姻关系与性交：也就是说，爱她身上属于人的特性，恨她作为妻子的属性。同样，他爱他的仇敌，不是作为仇敌，而是作为人；这样他愿同样的福祉临到仇敌如同临到自己，即愿他经过改正和更新而到达天国。对父亲、母亲和其他的血缘关系也应如此理解：我们恨他们身上属于人类在生和死中的命运，但爱那能与我们一同带到那领域的东西——在那里没有人说\Quote{我的父亲}，而是众人都对唯一天主说\BibleQuote{我们的父}；没有人说\Quote{我的母亲}，而是众人都对那另一个耶路撒冷说\BibleQuote{我们的母亲}；没有人说\Quote{我的兄弟}，而是每个人都对他人说\Quote{我们的兄弟}。但实际上，我们将作为一位新娘（当我们被聚集到合一中）与那位用自己的血将我们从这世界的玷污中救赎出来的主结婚。因此，基督的门徒必须恨那些在他人身上消逝的事物——他渴望同他们一起达到永恒常在的事物；而且他越爱他们本身，就越恨他们身上的这些事物。\par

42. 因此，一位基督徒可以和谐地与妻子生活，或者满足肉体的欲望（这是宗徒以准许而非命令说的），或者为了生儿育女（这在现在某种程度上可称赞），或者为了兄弟姊妹般的共融而无任何身体接触，有妻子如同没有妻子，如同在基督徒的婚姻中最卓越和崇高的那样：但如此对待她，即在她身上恨时间性关系之名，爱永恒福乐之望。因为我们毫无疑问恨那些我们至少希望将来不再存在的事物；例如，我们在现世的同一生命，若我们不恨它是暂时的，就不会渴望未来的生命，那生命不是由时间限定的。因为这生命由灵魂所取代，正如在那段话中说的：\BibleQuote{若有人不恨自己的灵魂，也不能作我的门徒。} 因为那腐朽的食物对此生命是必需的，主亲自说：\BibleQuote{难道灵魂不比食物更重要吗？} 即指食物所必需的生命。当主说祂将为羊群舍弃灵魂时，祂无疑指的是这生命，因为祂宣告祂将为我们而死。\par

% structure: p0061 source=h2 target=subsection reason=relative-heading-rank; base=h2

\subsection{第十六章}

43. 此处出现第二个问题：当主许可因奸淫的缘故休妻时，在本段中\Quote{奸淫}一词应在何种广义上理解——是理解为众人所共知的意义，即指在污秽行为中所犯的奸淫；还是按照圣经论及奸淫时的惯常用法（如前所述），指一切非法的败坏，诸如崇拜偶像、贪婪，因而自然也包括因其中所含非法情欲而犯的每样违法之事。但让我们征询宗徒，以免轻率发言。他说：\Quote{对于已婚者，我命令——其实不是我，而是主——妻子不可离开丈夫；但如果她离开了，就当持守独身，或与丈夫和好。}因为可能发生这样的情况：她离开正是出于主所许可的理由。或者，若女人有自主权因奸淫以外的其他理由休夫，而丈夫却没有此权，那么对于他后来所说的这句话：\Quote{丈夫也不可休妻}，我们当如何回答呢？为何他没有补充\Quote{除非为奸淫的缘故}（这是主所许可的），除非因为他希望我们理解为同样的规则：若他休妻（他因奸淫的缘故被许可如此行），就当不娶，或与妻子和好？因为丈夫与那样一个妇人（当无人敢于用石头砸她时，主曾对她说：\Quote{去吧，不要再犯罪了。}）和好，并非一件坏事。同样也因为，那位说\Quote{除非为奸淫的缘故，不可休妻}的主，若没有奸淫的情由，就强制他保留妻子；但若有奸淫，主并不强制他休妻，只是许可他，正如经上说：\Quote{女人在丈夫未死之前，不可另嫁；若在丈夫死前再嫁，便有罪；若丈夫死后不嫁，则无罪}，因为她没有被命令结婚，而只是被许可。因此，在所述的男女婚姻律法中，若有如此相似的规则，以致同一位宗徒不仅论到女人说：\Quote{妻子对自己的身体没有主权，丈夫却有}，也没有缄默不言，论到丈夫说：\Quote{同样，丈夫对自己的身体也没有主权，妻子却有}；那么，如果规则相似，就没有必要理解为女人休夫也仅限于奸淫的缘故，如同男人休妻一样。\par

44. 因此，应当思考我们应在何种广义上理解\Quote{奸淫}一词，并如我们开始所作的那样，征询宗徒。因为他接着说：\Quote{对其余的人，我说，不是主说。}此处，首先我们必须看出谁是\Quote{其余的人}，因为他先前是代表主向已婚者说话，如今却是以自己身份向\Quote{其余的人}说话：因而可能是指未婚者，但并非如此。因为他继续这样说：\Quote{若某弟兄有不信的妻子，而她乐意与他同住，就不可离弃她。}可见，他此时仍是对已婚者说话。那么，他说\Quote{其余的人}是何目的？除非他先前是对那些结合时双方在基督内的信德相同的人说话，而现在是对\Quote{其余的人}，即对结合时双方并非都信主的人说话？但他对他们说了什么？\Quote{若某弟兄有不信的妻子，而她乐意与他同住，就不可离弃她。若有女人有不信的丈夫，而他乐意与她同住，就不可离弃他。}因此，若他并非以主的身份颁下命令，而是以自己身份提出劝告，那么由此产生的好结果是：若有人行事相反，他并不违犯命令，正如他稍后论及贞女时所说，他没有主的命令，只提出他的意见；他如此赞美贞洁，凡愿意者皆可从中获益；但若他不这样做，也不能被判定为与命令相悖。因为有一类事是命令，另一类是劝告，还有一类是许可。妻子被命令不可离开丈夫；若离开，就当持守独身，或与丈夫和好：因此她不可有相反行为。但信主的丈夫被劝告，若他有不信的妻子且乐意与他同住，就不可离弃她：因此离弃她也是许可的，因为主并没有命令他不可离弃，只是宗徒的劝告：正如贞女被劝告不要结婚；但她若结婚，虽未遵循劝告，却未违反命令。许可出现于这句话：\Quote{但我说这话，是出于许可，并非出于命令。}因此，若不信的妻子可以被离弃（尽管不离弃更好），而按照主的命令，妻子不可被离弃，除非为奸淫的缘故，那么[由此]不信本身也就是奸淫。\par

45. 宗徒啊，你说什么？岂不是说，有不信的妻子的信主丈夫，若她乐意与他同住，就不应休她？他说，正是如此。因此，当主也命令不可休妻，除非为奸淫的缘故，为何你在此说\Quote{我说，不是主}？原因是：不信者所随从的拜偶像及其他所有有害的迷信，就是奸淫。主允许因奸淫的缘故休妻，但允许而非命令；祂给宗徒提供机会，劝告凡愿意的人不可休不信的妻子，以便或许她借此得以相信。因为他说：\Quote{因为不信的丈夫因妻子而成圣，不信的妻子因兄弟而成圣。}我想那时已有些妻子藉信主的丈夫接受信仰，有些丈夫藉信主的妻子接受信仰；虽未指名，他却用实例敦促，以加强他的劝告。接着他说：\Quote{否则你们的儿女就不洁净，但如今他们是圣洁的。}因为现在儿女是基督徒，他们因父母一方的请求或双方的同意而成圣；这若婚姻因一方信主而破裂，或若不忍受配偶的不信并给予相信的机会，便不会发生。因此，这是那位的劝告，我以为祂说了：\Quote{你额外花费的，我回来时必还给你。}\par

46. 此外，若不信是奸淫，拜偶像是不信，贪财是拜偶像，则无疑贪财也是奸淫。那么，如果贪财是奸淫，谁又能正确地将任何不法的欲情与奸淫范畴分开呢？由此我们看出，因着不法的欲情——不仅包括与别人配偶行污秽之举时所犯的，也包括任何不法的欲情，这些欲情使灵魂滥用身体而偏离天主的法律，遭受败坏和卑污的败坏——一个人可以无罪地休妻，妻子也可以无罪地休夫，因为主将奸淫的缘故作为例外；而根据上述考量，我们不得不将奸淫理解为普遍而包罗万象的。\par

47. 但当祂说\Quote{除非为奸淫的缘故}，祂并未指明是夫还是妻。因为不仅允许休犯奸淫的妻；而且凡休那甚至迫使他自己犯奸淫的妻的人，无疑是因奸淫的缘故休她。例如，若妻子强迫丈夫向偶像献祭，休这样妻子的人，不仅因她的缘故，也因自己的缘故休她：因她的缘故，因为她犯奸淫；因自己的缘故，免得他犯奸淫。然而，没有什么比一个人因奸淫休妻，而他自己也被证实犯奸淫更不公正了。因为想到这句话：\Quote{因为你论断别人，就定了自己的罪；你这论断人的，自己所行却一样。}因此，凡因奸淫而想休妻的人，应当先被清除奸淫之罪；对妇女我也作类似说明。\par

48. 关于祂所说的\Quote{凡娶被休的妇人就是犯奸淫}，可以问：被娶的妇人是否也与娶她的人一样犯奸淫？因为她也奉命独身，或与丈夫和好；但这是在她离开丈夫的情况下。然而，是她休夫还是被休有很大区别。若她休夫而另嫁，似乎出于改变婚姻关系的欲望离开前夫，这无疑是奸淫的心思。但若她被丈夫休弃——而她本愿与他在一起——那么娶她的人按主的宣告犯奸淫；但她是否也犯同样罪行则不明确——尽管很难发现，当男女同等地同意交合时，一个犯奸淫而另一个却不犯。此外，若他娶被休的妇人（她不是休夫而是被休）便犯奸淫，那么她使他犯奸淫，而这正是主所禁止的。由此我们推断，无论她被休或休夫，她必须独身，或与丈夫和好。\par

49. 又问：若妻子许可——或因她不能生育，或因她不愿行房——丈夫另娶一女，既非他人之妻，亦非离异者，他能否这样做而不犯奸淫？在旧约历史中有例证；但如今人类已过那阶段，领受了更大的诫命；为辨识那按最有序方式扶助人类的天主眷顾之施行世代的区别，应探究这些事，但不可作为生活规范来使用。但仍可问：宗徒所说\Quote{妻子对自己的身体没有主权，丈夫却有；同样，丈夫对自己的身体没有主权，妻子却有}，是否可以这样理解：在妻子（拥有丈夫身体主权）的许可下，丈夫可与另一非他人之妻亦非离异的女子交合？但这种意见不应采纳，免得似乎妻子在丈夫许可下也能行同样的事——而每个人的本能感觉都阻止这一点。\par

50. 然而，有时也会出现一些情况，妻子在丈夫同意下，似乎有义务为丈夫本人而这样做；例如，据说大约五十年前，在\Person{君士坦提乌斯}时代，在\Place{安提约基雅}就发生过这样的事。因为当时的行政长官\Person{阿辛狄努斯}（他曾一度担任执政官），在向某个公共债务人催讨一磅黄金时，不知出于什么动机，做了一件对于任何事都合法（或被认为合法）的官员来说常常危险的事——即以誓言和激烈的断言威胁说，如果对方不在他指定的某日前交出上述黄金，就处死他。因此，当那人被残酷地关押，无力摆脱那笔债务时，可怕的日子开始逼近。他碰巧有一位非常美丽的妻子，但没有钱来解救丈夫；这时，一个富人被这女子的美貌激起欲望，得知她丈夫处境危急，便派人告诉她，愿以一夜的代价（若她同意与他行淫）换取一磅黄金。她知道自己的身体不属于自己，而属于丈夫，便将消息告诉丈夫，说她愿意为丈夫这样做，但前提是丈夫本人——她身体在婚姻中的主人，一切贞洁所应归属于他——愿意她这样做，仿佛为了保全性命而处置自己的财产。丈夫感谢她，并命令她去做，绝不认为那是通奸的拥抱，因为要求这事的不是情欲，而是对丈夫的大爱，并且是出于他的命令和意愿。那女子来到富人的庄园，做了那好色之人所愿的事；但她只将自己的身体给了丈夫——他并不像平常那样渴望婚姻的权利，而是渴望生命。她得到了黄金；但那个给予者悄悄夺回了所给之物，换了一个装了泥土的类似袋子。然而，当女子回到家发现此事时，她冲出门外，要公开宣扬自己的所作所为，心中怀着迫使她那样做的同样的对丈夫的温柔爱意。她去找行政长官，坦白了一切，揭露了别人对她施行的欺诈。于是，行政长官首先宣判自己有罪，因为事情发展到这一步是由于他的威胁；然后，他仿佛在宣判别人的罪，决定从\Person{阿辛狄努斯}的财产中拿出一磅黄金缴入国库；而让那女子成为那块曾给她泥土代替黄金的土地的女主人。从这个故事中，我不作任何评价：各人可按自己的意愿作出判断，因为这历史并非来自神圣权威的典籍；但尽管如此，当这个故事被讲述时，人的直觉并不会像我们先前在没有任何事例而只陈述事情本身时所感到的厌恶那样，反对这女子在丈夫命令下的所作所为。然而，在\WorkTitle{福音}的这一段中，最需要牢记的是，奸淫的罪恶如此之大，以至于尽管已婚者之间以如此坚强的纽带相连，但只有这一个离婚的理由被排除；至于什么是奸淫，我们已经讨论过了。\par

% structure: p0070 source=h2 target=subsection reason=relative-heading-rank; base=h2

\subsection{第十七章}

51. 祂又说：\BibleQuote{你们又听见有对古人说的话：\InnerQuote{不可背誓，所起的誓，总要向主谨守。}我却告诉你们：什么誓都不可起。不可指着天起誓，因为天是天主的宝座；不可指着地起誓，因为地是祂的脚凳；也不可指着耶路撒冷起誓，因为耶路撒冷是大君王的京城。又不可指着你的头起誓，因为你不能使一根头发变白或变黑。你们的话该当是：是就说是，非就说非；其他多余的，便是出于邪恶。}\newline{}法利塞人的义是不背誓；这也由那命令不可起誓的祂所证实，就天国义的相关而言。因为正如那根本不说的人不可能说谎，同样那根本不起誓的人也不可能起假誓。然而，既然那呼求天主作证的人是在起誓，这一节必须仔细思考，免得宗徒似乎违背了主的诫命，他常常这样起誓，如他说：\BibleQuote{我现在写给你们的，看在天主面前，我没有说谎。}又说：\BibleQuote{我们的主耶稣基督的天主和父，那位永远受赞美的，知道我没有说谎。}类似性质的是那郑重的声明：\BibleQuote{天主是我所事奉的，我用我的心灵在祂圣子的福音中事奉祂，祂可以作证，我在祈祷中时常提到你们。}除非有人或许认为，只有在提到起誓的对象时才算是起誓；所以他没有用誓词，因为他没有说\Quote{指着天主}，而是说\BibleQuote{天主是见证。}这样想是可笑的；但由于那些好争论或理解迟钝的人，免得有人认为有区别，要知道宗徒也这样用了誓词，说：\BibleQuote{指着你们的夸耀，我是天天冒死。}不要有人以为这话是那样表达的，仿佛是说\Quote{你们的夸耀使我天天冒死}；正如有人说：\Quote{因他的教导，他成了博学}，即，因他的教导，他得以完全受教：希腊文本解决了这个问题，那里写道：\OriginalTerm{以你们的夸耀}{\GreekText{Νὴ} \GreekText{τὴν} \GreekText{καύχησιν} \GreekText{ὑμετέραν}}，这是只有起誓的人才会用的表达。因此，可以理解，主命令不可起誓是在这个意义上，免得有人把起誓当作好事而热衷追求，并因频繁使用起誓而陷于习惯性的发虚誓。因此，让那明白起誓不应算在好事中，而应算在必要事中的人，尽可能克制自己不去使用它，除非在必要时，当他看见人们因不易相信对他们有益的事，除非用起誓来保证。为此，有这样的话：\BibleQuote{你们的话该当是：是就说是，非就说非；}这是好的，是可取的。\BibleQuote{因为多余的，便是出于邪恶。}即，如果你被迫起誓，要知道这是出于必要，是因为你试图说服的那些人的软弱；这软弱当然是一种邪恶，我们每天祈祷从其中获救，当我们说：\BibleQuote{救我们免于凶恶。}因此祂没有说\Quote{多余的便是邪恶}；因为当你善用起誓时，你并没有行恶，起誓本身虽非善，却是必要的，为说服别人，你正试图为你某个有益的目的而打动他；而是\BibleQuote{出于邪恶}，是由于那人的软弱，使你被迫起誓。然而，除非亲身经历，没有人知道要戒除起誓的习惯，以及永不鲁莽地做那些必要时有时迫使他去做的事，是多么困难。\par

52. 但可能会问，为什么当祂说\Quote{我却告诉你们：什么誓都不可起}时，又加上\Quote{不可指着天起誓，因为天是天主的宝座}等，直到\Quote{不可指着你的头起誓}。我想原因是，犹太人认为，如果他们是凭着这类事物起誓，就不受誓言的约束；并且他们曾听见\Quote{所起的誓，总要向主谨守}，就认为如果指着天、地、耶路撒冷或自己的头起誓，就不受那向主的誓言的约束；这不是由于命令者的过失，而是因为他们没有正确理解。因此，主教导说，在天主的受造物中，没有一样是如此无价值，以致人可以认为能凭着它起假誓；因为从最高到最低的受造物，从天主的宝座开始到一根白头发或黑头发，都由天主眷顾所管理。祂说：\Quote{不可指着天起誓，因为天是天主的宝座；不可指着地起誓，因为地是祂的脚凳。}即，当你指着天或地起誓时，不要以为你的誓言不会使你受主的约束；因为你被证实是在指着那以天为宝座、以地为脚凳的那一位起誓。\Quote{也不可指着耶路撒冷起誓，因为耶路撒冷是大君王的京城。}这比祂说\Quote{我的城}更好的表达；虽然我们明白祂的意思是如此。而且，因为祂无疑是主，那指着耶路撒冷起誓的人便受誓言向主约束。\Quote{也不可指着你的头起誓。}现在，有谁能认为有什么东西比自己的头更属于自己呢？但我们怎能说头属于我们，当我们连使一根头发变白或变黑的能力都没有？因此，那甚至想指着自己的头起誓的人，也被誓言向天主约束，天主以不可言喻的方式掌管一切，并无所不在。这里也理解一切其他事物，当然不能一一列举；正如我们提到过的宗徒那句话：\BibleQuote{指着你们的夸耀，我是天天冒死。}并且为表明他受此誓言向主约束，他加了：\BibleQuote{这夸耀是在基督耶稣里。}\par

53. 然而（我这样说，是为了属血气的人），我们切不可认为，天堂被称为天主的宝座，大地被称为祂的脚凳，是因为天主如同我们坐下时那样，在天地之间安放了肢体；而是因为那个座位意指审判。由于在这整个宇宙的有机体中，天堂最为显赫，大地最为卑微——仿佛神圣的能力在优美卓越之处更加临在，但仍调节着最遥远、最低微的部分——故说祂坐在天上，践踏大地。但就属灵意义而言，天堂指圣善的灵魂，大地指有罪的人。既然属神的人判断一切，却不受任何人判断，他便合宜地被称为天主的座位；而罪人，经上对他说：\BibleQuote{你是土，还要归于土}，因为按那按各人功过分配合宜之物的公义，他被置于最低微的事物中，凡不愿留在法律之下的，便在法律之下受罚，便合宜地被视为祂的脚凳。\par

% structure: p0074 source=h2 target=subsection reason=relative-heading-rank; base=h2

\subsection{第十八章}

54. 现在，为总结这一段，有什么事比克服邪恶的习惯更为艰苦费力、更让信德的灵魂竭尽心力呢？这样的人应砍断阻碍天国的肢体，不要被痛苦压倒；在婚姻的忠贞中，他应忍受一切虽令人烦恼、却不涉及非法败坏之罪（即奸淫）的事：例如，若有人有妻子，或是不孕，或是身体畸形，或是肢体有缺陷——或瞎、或聋、或瘸，或有其他任何残疾——或为疾病、痛苦和衰弱所困，以及其他任何极其可怕的事（除奸淫外），他为了所应许的爱和婚姻的结合而忍受这一切；他不仅不应休掉这样的妻子，而且即使他没有这样的妻子，也不应娶一个被丈夫休弃的女人，即使她美丽、健康、富有、多产。如果这些事是不被允许的，那么靠近任何其他非法的拥抱就更不被视为合法了；他要这样逃避奸淫，使自己远离各种卑劣的败坏。让他说实话，不要以频繁的起誓，而是以道德的诚实来推荐它；至于那无数坏习惯的群众（其中只提几样以让人明白一切）起来反抗他，让他投靠基督徒争战的堡垒，仿佛从高处将他们打倒。但谁敢承担如此伟大的劳苦，除非那人被对义德的爱情所燃烧，以至于如同完全为饥渴所消耗，认为若不满足就没有生命，从而强取天主的国？否则，他将无法勇敢地忍受那些世俗之爱者视为劳苦、艰巨、完全难以摆脱恶习的一切事。所以，\BibleQuote{饥渴慕义的人是有福的，因为他们要得饱饫。}\par

55. 然而，当有人在这样的劳苦中遇到困难，在艰辛与崎岖中前进，被各种诱惑包围，感知到过去生活中的烦恼从四面八方涌现，害怕自己无法完成所承担的事时，让他热切地寻求忠告，以获得帮助。但除了这一条，还有什么忠告呢？即凡愿为自己的软弱获得天主帮助的人，应当承担他人的软弱，并尽可能帮助它。因此，让我们来看仁慈的诫命。温良的人和仁慈的人似乎是同一回事；但区别在于：温良的人（我们上面已说过）出于虔诚，不反驳那些针对他的罪过而提出的神圣判决，也不反驳那些他尚未理解的天主的话；但他对他所不反驳、不抵抗的人并不施加恩惠。而仁慈的人以这样的方式不抵抗，乃是为了纠正他因抵抗而会变得更坏的人。\par

% structure: p0077 source=h2 target=subsection reason=relative-heading-rank; base=h2

\subsection{第十九章}

56. 因此，主继续说：\BibleQuote{你们一向听说过：以眼还眼，以牙还牙。我却对你们说：不要抵抗恶人；而且，若有人掌击你的右颊，你把另一面也转给他。有人控告你，要拿你的内衣，你连外衣也让他拿去。有人强迫你走一千步，你就同他走两千步。求你的，就给他；向你借贷的，你不要拒绝。}法利塞人的较小义德是不在报复中超过限度，使谁也不还击超过他所受的；这是一大步。因为不容易找到一个人，当他挨了一拳时，只想还那一拳；当听到辱骂者的一句话时，满足于只回一句，且是相等的；而是更加过分地报复，或是由于愤怒的激动，或是因为他认为公正：先施害者应比未施害者受更重的伤害。这种精神在很大程度上被法律所约束，经上写道：\BibleQuote{以眼还眼，以牙还牙}；这些话指示一种尺度，使报复不超过伤害。这是和平的开端；但完全的和平是根本不想这样的报复。\par

57. 因此，在第一种超出法律的行为（即以较大的恶回报较小的恶）与主为了成全门徒而揭示的这一行为（即不以恶报恶）之间，有一条中间道路占有一定位置，即：偿还多少就收到多少；通过这一规定，按时代的安排，从最高的不和过渡到最高的和谐。请看因此，那些首先伤害他人、怀着损害和伤害欲望的人，与即使受到伤害也不报复的人，相隔多么遥远。然而，那人不首先伤害任何人，但受伤害时却以更大的伤害回报，无论是意愿还是行为，他已从最大的不义中退后，向最高的义迈进如此之多；但他仍未完全遵守梅瑟所颁布的法律所命令的。因此，偿还恰好与他所收到的一样多的人，已经宽恕了一些：因为伤害的一方并不应得与被他伤害的无辜者所遭受的同样多的惩罚。因此，这种不完全、绝不严厉、反而仁慈的正义，被那来成全法律而非废除法律者所完善。因此，还有两个中间步骤，祂留下让人理解，而祂选择谈论仁慈的最高发展。因为还有人可能这样做：他未能完全达到属于天国的诫命的那种高度；他的行为方式不是偿还同样多，而是更少；例如，以一拳代替两拳，或为被挖出的一只眼而割下一只耳朵。那超越此层、什么都不偿还的人，接近主的诫命，但仍未达到。因为在主看来，如果你对所受到的恶不行恶，还不够，除非你准备好接受更多。因此祂并没有说：\Quote{但我告诉你们，}你们不可以以恶报恶；即使这将是一条伟大的诫命；但祂说：\Quote{你们不要抵抗恶人；}这样，你不仅不要偿还可能加在你身上的，甚至不要抵抗其他的打击。因为这也是祂继续解释的：\Quote{凡有人打你的右脸，把另一面也转给他；}因为祂没有说，如果有人打你，你不要想打他；而是说，如果他要继续打你，你就把自己进一步献给他。关于同情，那些照料他们所深爱的人，如同自己的孩子，或在病中非常亲爱的朋友，或小孩，或疯子，这些人经常承受许多；如果他们的福祉需要，他们甚至准备承受更多，直到年老或疾病的软弱过去。同样，对于主——灵魂的医生——教导要关心邻舍的人，祂还能教他们什么呢，除了安静地容忍他们希望关心的那些人的软弱？因为一切邪恶都源于心灵的软弱：因为没有比在美德上完美的人更无害的了。\par

58. 但或许会问右颊是什么意思。因为我们发现在最可靠的希腊文抄本中正是这样读的，尽管许多拉丁文抄本只有\Quote{面颊}一词，而没有加上\Quote{右}。面容是每个人藉以被认出的；在宗徒的著作中我们读到：\Quote{你们若有人奴役你们，若有人吞噬你们，若有人侵占你们，若有人高举自己，若有人打你们的脸，你们竟然容忍！}紧接着他又说：\Quote{我说话如同出于羞辱，}他以此解释击打面容是指什么，即被藐视和受轻看。宗徒说这话并非要他们不容忍那些人，而是要他们更容忍他自己，因为他是如此爱他们，甚至愿意为他们耗尽自己。但由于面容不能称为右和左，但按照天主的评价和按照这个世界的评价，可能有一种价值，它被分派为右颊和左颊，以致基督的门徒若因作基督徒而必须受辱，倘若他拥有任何属于这世界的尊荣，就该更准备好在自身内承受羞辱。因此，这位同一位宗徒，若对他在这世界所拥有的尊严保持沉默，当人因他基督徒的名号而迫害他时，他就不会把另一面颊转向那击打右颊的人。因为当他说\Quote{我是罗马公民}时，他并非不愿意在那他视为最小的事上，被那些曾蔑视他内如此宝贵和赋予生命的名号的人所藐视。因为他后来难道因此较不屈服于锁链——那原是法律不准加在罗马公民身上的？或者他愿控告任何人这等伤害？若有人因罗马公民之名而宽免了他，他却并不因此不提供他们可击打的对象，因为他愿藉自己的忍耐医治那些他看见在他身上尊敬属于左而非右肢体的人的如此大悖逆。因为只需注意这一点：他做一切事怀着怎样的精神，对那些使他受苦的人，他行事多么仁慈和温和。因为当他奉大司祭之命被用手掌击打时，他肯定说：\Quote{天主将要打击你，你这粉白的墙！}这话在不明白的人听来似乎是侮辱，但在明白的人，却是一个预言。因为粉白的墙就是伪善，即一种矫饰，将司祭尊荣展现在前，而在此名号下，如同在白色遮盖下，隐藏内在的、可说是污秽的卑鄙。因为当人对他说：\Quote{你辱骂大司祭吗？}他回答：\Quote{弟兄们，我不知道他是大司祭；因为经上记载：不可诋毁你百姓的官长。}他奇妙地保守了属于谦卑的一切。在此他显示他说话时是多么平静，尽管他看似愤怒地说了那话，因为他回答得如此迅速和温和，这是那些愤怒和慌乱的人所不能做到的。而在这句话中，他对那些理解他的人说了实话：\Quote{我不知道他是大司祭。}如同在说：我认识另一位大司祭，我为祂的名受这些苦，辱骂祂是不允许的，而你们正在辱骂祂，因为在你们仇恨我的，不过是祂的名号。因此，人必须不以伪善的方式夸耀这些事，而是在心中准备好一切，以便他能咏唱那句先知的话：\Quote{我的心已准备妥当，天主，我的心已准备妥当。}因为许多人学会了献上另一面颊，却不知道爱那击打他们的人。但事实上，主自己——祂确实是首先实践祂所教导的诫命——当大司祭的仆人击打祂的面颊时，并没有献上另一面颊；反而说：\Quote{我若说得不对，你指证不对的地方；若说得对，你为什么打我？}然而，祂的心并非因此没有准备好，为了众人的得救，祂不仅准备被击打另一面颊，甚至准备整个身体被钉在十字架上。\par

59. 因此，接下来的话：\Quote{若有人愿告你，拿你的内衣，你连外衣也让他拿去。}应正确理解为一条关乎内心准备的诫命，而非外在行为的表现。但关于内衣和外衣所说的话，不仅应在这些事上实行，也应在每一样我们按任何权利称为我们暂时拥有的东西上实行。因为若这条命令是针对必需品而给予的，那么我们更应轻视那些多余之物！然而，那些我称为属于我们的东西，都应按主亲自给予诫命的那种范畴来理解，当祂说：\Quote{若有人愿告你，拿你的内衣。}因此，所有那些可能被控告并使其权利从我们转给控告者或他所代表的人的东西，都应如此理解；例如，衣服、房屋、地产、驮兽，总之各类财产。但至于是否也包括奴隶，这是一个重大问题。因为基督徒不应像拥有一匹马或金钱那样拥有奴隶：虽然可能一匹马比一个奴隶更值钱，而某些金器银器比两者都贵得多。但关于那个奴隶，如果他被其主人教育和管理，方式比想要带走他的人更加正直、更体面、更促进敬畏天主，我不知道是否有人敢说，他应该被像一件衣服那样被轻看。因为人应当爱近人如己，因为他被万有的主命令（如随后所示）也要爱自己的仇敌。\par

60. 必须细心留意：每一件内衣固然是一件衣服，但并非每一件衣服都是一件内衣。因此，\Quote{衣服}一词比\Quote{内衣}一词含义更广。所以我认为祂如此表述：\BibleQuote{若有人控告你，要拿你的内衣，就让他连外衣也拿去}，就如同祂说：无论谁想拿你的内衣，就把你其余的衣服都给他。所以有些人将\Emph{pallium}一词解释为希腊文此处所用的\OriginalTerm{外衣}{\GreekText{ἱμάτιον}}。\par

61. 祂说：\BibleQuote{有人强迫你走一里路，你就同他走两里}。这当然不是指你实际要步行那么远，而是指你在心意上准备好这样做。因为在具有权威的基督传记中，你找不到圣人们或主亲自（当祂屈尊降生成人，以人性生活，给我们树立生活榜样时）这样做过；同时，几乎所有地方你都会看到他们准备好平静地忍受任何被不义强加的事。但我们是否应该认为这句话只是为表达\BibleQuote{同他走两里}而说？还是祂更愿意使三里之数得以完全——此数具有完美的含义，从而让每个人在这样做时记住，他是通过怜悯地承受他所愿治愈之人的软弱来履行完美的义？为此祂似乎也用三个例子来推荐这些训诫：第一，若有人打你的脸颊；第二，若有人要拿你的外衣；第三，若有人强迫你走一里；在这第三个例子中，在原有的一里上加了一倍，如此便完成了三个一组。如果这段经文中的这个数字（如前所述）不代表完美，那么可以这样理解：祂在颁布训诫时，仿佛从较可容忍的开始，逐步推进，直到达到忍受加倍的程度。因为首先，祂希望在右颊被打后，再转左颊，这样你就准备好忍受比你已受的为轻的事。因为无论\Quote{右}代表什么，总比\Quote{左}所代表的更加珍贵；如果一个人已经在更珍贵的事上忍耐了，那么在较不珍贵的事上也忍耐，就是较轻的事。其次，对于想要拿人外衣的人，祂吩咐连外衣也给他：这要么同样多，要么不更多；但不是加倍。第三，关于一里路，祂说加两里，吩咐你要忍受甚至加倍之多：这表明无论恶人想从你这里拿走的东西比你原有的少一点、同样多还是更多，你都应以平静的心忍受。\par

% structure: p0084 source=h2 target=subsection reason=relative-heading-rank; base=h2

\subsection{第20章}

62. 的确，在这三类例子中，我看没有什么伤害被遗漏了。因为我们所遭受的一切不义之事可分为两类：一类是不可弥补的，另一类是可弥补的。但在不可弥补的情况下，通常寻求报复性补偿。被打后还手有何益处？受伤的身体部位会因此恢复原状吗？然而激动的灵魂渴望这样的慰藉。但健康坚固的灵魂并不以此为乐；反倒认为别人的软弱应被怜悯地忍受，而非藉惩罚他人来抚慰自己的（根本不存在的）软弱。\par

63. 我们并非因此被禁止施加那种有益于改正、且由怜悯本身所要求的惩罚；这也不妨碍我们准备为所欲纠正之人的缘故忍受更多。但只有那种因爱德之大而克服了报复者惯常的仇恨的人，才适合施加这种惩罚。因为不必担心父母会恨自己的小孩：当他犯错时，父母打他，使他不再继续犯错。当然，爱德的完美是通过效法圣父本身而在下面的话中展示给我们的：\BibleQuote{爱你们的仇敌，善待恨你们的人，为迫害你们的人祈祷}；然而先知论及祂时说：\BibleQuote{主所爱的，祂必管教；祂鞭打所收纳的每一个儿子}。主也说：\BibleQuote{那不知道主人意愿而做了该打之事的仆人，必少受责打；那知道主人意愿而做了该打之事的仆人，必多受责打}。因此，只要求那按事物自然秩序被赋予权力的人施行惩罚，并且像父亲对自己尚年幼而无法怀恨的孩子那样怀着同样的善意去惩罚。因为由此引出最恰当的例子，足以表明罪可以用爱德来惩罚，而非放任不管；这样，施行惩罚的人希望受罚者不是因惩罚而悲惨，而是因改正而幸福，同时准备（若有需要）平静地忍受更多由他所愿纠正之人施加的伤害，无论他是否有能力约束那伤害者。\par

64. 但是，伟大而圣洁的人，尽管他们当时极好地知道，那种使灵魂与身体分离的死亡不足畏惧，却依照那些可能畏惧死亡者的情感，以死亡惩罚某些罪行，既因活着的人被一种有益的恐惧所震慑，也因死亡本身不会伤害那些受到死亡惩罚的人，而是罪——若他们继续活着，罪会增多。他们并非鲁莽地判断，因为天主已将如此判断的权力赐给了他们。因此，厄里亚亲手并呼求天火下降而处死了许多人；其他许多伟大而神圣的人，也出于同样关怀人类福祉的精神，无私地这样做。当门徒们引用了这位厄里亚的例子，向主提及他所做的事，希望能赐给他们自己权力，从天降火消灭那些不接待主的人时，主责备了他们，不是责备圣先知的榜样，而是责备他们在报复方面的无知——他们的知识尚属初阶；主看出他们并非以爱德渴望纠正，而是以仇恨渴求复仇。因此，当主教导了他们爱近人如己的道理，并当圣神倾注之后——主在升天十天后，按所应许的，从上派遣了圣神——并不缺少这样的报复行为，虽然比旧约中稀少得多。因为在旧约中，他们多半是仆人，被恐惧所压制；而在新约中，他们多半是自由人，被爱德所滋养。因为，正如我们在\BibleBook{宗徒}{大事录}中读到，连宗徒伯多禄的话也使阿纳尼雅和他的妻子仆倒死去，没有再活过来，而是被埋葬了。\par

65. 然而，如果那些反对旧约的异端者不相信这卷书，让他们默观宗徒保禄罢——他的著作他们和我们一起诵读——他谈及某个罪人时说：\Quote{使这个人交与撒殚，败坏他的肉体，为使他的灵魂得救。}如果他们在那里不愿理解为死亡（也许这不明确），那么让他们承认，宗徒通过撒殚施加了某种惩罚；而他这样做并非出于仇恨，而是出于爱德，这从附加的话\Quote{为使他的灵魂得救}可以清楚看出。或者让他们注意我们在那些他们自己极为推崇的书中所说的——那里记载说，宗徒多默曾为一个用手掌打了他的人咒求一种极其残酷的死亡刑罚，同时却把他的灵魂交托给天主，使它在来世得以幸免——那人被狮子杀死后，狗把他那与身体分离的手叼到宗徒宴席的桌前。我们不信这记载是允许的，因为它不在大公教会正典之中；然而，那些对旧约中的肉体惩罚极其猛烈地咆哮（不知何故盲目）的人，却既诵读这记载，又尊崇它为完全纯正、完全真实，他们全然不知这些惩罚是在什么精神下、在时代有序分派的哪个阶段施加的。\par

66. 因此，在这类须以惩罚赎补的伤害中，基督徒将持守这样的尺度：遭受伤害时，心思不会上升到仇恨，反而出于对被伤害懦弱的怜悯，甚至愿意承受更多；同时也不忽略纠正，这纠正可以用劝告、权威或权力来施行。另有另一类伤害，可以完全补偿，其中有两种：一种涉及金钱，另一种涉及劳务。因此，附有例子：关于外套和长袍的例子属于前者，关于强制服役一里和两里的例子属于后者；因为一件衣服可以还回，你所帮助过的人也可以在必要时帮助你。除非这种区分更应如此理解：第一个假设的情况（关于被打的脸颊）指所有由恶人施加的伤害，其补偿只能通过惩罚；第二个假设的情况（关于衣服）指所有无需惩罚即可补偿的伤害；因此，也许附加了\Quote{如果有人控告你}，因为通过判决被夺去的东西，不被认为带有如此程度的暴力以致应受惩罚；而第三种情况则兼有两者，使补偿既可在无惩罚也可在有惩罚下进行。因为那人暴力索要非分之劳，却未经任何法律程序，如同那邪恶地强迫人同他走、以不法手段迫使不愿者帮助他的人，他既能为其邪恶付出代价，也能偿还劳务——如果受屈者再次要求的话。因此，在这所有类别的伤害中，主教导说，基督徒的态度应当极其忍耐和怜悯，并充分准备承受更多。\par

67. 但是，仅仅不伤害人是一件小事，除非你也尽可能地施予恩惠，因此祂接着说：\BibleQuote{凡求你的，就给他；有人向你借贷的，不要转身离去。} \BibleQuote{凡求你的，}祂说；并不是说，凡他所求的一切都给他。如此，你应当给予那些你能够诚实且公正地给予的东西。因为假如他求取金钱，并企图用它来压迫一个无辜的人，那会怎样呢？简言之，假如他求取某种不洁之物，那又当如何呢？但不必列举许多例子，事实上数不胜数；可以肯定的是，应当给予那既不伤害你自己也不伤害对方的东西，只要人能知道或推测；至于你公正地拒绝了他的请求的人，应将公义本身告诉他，从而不让他空手离去。如此，你便给予每一个求你的人，尽管你并不是总给予他所求的；当你纠正了那个提出不正当请求的人时，你有时会给予更好的东西。\par

68. 然后，关于祂所说的：\BibleQuote{有人向你借贷的，不要转身离去，}这应归诸于心态；因为天主喜爱乐意捐助的人。此外，每一个接受某种东西的人都是借贷者，即使他自己不打算偿还；因为既然天主更加报答怜悯人的人，那么凡行善的人就是以利息借贷。或者，如果除了将借贷者理解为接受某物并打算偿还的人之外，似乎没有更好的理解，那么我们就必须明白，主已经包括了两种施恩的方式。因为我们要么是出于行善而赠送礼物，要么是借给将偿还我们的人。而许多人虽然将神圣的赏报摆在面前，准备赠送礼物，却变得迟疑，不愿答应借贷的请求，仿佛他们命中注定从天主那里得不到任何回报，因为接受者已经偿还了所给的东西。因此，神圣的权威正确地劝勉我们这种施恩的方式，说道：\BibleQuote{并且，有人向你借贷的，不要转身离去：}\Emph{即}，不要将你的善意与那祈求的人疏远，因为你的金钱将毫无用处，而且天主不会报答你，因为那人已经偿还了；但当你出于对天主诫命的尊重而这样做时，对于颁布这些诫命的那一位来说，这不会不结出果实。\par

% structure: p0092 source=h2 target=subsection reason=relative-heading-rank; base=h2

\subsection{第二十一章}

69. 接着，祂继续说：\BibleQuote{你们曾听到说过：你应爱你的近人，恨你的仇敌；我却对你们说：爱你们的仇敌，善待恨你们的人，为迫害你们的人祈祷；这样你们才能成为你们在天之父的子女；因为祂使自己的太阳升起，光照恶人也光照善人，降雨给义人也给不义的人。所以，如果你们只爱那些爱你们的人，你们有什么赏报呢？连税吏不也是这样作的吗？如果你们只问候你们的兄弟，你们做了什么比人更多的事呢？连外邦人不也是这样作的吗？所以，你们应当是完美的，如同你们在天之父是完美的一样。}因为没有这种爱，即我们被命令连仇敌和迫害者也要爱的爱，谁能完全实行那些上文所提及的事呢？此外，那种怜悯的圆满——灵魂在困苦中最被顾及的怜悯——不能超越爱仇敌的境界；因此结束的话是：\BibleQuote{所以，你们应当是完美的，如同你们在天之父是完美的一样。}但应这样理解：天主作为天主是完美的，而灵魂作为灵魂是完美的。\par

70. 然而，法利塞人的义德确实有一个进阶，属于旧法律，这可以从以下考虑看出：许多人也恨那些爱他们的人；例如，放荡的子女恨他们的父母，因为父母阻止他们放荡。因此，那爱近人的人，尽管他还恨仇敌，已上升了一个阶梯。但在那位来成全法律而非废除法律者的国度里，他将在仁慈和善意上达到圆满，当他将之推展到爱仇敌时。因为先前的阶段虽然也算是一点进步，却如此微小，连税吏也能达到。至于法律中所说的：\BibleQuote{你要恨你的仇敌，}这不应被理解为对义人发出的命令，而更应该是对软弱之人的许可。\par

71. 这里确实产生了一个不可忽视的问题：对于主的这条诫命，即祂劝勉我们爱仇敌、善待恨我们的人、为迫害我们的人祈祷，圣经的其他许多部分，在那些不够勤勉和冷静地思考的人看来，似乎与之相悖；因为在先知书中，有许多对仇敌的诅咒，被认为是咒骂：例如，那一句：\BibleQuote{愿他们的餐桌成为网罗，}以及那里所说的其他事；还有那一句：\BibleQuote{愿他的儿女成为孤儿，他的妻子成为寡妇，}以及同一位圣咏的作者，在同一篇圣咏中，关于犹达斯的事，在他之前或之后所说的其他话。圣经各部分还有许多其他陈述，可能既与主的这条诫命相悖，也与宗徒的那条诫命相悖，那诫命说：\BibleQuote{你们要祝福，不可诅咒；}而关于主，也记载祂诅咒了那些不接受祂话的城市；而且上述宗徒也这样谈到某个人：\BibleQuote{主要照他的行为报应他。}\par

72. 但这些困难很容易解决，因为先知是通过诅咒来预言将要发生的事，这不是在祈求他所愿望的，而是出于预知者的精神。\OriginalTerm{诅咒}{imprecation}同样，主也是这样，宗徒也是这样；尽管在他们的言辞中，我们也没有发现他们愿望了什么，而是他们预言了什么。因为当主说：\BibleQuote{祸哉，你，\Place{葛法翁}！}祂不是别的，只是说某种灾祸会因她的不信而临到她；而这事的发生，主并非\OriginalTerm{恶意地}{malevolently}愿望，而是凭祂的神性预见到了。宗徒没有说\Quote{愿主报答}，而是说：\BibleQuote{主要照他的工作报答他}；这是预言者的言辞，而非发出诅咒者的言辞。同样，关于我们已经谈过的犹太人的伪善，他看见他们的毁灭即将来临，就说：\BibleQuote{天主必要打你，你这粉饰的墙。}但先知们尤其习惯于以发出诅咒的形式预言未来的事件，正如他们常常以过去时的形式预言将要发生的事一样：例如这段经文：\BibleQuote{外邦为什么叛乱，万民为什么妄想？}他没有说\Quote{外邦将要叛乱，万民将要妄想？}虽然他并不是在提及这些事如同已成过去，而是期待它们尚未来到。同样还有这段：\BibleQuote{他们分了我的衣裳，为我的内衣拈阄。}因为这里他也没有说\Quote{他们将分我的衣裳，将为我的内衣拈阄。}然而，除了那些不明白言谈方式的变化丝毫不减少事实的真实性，而且大大增强了我们心中的印象的人以外，没有人指责这些言辞。\par

% structure: p0097 source=h2 target=subsection reason=relative-heading-rank; base=h2

\subsection{第二十二章}

73. 但我们面前的问题因\Person{若望}宗徒所说的而变得更加紧迫：\BibleQuote{谁若看见自己的弟兄犯了不至于死的罪，就应当祈求，主必将生命赐给他，就是给那些犯不至于死的罪的人。有至于死的罪，我不说要为他祈求。}因为他清楚地表明，有某些弟兄，我们不被命令为他们祈祷，尽管主命令我们甚至为迫害我们的人祈祷。而且，除非我们承认，在弟兄中有些罪比敌人的迫害更严重，否则手头的问题就无法解决。此外，\Quote{弟兄}指基督徒，这可以用圣经中的许多例子证明。然而，最清楚的是宗徒所说的：\BibleQuote{因为不信主的丈夫因妻子而成圣，不信主的妻子因\Emph{弟兄}而成圣。}他没有加上\Emph{我们的}一词，但认为这是显而易见的，因为他希望用\Emph{弟兄}一词来指代一位有不信主妻子的基督徒。因此，稍后他说：\BibleQuote{但若不信主的人离去，就由他离去：在这种情形下，弟兄或姊妹不受束缚。}因此，我认为一个弟兄的罪至于死，是指当一个人通过我们的主耶稣基督的恩宠认识了天主之后，攻击弟兄团体，并被嫉妒的火焰所驱使，反对那使他与天主和好的恩宠本身。但罪不至于死，是指若有人没有撤回对弟兄的爱，而是因性情的某些软弱未能履行弟兄应尽的义务。为此，我们的主在十字架上说：\BibleQuote{父啊，宽恕他们，因为他们不知道他们做什么。}因为他们尚未成为圣神恩宠的分享者，还没有进入圣洁的弟兄团契。在\WorkTitle{宗徒大事录}中，圣\Person{斯德望}为那些用石头砸他的人祈祷，因为他们尚未相信基督，也没有反对那共同的恩宠。因此，我相信\Person{保禄}宗徒没有为\Person{亚历山大}祈祷，是因为他已经是一个弟兄，并且犯了至于死的罪，即因嫉妒攻击弟兄团体。但对于那些没有断绝爱，而是因惧怕而退却的人，他祈求他们得到宽恕。因为他这样表达：\BibleQuote{铜匠\Person{亚历山大}加给了我许多恶事；主要照他的作为报答他。你们也要提防他，因为他极力反对我们的话。}然后他加上为谁祈祷，这样说：\BibleQuote{在我初次辩护时，没有人支持我，众人都离弃了我；愿天主不将这罪归于他们。}\par

74. 正是他们罪过的这一区别，将出卖者\Person{犹达斯}与否认者\Person{伯多禄}分开：这并不是说一个悔改者不应被宽恕，因为我们不可与主的声明冲突，祂命令当弟兄请求弟兄宽恕时，应当宽恕他；而是因为与那罪过相连的毁灭如此之大，以至于他无法忍受祈求宽恕的卑下，即使被良心所迫而承认并揭露自己的罪。因为当\Person{犹达斯}说：\BibleQuote{我犯了罪，因为我出卖了无辜的血。}然而，对他而言，在绝望中跑去上吊比在谦卑中祈求宽恕更容易。因此，知道天主宽恕哪种悔改是非常重要的。因为许多人更容易承认自己犯了罪，并且对自己如此愤怒，以至于强烈希望自己没有犯罪；但他们却不肯谦卑心灵、使之痛悔并恳求宽恕；我们必须认为他们有这样的心态，因为他们因自己罪过的重大而感觉自己已被定罪。\par

75. 这也许就是干犯圣神的罪，\Emph{即}在领受圣神的恩宠后，因恶意和嫉妒而做出与兄弟之爱相违背的事——我主说，这种罪不论在今世或来世都不得赦免。因此，可以问：当犹太人说我主是靠魔王贝耳则步驱魔时，他们是否干犯了圣神？我们应当理解为这是针对我主自己说的，因为祂在另一处论及自己说：\BibleQuote{如果家主被称为贝耳则步，何况祂的家人呢？} 还是说，由于他们出于极大的嫉妒，对如此明显的恩惠忘恩负义，尽管他们还不是基督徒，但因其嫉妒之深，我们可否认为他们干犯了圣神？后一种解释肯定不能从我主的话中得出。因为虽然祂在同一处说：\BibleQuote{凡说话干犯人子的，可得赦免；但若说话干犯圣神的，不得赦免，无论今世或来世。} 然而，祂劝诫他们的目的似乎是让他们来到祂的恩宠中，并在领受恩宠后不再像现在这样犯罪。因为他们现在说了干犯人子的话，如果他们悔改、信从祂并领受圣神，这罪可得赦免；但如果在领受圣神后，他们仍选择嫉妒兄弟情谊，并攻击已领受的恩宠，那么这罪无论今世或来世都不得赦免。因为如果祂认为他们已定罪到毫无希望的地步，祂就不会认为他们仍应受劝诫，正如祂接下来所说：\BibleQuote{或者使树好，它的果子也好；或者使树坏，它的果子也坏。}\par

76. 因此，应当明白：我们应当爱仇敌，善待恨我们的人，为迫害我们的人祈祷；同时也要明白，对于某些弟兄的罪，我们没有被命令为他们祈祷；免得因我们缺乏技巧，而使圣经似乎自相矛盾（这是不可能发生的）。但是，是否正如我们不为某些人祈祷，我们同样也要反对某些人，这一点尚不十分明确。因为经上普遍说：\BibleQuote{祝福，而不诅咒；} 又说：\BibleQuote{不以恶报恶。} 再者，你不为某人祈祷，并不因此就反对他；因为你可以看到他的惩罚是确定的，他的救恩完全无望；你不为他祈祷，不是因为你恨他，而是因为你感到你对他毫无益处，并且你不希望你的祈祷被最公义的审判者拒绝。但是，关于那些人，我们有启示说圣人们曾为他们祈祷，但不是为了使他们从错误中转回（因为这样祈祷是为他们好），而是为使最终的定罪临到他们：不是像先知为我主的出卖者所做的祈祷；因为那是预言未来，而不是希望惩罚；也不是像宗徒对亚历山大的祈祷；关于这一点已经说得够多了；而是像我们在 \BibleBook{若望}{默示录} 中读到殉道者祈祷得到报复；而那位著名的首位殉道者却为那些用石头砸他的人祈祷，求主赦免他们。\par

77. 但我们不必为此情况所动摇。因为关于那些穿白衣的圣人，当他们祈求得报应时，有谁敢断言他们是反对那些人本身，还是反对罪恶的统治权？因为罪恶统治权的推翻本身就是对殉道者真正的报应，且充满公义与仁慈，正是在这个统治权下，他们遭受了如此巨大的痛苦。宗徒也为了推翻它而奋斗，说：\BibleQuote{所以不要让罪恶在你们必死的身体上为王。} 但罪恶的统治权被摧毁和推翻，部分是通过人的悔改，使肉体服从心神；部分是通过对坚持犯罪之人的定罪，使他们受到公义的处置，以致不能再扰乱与基督一同为王的义人。请看宗徒保禄；他这么说：\BibleQuote{所以我这样奔跑，不是如同无定向的；我这样打拳，不是如同打空气的；而是痛击我身，使它屈服。} 难道你在其中不看到他以自身报复了殉道者斯德望？因为他确实在打倒、削弱、屈服并约束自己身上那曾使祂迫害斯德望和其他基督徒的原则。那么，谁能证明圣殉道者们不是向主祈求这样的一种报应呢？同时，为了得以报复，他们也可合法地盼望这个世界的终结，因为他们曾在此世承受了如此多的殉道之苦。如此祈祷的人，一方面为那些可能治愈的敌人祈祷，另一方面并不反对那些选择不治的人；因为天主在惩罚他们时，并非恶意的折磨者，而是最公义的安排者。因此，毫不犹豫地，让我们爱我们的仇敌，善待恨我们的人，并为迫害我们的人祈祷。\par

% structure: p0103 source=h2 target=subsection reason=relative-heading-rank; base=h2

\subsection{第二十三章}

78. 然后，关于接下来的这句话：\BibleQuote{好使你们成为你们在天之父的子女}，应当按照\Person{若望}所说的这句\BibleQuote{他赐给他们权能，成为天主的子女}的规则来理解。因为一位是按本性为子的，他丝毫不知罪；但我们因接受权能而成为儿子，在程度上我们实行祂所命令我们的事。因此，宗徒的教导将我们蒙召承受永生的产业称为\OriginalTerm{义子}{adoption}，好使我们与基督\OriginalTerm{同为承继者}{joint-heirs with Christ}。所以，我们藉着\OriginalTerm{属神的重生}{spiritual regeneration}成为儿子，并被收养进入天国，不是作为外人，而是作为祂所造所创的：因此，一个是祂以全能创造我们（我们原本虚无）的恩惠，另一个是祂收养我们，使我们作为儿子，因分享而能与祂一起享有永生。所以，祂并不说：你们做这些事，因为你们是儿子；而是说：你们做这些事，为要成为儿子。\par

79. 但祂藉着\OriginalTerm{独生子}{Only-begotten}亲自召叫我们相似祂时，祂召叫我们达致祂自己的肖像。因为祂，正如接下来所说的，\BibleQuote{他使太阳上升，光照恶人，也光照善人；降雨给义人，也给不义的人}。你要理解祂的太阳不是那肉眼可见的，而是那智慧，经上论及她说：\BibleQuote{她是永恒之光的辉映}；经上也说：\BibleQuote{义德的太阳在我身上升起}；又说：\BibleQuote{但为你们敬畏上主名号的人，义德的太阳将要升起}；这样你也当理解那雨水是用\OriginalTerm{真理的教导}{doctrine of truth}来浇灌，因为基督向善人也向恶人显现，并为善人也为恶人所宣讲。或者你宁愿理解那摆在人眼前，不仅人，还有牲畜，可见的太阳，以及那产生果实的雨水，这些果实是为滋养身体的，我认为后一种解释更可能：这样，那\OriginalTerm{属神的太阳}{spiritual sun}只升在善人和圣人身上；因为正是这件事，恶人在那称为\BibleBook{撒罗满}{智篇}的书中哀叹：\BibleQuote{太阳没有在我们身上升起}；而那\OriginalTerm{属神的雨水}{spiritual rain}只浇灌善人；因为恶人是指那葡萄园，经上说：\BibleQuote{我也命令云不再降雨在上面}。但无论你理解哪一种，都是出于天主的伟大良善，我们受命要效法，如果我们愿意成为天主的子女。因为谁如此忘恩，以至于不感到那可见的光和物质的雨在这现世带来何等大的安慰？我们看见这安慰在此生同样赐予义人和罪人。但祂没有说\BibleQuote{他使太阳上升，光照恶人，也光照善人}；而是加上了\BibleQuote{祂的}一词，\Emph{即}祂自己所造所立的，祂在造它们时没有从任何人取什么，正如\BibleBook{}{创世纪}中关于一切光体所记载的；祂能合宜地说，凡祂从无中创造的万物都是祂自己的：这样我们因此受到告诫，当以何等大的慷慨，按照祂的诫命，将那些不是我们创造、而是从祂的恩赐中领受的东西，给予我们的仇敌。\par

80. 但是，谁能准备忍受来自弱者的伤害，只要对他们的救恩有益；宁可忍受他人更多的不义，也不报复自己所受的；凡向他求什么的，就给他，或给他所求的，如果他有，而且能够正当地给予，或给予善意的建议，或显示仁慈的态度，不拒绝愿意借东西的人；爱他的仇敌，善待恨他的人，为迫害他的人祈祷——我说，谁做这些事，除非是那完全而彻底\OriginalTerm{慈悲的}{merciful}人？而且依靠那说\BibleQuote{我喜欢仁爱，而不是祭献}的祂的帮助，藉着那劝言，能够避免不幸。所以，\BibleQuote{怜悯人的人是有福的，因为他们要得到怜悯。}但现在，我认为更方便的是，在这点上，读者因如此长的卷帙而疲倦，可以稍事休息，在另一卷中考虑剩下的内容。\par

\SourceRecord{Translated by William Findlay.；Nicene and Post-Nicene Fathers, First Series；Vol. 6.；Edited by Philip Schaff.；Buffalo, NY: Christian Literature Publishing Co.,；1888.；<http://www.newadvent.org/fathers/16011.htm>.}

% structure: p0001 source=h1 target=section reason=page-title

\section{论山中圣训，第二卷}

\Emph{玛 6-7}\par

% structure: p0003 source=h2 target=subsection reason=relative-heading-rank; base=h2

\subsection{第一章}

1. 第一书以仁慈为题结束，接着便是洁净心灵，本书即以此开始。洁净心灵，犹如洁净那只用以看见天主的眼睛；保持这只眼睛的纯一，其审慎之程度应与所见之对象的尊贵相当。然而，即使这只眼睛在很大程度上得到了洁净，仍难以防止某些污秽不知不觉地渗入，这些污秽往往伴随着我们善行而来——例如，人的赞美。如果，不正直地生活是有害的；那么，正直地生活，却不愿受人赞美，这岂非与人的事业为敌？人的事业肯定是越痛苦，人的正直生活就越不能给人带来快乐？因此，如果你生活在他们中间的那些人，在你正直地生活时不赞美你，他们是在错误中；但如果他们赞美你，你就处于危险中，除非你有一颗如此纯一纯洁的心，以至于在你正直行事时，你并非因人的赞美而如此行；你反而更祝贺那些赞美正义之事的人，因为他们喜欢善事，而不是祝贺你自己；因为即使无人赞美你，你也会正直地生活；并且你明白，这赞美对赞美你的人有益，仅当他们所尊荣的不是你在善行中的自己，而是天主，因为每个生活良善的人都是祂最神圣的殿宇；这样，达味所说的话便应验了：\BibleQuote{我的灵魂要在上主内受赞美；谦卑的人听见，也必喜乐。}因此，纯洁的眼睛在行事正直时不看人的赞美，也不应参考这些赞美，\Emph{即}不以取悦人为目的做任何正确的事。如果你只着眼于人的赞美，你也会倾向于假冒为善；因为人看不见内心，也可能赞美虚假的事。那些这样做的人，\Emph{即}假冒良善的人，是二心之人。因此，没有人有纯一的心，\Emph{即}纯洁的心，除非他超越人的赞美；当他生活良善时，只仰望天主，并努力取悦那唯一监察良心者。而且，从良心的纯洁所发出的一切，越不渴望人的赞美，就越值得赞美。\par

2. 祂说：\BibleQuote{所以，你们应当小心，不要在众人面前行你们的义，故意叫他们看见；}\Emph{即}，小心不要怀着这种意图行义，也不要把你们的幸福建立在被人看见上。\BibleQuote{不然，你们在天父那里就没有赏报了。}不是若被人看见；而是若怀着被看见的意图行义。因为，[如果是前者]，那么在这篇讲道开始时所说的话又作何解释：\BibleQuote{你们是世界的光。建在山上的城是不能隐藏的。人点灯，不放在斗底下，而是放在灯台上，就照亮全家的人。你们的光也当这样照在人前，叫他们看见你们的好行为}？但祂并没有将此作为目的；因为祂接着说了：\BibleQuote{并将光荣归于你们在天之父。}但在这里，祂责备的是，如果我们善行的目的止于此，\Emph{即}我们行义只为了被人看见；在祂说了\BibleQuote{你们要小心，不要在人前行你们的义}之后，便没有加什么。由此可见，祂说这话，并不是要阻止我们在人前行义，而是恐怕我们偶尔在人前行义是为了被人看见，\Emph{即}把目光放在这上面，并以此为我们的目的。\par

3. 宗徒也说：\BibleQuote{如果我仍取悦于人，我就不是基督的仆人了；}而他在另一处又说：\BibleQuote{要事事叫众人喜欢，如同我一样事事叫众人喜欢。}不明白的人以为这是矛盾；而解释是：他说他不取悦于人，因为他习惯于行事正义，并非以取悦人为明确目的，而是以取悦天主为目的，他愿意藉着取悦人这件事，使人心转向对天主的爱。因此，他说他不取悦于人是对的，因为在那件事上他旨在取悦天主；他权威性地教导我们应当取悦人也是对的，但这并不是为了将这寻求为善行的赏报；而是因为一个人若不愿成为他所希望得救之人的榜样，就不能取悦天主；但没有人能效法一个不讨他喜欢的人。因此，正如一个人若说：在这寻船的工作中，我不是寻船，而是寻我的祖国，那并非荒谬之言；同样，宗徒也可以恰当说：在这取悦人的工作中，我不是取悦人，而是取悦天主；因为我并不以取悦人为目的，我的目标是使那些我希望得救的人能效法我。正如他论及为圣徒所作的奉献时所说：\BibleQuote{我不是求礼物，而是求果实；}\Emph{即}，我寻求你们的礼物时，我寻求的不是它，而是你们的果实。因为通过这个证明，可以看出他们朝向天主进步了多少，当他们甘心献上那被寻求之物时，不是因为献祭本身所带来的快乐，而是为了爱德的共融。\par

4. 虽然祂接着又说：\BibleQuote{不然，你们在天父那里就没有赏报了。}祂所指出的无非是：我们应当警惕，不可将人的赞美作为我们行为的赏报，\Emph{即}不可认为我们可以因此获得真福。\par

% structure: p0008 source=h2 target=subsection reason=relative-heading-rank; base=h2

\subsection{第二章}

5. \BibleQuote{因此，当你施舍时，}祂说，\BibleQuote{不可在你前面吹号，像假善人在会堂及街道上所行的，为要获得人的光荣。}祂说：不可像假善人那样渴望被人认识。显然，假善人心里并没有他们所呈现在人眼前的东西，因为假善人是伪装者，犹如在剧场中扮演其他角色的人。例如，在悲剧中扮演\Person{阿伽门农}或任何其他历史或传说中的人物，他其实不是那个人本身，而是扮演他，被称为假善人。同样，在教会或人类生活的任何层面，凡希望显得与真实不符的人都是假善人。因为他假装为义人，却没有显露出真正的自己；他将全部行为的果实放在人的赞美上，连假善人也能得到这种赞美，因为他们欺骗了那些以为他们是好人并称赞他们的人。然而，这样的人从洞察人心的天主那里得不到赏报，除非是他们欺骗的惩罚；但从人那里，祂说，\BibleQuote{他们已经获得了他们的赏报。}最公义的话将对他们说：\BibleQuote{离开我，你们这些行骗的人；你们称呼我的名，却不实行我的作为。}因此，那些施舍只是为了获得人的光荣的人，已经得到了他们的赏报；这不是说他们获得了人的赞美，而是说他们故意为了获得这种赞美而行善，这正如上文所讨论的。因为行善者不应寻求人的赞美，而应让赞美伴随行善者，以便那些能效法他们所赞美之事的人能够获益，而不是让被赞美的人以为他们对他有什么益处。\par

6. \BibleQuote{但你施舍时，不要让左手知道右手所行的。}如果你将不信者理解为左手，那么取悦信徒似乎就无可厚非；然而我们完全被禁止将我们善行的果实和目的放在任何人的赞美上。但关于这一点，那些因你的善行而喜悦并效法你的人，我们不仅应在信徒眼前，也应在不信者眼前行动，好因我们所行的善工——这些善工应受赞美——他们能光荣天主，并获得救恩。但如果你认为左手是指敌人，那么当你施舍时，你的敌人不应知道；那么为何主自己在祂的敌人犹太人围绕时，仍仁慈地医治人呢？为何宗徒\Person{伯多禄}在丽门怜悯并医治瘸子时，也招来了敌人对他自己和其他基督门徒的愤怒呢？再者，如果必须不让敌人知道我们施舍，那么我们又该如何对待敌人本身，以实践那诫命：\BibleQuote{如果你的仇人饿了，给他吃；如果他渴了，给他喝。}？\par

7. 第三种意见是属血肉的人通常所持的，荒谬可笑，若非我发现不少人陷入此错误，我就不提了；他们说\Quote{左手}指的是妻子。因为在家庭事务中，女人通常更吝啬钱财，所以丈夫出于怜悯而花钱在穷人身上时，应瞒着妻子，以免发生家庭争吵。仿佛只有男人才是基督徒，而这条训诲不是对女人说的！那么，女人被命令向哪只左手隐瞒她的慈悲善行呢？丈夫也是他妻子的左手吗？这种说法极其荒谬。或者，若有人以为他们互为左手；那么若配偶一方擅自花费家庭财产而违背另一方意愿，这样的婚姻就不是基督徒的婚姻。但无论哪一方，若愿意按照天主的命令施舍，发现有人反对，那人就必然是反对天主命令的敌人，因此应被视为不信者。关于这类人的命令是：信主的丈夫应借善行和品行赢得妻子，信主的妻子应借善行和品行赢得丈夫。因此，他们不应向对方隐瞒善行，因为善行能相互吸引，使一方能吸引另一方加入基督徒的信仰共融。也不应为了取悦天主而行窃。然而，若因对方的软弱尚且不能平静地忍受那些虽非不义或非法的事，而需要隐瞒一些事，那么在此处\Quote{左手}并非此意，这从整个段落中可以很容易看出，并同时会发现祂所称的左手是什么。\par

8. \Quote{你们要谨慎} 祂说，\Quote{不可在人前显出你们的义德，为叫他们看见；否则，你们在你们天上的圣父那里就没有赏报了。} 这里祂一般性地提到了义德，然后详细地加以说明。因为以施舍的方式所做的事是义德的一部分，为此祂将二者联系起来说：\Quote{所以，当你施舍时，不要在你前面吹号，像假善人在会堂及街市上所行的一样，为获得人的光荣。} 这指向祂先前所说的：\Quote{你们要谨慎，不可在人前显出你们的义德，为叫他们看见。} 但随后的话：\Quote{我实在告诉你们，他们已获得了他们的赏报。} 则指向祂上面所说的另一句话：\Quote{否则，你们在你们天上的圣父那里就没有赏报了。} 然后是：\Quote{但你施舍时。} 祂说：\Quote{至于你们，} 除了指不要像他们那样之外，还能指什么呢？那么，祂命令我做什么呢？祂说：\Quote{但你施舍时，} \Quote{不要让你的左手知道你右手所行的。} 因此，那些人的行为使得左手知道右手所做的事。所以，在他们身上受责备的，正是你要禁止做的。受责备的正是这一点：他们行事是为了寻求人的称赞。因此，左手似乎没有比喜爱赞美更合适的含义了。而右手则意味着遵守天主诫命的意向。因此，当施舍者的意识中掺杂了追求人称赞的欲望时，左手就知道了右手的工作：\Quote{所以，不要让左手知道你右手所行的；} \Emph{即}，当你努力履行天主的命令而施舍时，不要让你的意识中掺杂追求人称赞的欲望。\par

9. \Quote{好使你的施舍隐而不露。} 除了指一种良好的良心，不能被人眼看见，也不能以言语揭示之外，\Quote{隐而不露} 还能指什么呢？事实上，大多数的人说了许多谎言。因此，如果右手在内心隐密地行动，所有外在的、可见的和暂时的事物，都属于左手。所以，让你的施舍深藏于自己的意识中，许多人凭着善意施舍，即使他们没有金钱或其他任何可以施舍给需要的人的东西。但许多人在外表上施舍，内心却不施舍，他们或是出于野心，或是为了某些暂时的目标，希望显得仁慈，这些人只能算作左手在工作。另一些人则似乎处在两者之间的中间位置：他们以指向天主的意向施舍，然而在这美好的愿望中也掺杂了某种对赞美的渴望，或对某种易腐朽的暂时事物的贪求。但我们的主更严厉地禁止只有左手在我们内工作，祂甚至禁止左手与右手的工作混合：这就是说，我们不仅要避免单单因暂时的欲望而施舍；而且在这项工作中，我们甚至不能这样仰望天主，以致将对外在利益的追求与其混合或结合。因为这里讨论的是心灵的净化，心若不纯洁，就不会洁净。然而，如果它侍奉两个主人，既不完全因追求永恒而洁净其目光，又因对可朽坏且暂时之物的爱而蒙蔽它，那它如何能纯洁呢？所以，\Quote{你的施舍} 因此，\Quote{要隐而不露；你的圣父，在暗中看见的，必要报答你。} 这完全是最公义、最真实的。因为如果你期待从那位独一的洞察良心者那里获得赏报，就让良心本身足够使你配得赏报。许多拉丁文抄本如此记录：\Quote{你的圣父，在暗中看见的，必要公开报答你。} 但因为我们在更早的希腊文抄本中没有找到\Quote{公开} 一词，所以我们认为对此无需多言。\par

% structure: p0014 source=h2 target=subsection reason=relative-heading-rank; base=h2

\subsection{第三章}

10. \Quote{当你们祈祷时，} 祂说，\Quote{不要像假善人一样；因为他们喜欢站在会堂里和街头祈祷，为叫人看见。} 同样在这里，错误不是被人看见，而是为了被人看见而做这些事；反复说同样的话是多余的，因为只要遵守一条规则即可，从这条规则我们学到，我们应当恐惧和避免的不是人们知道这些事，而是怀着这样的意图去做，即在这些事中追求取悦人的效果。我们的主自己也同样保留了相同的话，当祂同样加上：\Quote{我实在告诉你们，他们已获得了他们的赏报。} 藉此表明祂禁止这种行为——追求那种愚昧人因受人称赞而喜乐的赏报。\par

11. 祂说：\BibleQuote{但你们祈祷时，要进入你们的卧室。}这些卧室不就是我们的心吗？正如\BibleBook{}{圣咏}上说的：\BibleQuote{你们在你们心中所言的，在你们的床上须要懊悔。}祂又说：\BibleQuote{关上门后，要祈祷你们在暗中的父。}进入我们的卧室还是小事，如果门向无礼者敞开，外面的事物就会粗鲁地冲进来袭击我们内在的人。我们已经说过，外面的一切都是暂时的、可见的事物，它们通过门——即肉体的感官——进入我们的思想，并以一群虚幻的幻象喧闹地打扰祈祷者。因此，必须关门，即抵抗肉体的感官，以便向父献上属神的祈祷，这祈祷是在内心深处向在暗中的父献上的。祂说：\BibleQuote{你们的父在暗中看见，必要报答你们。}而这句话必须以这样的结语作为总结；因为在此阶段，劝告不是我们应当祈祷，而是我们应当如何祈祷。前面的劝告也不是我们应该施舍，而是我们应怀着怎样的精神施舍，因为祂正在教导我们如何净化心灵，而唯有以纯洁的智慧之爱，一心一意追求永生，才能净化心灵。\par

12. 祂说：\BibleQuote{但你们祈祷时，不要唠叨，像外邦人一样；因为他们以为话多了必蒙垂听。}如同伪善者的特点是祈祷时向人炫耀，他们的果实是讨人喜悦，同样，外邦人的特点是认为话多了必蒙垂听。事实上，每一种唠叨都来自外邦人，他们努力操练舌头而不是净化心灵。他们试图将这种无用的努力甚至转移到通过祈祷影响天主的事上，以为法官就像人一样，可以被言语说服。\BibleQuote{因此，不要像他们，}唯一的真正导师说。\BibleQuote{因为你们的父在你们祈求祂以前，已经知道你们需要什么。}因为如果使用许多言语是为了让无知的人得到教导和教诲，那么对于知道万物的祂来说，有什么必要呢？因为一切存在的事物，因着它们的存在本身，就向祂说话，并显示它们是被造就的；而那些未来的事物，在祂的知识和智慧中并不隐藏，在祂内，过去和将来的事都是现在的，且不能过去。\par

13. 然而，既然即将由祂亲自说出的话语——即使很少——用来教导我们如何祈祷，那么有人可能会问：对于那位在一切事发生之前就知道一切、并且如前所述在我们祈求之前就知道我们需要什么的祂，为什么还需要这些有限的话语呢？首先，答案是：我们应当向天主陈情，以获得我们所愿的，不是靠言语，而是靠我们心中怀有的思想，以及我们思想的指向，怀着纯洁的爱和诚恳的愿望；但是，我们的主以言语教导了我们这些思想本身，以便我们通过记忆在祈祷时能想起这些思想。\par

14. 但是，我们还可以问（无论我们用思想还是用言语祈祷），如果天主已经知道我们需要什么，为什么还需要祈祷本身呢？除非是祈祷本身的努力能使我们的心灵平静和净化，并使之更有容量去接受那些属神地倾注在我们内的恩赐。因为并不是我们祈祷的迫切使天主垂听我们，祂随时准备赐予我们祂的光，不是物质的光，而是理智和属神的光：但我们并不总是准备好接受，因为我们倾向于其他事物，并因对暂时之物的贪求而陷入黑暗。因此，在祈祷中，我们的心转向祂，祂永远准备给予，只要我们愿意接受祂所赐的；在转向的行为中，内在的眼睛得到净化，因为所贪求的暂时之物被排除，使纯洁的心能承受那神圣照耀的、没有日落或变化的纯光：不仅承受，而且得以停留在其中；不仅没有烦恼，而且带着不可言喻的喜乐，在其中真实而真正幸福的生命得以成全。\par

% structure: p0020 source=h2 target=subsection reason=relative-heading-rank; base=h2

\subsection{第四章}

15. 但现在我们须考虑，我们藉以学习祈祷何事、并获得所求的祂，教导我们祈求什么。祂说：\BibleQuote{你们应当这样祈祷：} \BibleQuote{我们的天父！愿你的名被尊为圣，愿你的国来临，愿你的旨意承行在人间，如同在天上。求你今天赏给我们日用的食粮；求你宽免我们的罪债，如同我们也宽免得罪我们的人；不要让我们陷入诱惑，但救我们免于凶恶。} 既然在一切祈祷中，我们须先求得我们所祈求者的善意，然后说出我们所求；善意通常透过向祈祷对象献上赞美而获得，这通常放在祈祷的开端；而在此特别之处，我们的主命我们只说\BibleQuote{我们的天父}。因为赞美天主的话很多，遍及全部圣经各处，每个人阅读时自能思考；然而，在\Place{以色列}子民中，从未有诫命叫他们说\BibleQuote{我们的父}，或叫他们以父亲之名向天主祈祷；而是以主之名向他们启示，因为他们仍是仆人，即仍\Emph{i.e.}按肉身生活。我这样说，是因为他们接受了律法的诫命，并奉命遵守；因为先知们常指出，这位我们的主原本也能作他们的父，只要他们没有偏离祂的诫命：例如，我们有这句话：\BibleQuote{我养育了子女，使他们成长，但是他们背叛了我}；以及另一句：\BibleQuote{我曾说过：你们是神，都是至高者的子女}；还有这一句：\BibleQuote{我既是父亲，我的尊敬在哪里？我既是主人，我的敬畏在哪里？} 以及许多其他语句，其中犹太人被指责因其罪过表明他们不愿成为儿子；至于那些预言未来基督徒子民的话，说他们将拥有天主为父，则暂不提及，正如那福音的话：\BibleQuote{他赐给他们成为天主子女的权能}。宗徒\Person{保禄}又说：\BibleQuote{承继人虽然是全部产业的主人，但在未成年的时候，与奴隶没有分别}；并提到我们领受了做义子的圣神，藉此我们呼喊：\BibleQuote{藉着这圣神我们呼喊说：\InnerQuote{阿爸，父呀！}}\par

16. 既然我们蒙召得永生的产业，与基督同为承继人，并获致义子的名分，并非由于我们的功德，而是由于天主的恩宠；我们在祈祷的开端就提到这同一恩宠，当我们说\BibleQuote{我们的父}时。藉此称谓，既激发爱德——因为对儿子来说，有什么比父亲更亲爱呢？——又激发恳求的姿态，当人对天主说\BibleQuote{我们的父}时；并对获得所求之事有某种确信；因为在祈求任何事物之前，我们已领受了如此大的恩赐，即被允许称天主为\BibleQuote{我们的父}。既然祂已赐予这根本之事，即他们成为儿子，如今还有什么不赐予祈求的儿子呢？最后，何等大的忧虑占据人心，使那说\BibleQuote{我们的父}的人，不致不配如此伟大的父亲！因为如果一个平民被那位元老本人允许称一位年长的元老为父；无疑他会战栗，不敢轻易去做，思量自己出身的卑微、资源的匮乏、以及平民身份的卑贱；更何况我们称天主为父时，岂不更当战栗，如果我们品格中有如此大的污点和卑劣，以至于天主更合理地驱逐这些远离祂自己，胜于那位元老驱逐任何乞丐的贫穷！因为，的确，他（元老）藐视乞丐身上那些连他自己也可能因人事变迁而陷入的东西；但天主从不陷于品格的卑劣。感谢祂的仁慈，祂要求我们这样作，即祂应作我们的父——这种关系不是靠我们的任何付出，而唯独靠天主的善意。在此也有对富人和贵族的告诫，就这个世界而言，当他们成为基督徒后，不应骄傲地对待穷人和出身低微的人；因为他们与他们一同称天主为\BibleQuote{我们的父}——除非他们承认自己互为弟兄，否则不能真诚而虔诚地使用这个称谓。\par

% structure: p0023 source=h2 target=subsection reason=relative-heading-rank; base=h2

\subsection{第五章}

17. 因此，那蒙召得永生产业的新子民，当使用新约的话，说\BibleQuote{我们的天父}，即\Emph{i.e.}在圣洁和正义者中。因为天主不受空间限制。因为天固然是世间更高的物质天体，但仍属物质，因此只能存在于某一特定地方；但若认为天主的地方在天上，即世界的较高部分，那么飞鸟比我们更有价值，因为它们的生命更接近天主。但经上并未记载：主亲近身材高大的人，或亲近住在山上的人；而是记载：\BibleQuote{主亲近心灵破碎的人}，这更指向谦卑。但正如罪人被称为土，经上对他说：\BibleQuote{你本是土，还要归于土}；同样，义人可被称为天。因为经上对义人说：\BibleQuote{天主的宫殿是圣的，你们就是这宫殿}。因此，如果天主居住在他的圣殿，而圣人们就是他的圣殿，那么\BibleQuote{你在天上}的表达正当地被用来指：你在圣人们中。这样的类比最为合适，以致在灵性上，义人与罪人之间的区别，如同物质上天与地之间的区别那样大。\par

18. 为了表明这一点，当我们站立祈祷时，我们转向东方，那是太阳升起的地方：这并不是因为天主也住在那里，好像那无所不在、不是占据空间而是以祂威能权能临在的主，遗弃了世界的其他部分；而是为了提醒心灵转向一个更卓越的本性，\Emph{即}转向天主，当它那属地的身体转向一个更卓越的身体，\Emph{即}属天的身体时。这也适合宗教的不同阶段，并且极为有益，在所有大小人的心中，都应培养对天主的相称观念。因此，对于那些仍被可见之美吸引，无法思考任何无形体之物的人而言，既然他们必然偏爱天堂胜过地上，如果他们相信天主——他们至今仍以有形的方式思考祂——是在天上而非地上，那么他们的看法是更可容忍的：这样，当将来有一天他们学到灵魂的尊贵甚至超过天上的身体时，他们就会在灵魂中而不是在天上的身体中寻求祂；当他们学到罪人与义人的灵魂之间有多么大的距离时，正如他们当初尚属肉性智慧时，不敢将祂放在地上，而是放在天上一样，此后他们将以更好的信德或理智，再次在义人的灵魂中而非在罪人的灵魂中寻求祂。因此，当经上说：\BibleQuote{我们的天父在天上}时，这被正确地理解为是在义人的心中，如同在祂的圣殿中。同时，这样祈祷的人也愿意他所呼求的那一位也居住在他自己内；当他为此努力时，他就实践了正义——这是一种使天主被吸引而居住在灵魂中的服务。\par

19. 现在让我们看看应该祈求什么。因为已经陈述了向谁祈祷，以及祂住在哪里。首先，在祈求的事物中，有这句恳求：\BibleQuote{愿你的名受显扬。}这祈求并不是因为天主的名字本来不圣洁，而是为了使人尊它为圣；\Emph{即}，使天主如此为人们所知，以致他们认为没有什么比它更圣洁，并且更畏惧触犯它。因为经上说：\BibleQuote{在犹太，天主为人所知；在以色列，祂的名号是伟大的，}我们不应这样理解，好像天主在一个地方小，在另一个地方大；而是在那里祂的名号是伟大的，那里的人按照祂威严的伟大来称呼祂。同样，在那里祂的名号被称为圣洁的，那里的人以敬畏和惧怕触犯祂的心来称呼祂。这就是现今正在发生的事情，福音通过在各个民族中为人所知，藉着祂圣子的管理，赞颂唯一天主的名号。\par

% structure: p0027 source=h2 target=subsection reason=relative-heading-rank; base=h2

\subsection{第六章}

20. 接下来是：\BibleQuote{愿你的国来临。}正如主自己在福音中所教导的，审判的日子将在福音传遍万民时发生：这属于天主之名的圣化。因为这里\BibleQuote{愿你的国来临}这个说法也并非暗示天主现在没有统治。但也许有人会说\BibleQuote{来临}的意思是\Emph{在地上}；好像祂现在真的没有在地上统治，也从创世以来没有统治过一样。因此，\BibleQuote{来临}应理解为\Emph{向人类显现}。因为就像光一样，虽然存在，但对瞎子和闭眼的人却是缺席的；同样，天主的国虽然从未离开大地，但对那些不认识它的人却是缺席的。但当祂的独生子从天降临时，将无人能不认识天主的国，祂不仅以理智可领会的方式，而且以可见的方式，以神人的身份来临，审判生者死者。在那审判之后，\Emph{即}当义人与不义之人被区分和分离的过程完成后，天主将如此居住在义人之中，以致没有人需要人的教导，而所有人都将如经上所记载的：\BibleQuote{都是受天主教导的。}那时，幸福的生命的各个方面将在圣人们中永远完美，正如现在最神圣和最有福的天使们因唯独天主是他们的光明而变得智慧和有福一样；因为主也向祂自己的人应许了这事：\BibleQuote{在复活时，}祂说，\BibleQuote{他们将如同天上的天使一样。}\par

21. 因此，在那我们诵念\Quote{愿你的国来临}的祈求之后，接着是\Quote{愿你的旨意奉行在人间，如同在天上}：\Emph{即}，就如同你的旨意在于天上的天使们，使他们完全依附你，充分享有你，没有错误蒙蔽他们的智慧，没有痛苦妨碍他们的福乐；同样，愿它在你在地上的圣人们身上承行，他们虽以身体而论出自尘土，且虽然要进入天上的居所和变换，却仍将从尘土被提升。对此，天使的赞歌也提供了参照：\BibleQuote{天主受享光荣于高天，主爱的人在世享平安}；如此，当我们的善意先行，跟随那呼召者时，天主的旨意就在我们内得以成全，如同在天上的天使们一样；这样就没有任何敌对阻挡我们的福乐：这就是平安。\BibleQuote{愿你的旨意奉行}也可以正确地理解为：让你的诫命得到服从：\BibleQuote{如同在天上，也在地上}，\Emph{即}如同天使们，也如同人们。因为，当天主的诫命被遵守时，天主的旨意便被承行，主自己如此说，当他肯定：\BibleQuote{我的食物就是承行派遣我者的旨意}；且常言：\BibleQuote{我来不是为执行自己的旨意，而是为执行派遣我者的旨意}；又当他说：\BibleQuote{你看我的母亲和我的兄弟！凡承行天主旨意的，就是我的兄弟、姊妹和母亲}。因此，至少在那些承行天主旨意的人中，天主的旨意得以成就；不是因为他们在促使天主愿意，而是因为他们做了他所愿意的，\Emph{即}他们按他的旨意行事。\par

22. 也有另一种诠释：\BibleQuote{愿你的旨意奉行在人间，如同在天上}——如同在圣善和正义的人身上，也在罪人身上。此外，这还可以从两个方面理解：要么是我们应为仇敌祈祷（因为不顾他们的意志，基督徒和天主教的名字仍然传播？），因此说：\BibleQuote{愿你的旨意奉行在人间，如同在天上}——仿佛意思是：正如义人承行你的旨意，同样，也让罪人承行，以便他们能归向你；要么是这一含义：\BibleQuote{愿你的旨意奉行在人间，如同在天上}，以便每个人得到应得的；这将在最后审判时发生，义人得赏报，罪人受惩罚——当绵羊与山羊分开时。\par

23. 另一种诠释也不荒谬，反而完全符合我们的信德和望德，即我们将天地理解为精神与肉体。既然宗徒说：\BibleQuote{以理智我本身服从天主的法律，但以肉体服从罪的法律}，我们看见天主的旨意是在理智中承行的，\Emph{即}，在精神中。但当死亡被胜利吞噬，这必朽的穿上不朽，这将在肉身复活及义人应许的改变时发生，正如同一宗徒的预言，愿天主的旨意承行在地上如同在天上；\Emph{即}，以这样一种方式：正如精神不反抗天主，而是追随并执行他的旨意，同样，身体也不反抗精神或灵魂——目前它被身体的软弱所困扰，并倾向于肉体的习惯：这将是永生中完美平安的一个要素，即不仅意愿与我们同在，而且行善的实行也与我们同在。他说：\BibleQuote{因为愿意行善的意愿已在我内，但如何去行善我却找不到}：因为在地上还未如同天上，\Emph{即}在肉体中还未如同在精神中，天主的旨意得以承行。因为即使在我们的苦难中，天主的旨意也得以承行，当我们通过肉体承受那些因我们的必朽而应得的痛苦时，这是我们的本性因罪而应得的。但我们应为此祈祷，愿天主的旨意承行在地上如同在天上；以便正如我们怀着内心按内在之人喜爱法律一样，同样，当我们的身体发生改变后，我们没有任何部分会因地上的忧伤或快乐而反对这种喜悦。\par

24. 那种观点也不违背真理，即我们要理解这句话，\BibleQuote{愿你的旨意奉行在人间，如同在天上}，如同在我们的主耶稣基督自己身上，也如同在教会中：仿佛有人说，如同在履行圣父旨意的男人身上，也如同在许配给他的女人身上。因为天和地适合被理解为男人和妻子；因为地因天的滋润而结果实。\par

% structure: p0033 source=h2 target=subsection reason=relative-heading-rank; base=h2

\subsection{第七章}

25. 第四项恳求是：\BibleQuote{我们今日的日用粮，求祢今天赐给我们。}日用粮可以指维持现世生活所需的一切，主曾教导：\BibleQuote{不要为明天忧虑；}因此加上了\Quote{求祢今天赐给我们。}或者，它可指我们每日领受的基督圣体圣事；又或是指灵性神粮，主曾就此说：\BibleQuote{你们不要为那可朽坏的食粮劳碌，}又说：\BibleQuote{我是从天上降下的生活的食粮。}这三种理解中哪一种更合理，需加考虑。因为或许有人会奇怪：既然主亲自说：\BibleQuote{不要为吃什么、穿什么忧虑，}为何我们还要祈求获得此生所需之物——例如食物和衣裳？难道一个人能不为他所祈求获得之物忧虑吗？祈祷本应怀着如此恳切之心，乃至一切关于关上门户的教训以及\BibleQuote{你们先该寻求天主的国和它的义德，这一切自会加给你们}的命令都与此相关？诚然，他并未说：先寻求天主的国，然后再寻求那些东西；而是说：\Quote{这一切自会加给你们，}即即使你们不寻求，也会赐下。但我不知道能否找出一个合理的解释，说明一个人怎能不寻求他恳切祈求天主赐予之物。\par

26. 关于主的圣体圣事——为使那些多在东方地区、不每日领受主的晚餐的人不致提出异议，他们称这饼为日用粮——为使他们缄默，不以教会权威本身来辩护自己的看法，因为他们如此行却不致招致丑闻，且管辖他们教会的负责人并未禁止，他们不顺从也未被定罪；由此证明在这些地区并不将此理解为日用粮。因为若果真如此，那些不每日领受的人就应被指控犯了大罪。但如所述，且不从这类人的情形作任何论证——至少以下思考应见于反省者：我们从主领受了祈祷的规则，不可增添或删减任何内容。既然如此，谁敢说我们只应使用一次天主经？或者，即便我们在领受主圣体前使用过两次或三次，之后在当日其余时辰就确定不再如此祈祷？因为我们将不能再就所已领受的饼说：\BibleQuote{求祢今天赐给我们；}或者人人都能强迫我们在当日最后一刻举行那圣事。\par

27. 因此，余下的解释是：我们应将日用粮理解为灵性食粮，即天主的诫命，我们应每日默想并努力遵行。因为主正就此教训说：\BibleQuote{你们不要为那可朽坏的食粮劳碌。}这食粮目前称为日用粮，因为今世生命由离去又复返的日子度量。的确，只要灵魂的渴望时而转向高处，时而转向低处——\Emph{i.e.}时而转向灵性事物，时而转向肉体事物——就如那有时得食、有时挨饿的人一样，日用粮便每日必需，使饥饿者得以恢复，跌倒者得以扶起。因此，正如在这生命中，即在巨大转变之前，我们的身体因感到亏损而靠食物恢复，同样，灵魂因对天主的渴求被暂时欲望所耗损，也当靠诫命之粮重获力量。此外，这\Quote{求祢今天赐给我们}是在称为\Quote{今天}的时期内说的，\Emph{i.e.}在今世生命之中。因为此世之后，我们将获得如此丰盛的灵性食粮，直至永恒；那时它不再称为日用粮，因为不存在流逝的时间，即日复一日从而可称为\Quote{今天}者。但正如经上说：\BibleQuote{今天，如果你们听从他的声音，}宗徒在《希伯来书》中解释为：\Quote{只要还有\Quote{今天}之称，}此处\Quote{求祢今天赐给我们}也应作此理解。然而，若有人愿将本句也理解为身体所需的食粮或主的圣体圣事，则应将三者结合理解：即我们当同时祈求日用粮——身体所需之食粮、可见的祝圣之饼，以及天主圣言的不可见之粮。\par

% structure: p0037 source=h2 target=subsection reason=relative-heading-rank; base=h2

\subsection{第八章}

28. 第五个祈求如下：\Quote{免我们的债，如同我们免了人的债。}显然，债是指罪过，这可以从主亲自所说的那句话看出：\Quote{你决不能从那里出来，直到你还清最后一文钱}；也可从祂称那些被报告说被杀的人为欠债者这一事实看出，无论是被塔压死的人，还是黑落德将他们的血与祭品混杂的人。因为祂说，人们以为这是因为他们过于欠债，\Emph{即}过于犯罪，祂接着说：\Quote{我告诉你们：不是的；但你们若不悔改，都要同样丧亡。}因此，这里不是要求人免除金钱上的债务，而是免除别人可能犯下的任何冒犯自己的罪过。因为我们被命令免除金钱上的债务，是由前面已经给予的那条诫命：\Quote{若有人愿意告你，拿你的里衣，连外衣也让他拿去}；而且不必对每一个金钱上的债务人免除债务；而只对不愿还债的人，甚至到愿意告到法庭的程度。因为\Quote{天主的仆人}，如宗徒所说，\Quote{不可争讼。}因此，对于不愿自愿或经请求还清所欠金钱的人，应当予以免除。因为他不愿还债的原因有两个：要么他没有钱，要么他贪婪并觊觎他人的财物；这两者都属于贫穷状态：前者是物质上的贫穷，后者是心态上的贫穷。所以，谁若免除这样人的债务，就是免除穷人的债务，并实践基督徒的工作；同时那条原则仍然有效，即他应该心里准备失去所欠的钱。因为如果他已经尽了各种努力，安静而温和地要求归还，与其说是为了金钱利益，不如说是为了使那人改正（毫无疑问，有能力支付却拒不支付对那人是有害的），那么他不仅不会犯罪，反而做了一件大善事，试图阻止那人（他想要通过他人的金钱获利）在信仰上遭受船难；后者是如此严重，根本无法相比。由此可见，在这第五个祈求中，当我们说\Quote{免我们的债}时，这些话并非指金钱，而是指所有别人冒犯我们的方式，因而也间接指金钱。因为一个有能力支付却拒绝支付所欠钱的人，是冒犯了你。如果你不免除这个罪，你就不能说\Quote{免我们的债，如同我们免了人的债}；但如果你宽恕了它，你就看到那被命令献上这样祈祷的人，也被提醒要免除金钱上的债务。\par

29. 这也可以这样解释：当我们说\Quote{免我们的债，如同我们免了人的债}时，只有当我们不宽恕那些请求宽恕的人时，我们才被证明违背了这条规则，因为我们也希望在我们恳求祂宽恕时，得到我们最慈爱的圣父的宽恕。但是，另一方面，那命令我们为仇敌祈祷的诫命，并非命令我们为那些请求宽恕的人祈祷。因为已经处于这种心态的人不再是仇敌。然而，一个人若没有宽恕某人，就不可能真实地说他为他祈祷。因此，我们必须承认，凡冒犯我们的罪过都应当被宽恕，如果我们希望我们的圣父宽恕我们冒犯祂的罪过。因为关于报复的主题，如我所想，已经充分讨论过了。\par

% structure: p0040 source=h2 target=subsection reason=relative-heading-rank; base=h2

\subsection{第九章}

30. 第六个祈求是：\Quote{不要让我们陷入诱惑。}有些抄本用\Quote{带}字，我认为这意思相同，因为两种译法都源自同一个希腊词。但许多人在祈祷中如此表达：\Quote{不要让我们被带入诱惑}；也就是说，解释\Quote{带}一词的意义。因为天主自己并不带入，而是任凭那人被带入诱惑，祂按照一个至为深奥的安排和那人的功过，撤回了祂的助佑。而且，常常出于明显的理由，祂判定那人应受这样的抛弃，并被允许带入诱惑。但被带入诱惑是一回事，受试探是另一回事。因为没有试探，没有人能受考验，无论是对自己，如经上所载：\Quote{未曾受过试探的人，他知道什么呢？}或是对别人，如宗徒所说：\Quote{你们在我肉身中所受的试探，你们没有轻看}；因为从这一事实，他得知他们是坚定的，因为他们没有被那些按肉身发生在宗徒身上的试炼所动摇而远离爱德。因为甚至在一切试探之前，我们已被天主所知，祂在一切事发生之前已知道一切。\par

31. 因此，当经上说 \BibleQuote{上主你的天主试探（考验）你，要知道你是否爱祂}，\Quote{要知道}一词是用于表示实际情况，即祂为使你们知道。正如我们称一个日子为快乐的日子，因为它使我们快乐；称寒霜为呆滞的寒霜，因为它使我们呆滞；以及无数类似的说法，无论是在日常用语中，还是在学者的论述中，或是在圣经中都可以找到。而那些反对旧约的异端者，不明白这一点，以为那人——经上说 \BibleQuote{上主你的天主试探你}——身上似乎应当被加上无知的烙印；仿佛福音中没有记载有关主的话说：\BibleQuote{祂说这话，是要试探（考验）他，因为祂自己知道要做什么}。因为如果祂了解祂所试探的那人的心，那么祂试探他究竟想看见什么呢？但实际上，这样做是为了使那受试探的人认识自己，并因众人被主的饼充满而谴责自己的绝望，因为他曾以为他们没有足够的食物。\par

32. 因此，这里的祈祷不是求不受试探，而是求不陷入诱惑：正如如果有人必须经过火的考验，他应祈祷的不是不被火触碰，而是不被烧毁。因为 \BibleQuote{炉火试探陶匠的器皿，忧患的试探显明义人}。因此，若瑟曾受淫乱的诱惑试探，但他没有陷入诱惑。苏撒纳受了试探，但她没有被引入或陷入诱惑；还有许多其他男女圣者：但最突出的乃是约伯，关于他在上主天主面前那令人赞叹的坚忍，那些旧约的异端敌人，当想以亵渎之口嘲笑圣经时，特别挥舞这个武器，说撒殚请求让他受试探。因为他们向无能理解此类事情的无知者提出这样的问题：撒殚怎能与天主说话？他们不明白（因为他们不能，因被迷信和争论盲目）天主并不以祂形体性的体积占据空间，因而存在于一处而不在另一处，或至少有一部分在这里，另一部分在那里；而是祂在祂的尊威中无所不在，不分部分，而是处处完整。但如果他们按肉体理解这话：\BibleQuote{上天是我的宝座，下地是我的脚凳}——我们的主也引用这段经文，当祂说：\BibleQuote{你们总不可发誓：不可指天，因为天是天主的宝座；不可指地，因为地是他的脚凳}——那么魔鬼既在地上，站在天主的脚前，并在祂面前说什么，这有什么可惊奇的？因为他们何时才能明白，没有一个灵魂（无论多么邪恶），只要还能以某种方式推理，在天主的良心里不说话？因为除了天主，谁把自然法律写在人心中呢？——就是那法律，关于它，宗徒说：\BibleQuote{几时，没有法律的外邦人，顺着本性去行法律上的事，他们虽然没有法律，但自己对自己就是法律。如此，他们显示了法律的功用刻在他们的心中，他们的良心也为之作证，因为他们的思想互相矛盾，有时控告，有时辩护，在那天主审判人隐微事的日子。}因此，如同对于每一个理性灵魂，尽管被情欲盲目，但仍在思考和推理，我们将其推理中任何真实的部分，不归功于它本身，而归功于真理之光本身——这光即使微弱，也按其能力照亮它，以便它通过推理能感知一些真理——那么，如果堕落魔鬼的败坏精神，虽被贪欲扭曲，却被描述为从天主本身的声音，即从真理本身的声音，听到了关于它打算试探的义人所抱有的任何真实想法，这有什么可奇怪的呢？但凡虚假的，都应归于那贪欲，魔鬼由此得名。虽然也有这样的情况：天主常藉有形的受造物向善人恶人说话，作为万有的主和治理者，并按各行为的功过安排一切：例如藉天使（他们也以人形显现）和藉先知说：\Quote{上主这样说}。那么，如果天主不是仅仅在思想中，而是藉某种适合此工作的受造物，被称为与魔鬼说了话，这有什么可惊奇的？\par

33. 而且，不要让他们以为天主与他说话有损天主的尊威和公义：因为祂与一个天使性的精神说话（虽然它是愚昧和贪欲的），正如祂与一个愚昧和贪欲的人的精神一样。或者让这些人自己告诉我们，祂如何与那个富有的（贪得无厌的人）说话，那人最愚昧的贪得无厌，祂想要谴责，说：\BibleQuote{糊涂人哪！今夜就要索回你的灵魂，你所备置的，将归谁呢？}主自己在福音中确实这样说，那些异端者，无论愿意与否，都要向福音低头。但如果他们对撒殚向天主请求试探义人这件事感到困惑——我不解释这事如何发生，但我迫使他们解释：为什么主自己在福音中对门徒说：\BibleQuote{看，撒殚求得了你们，好筛你们像筛麦子一样；}又对伯多禄说：\BibleQuote{但是，我已为你祈求，为叫你的信德不致丧失。}当他们向我解释这事时，他们同时也就解释了他们所问我的问题。但如果他们不能解释这事，就不要鲁莽地胆敢在任何书中指责他们在福音中读到而无反感的事。\par

34. 诱惑之所以藉撒殚发生，并非因他的能力，而是因主的许可，或是为惩罚人的罪，或是为按主的慈悲来考验和操练他们。各人所陷诱惑的性质有极大差异。因为出卖主的犹大，他所陷入的诱惑，与伯多禄因恐惧而否认主时所陷入的，性质不同。也有属于人之常情的诱惑，我认为，就是当每个心怀善意的人，因人性软弱，在某项计划中犯错，或在热心劝勉弟兄归正时，比基督徒应有的平和稍有过激而恼怒弟兄。关于这些诱惑，宗徒说：\BibleQuote{你们所受的诱惑，无非是人所常受的；}同时他又说：\BibleQuote{天主是信实的，必不叫你们受诱惑超过你们所能受的；在受诱惑的时候，必给你们开一条出路，叫你们能忍受得住。}在这句话中，他清楚地表明，我们不应祈求不受诱惑，而应祈求不陷于诱惑。因为如果我们遇见了无法承受的诱惑，那就是陷于诱惑了。但危险的诱惑——陷入并被引入其中对我们是有害的——或来自顺境或逆境的世事，只有被顺境的快乐所俘虏的人，才会被逆境的折磨所击倒。\par

35. 第七项即最后一项祈求是：\BibleQuote{但救我们免于凶恶。}因为我们不仅要祈求不要陷于那已经摆脱的凶恶（这是第六项所求的），也要祈求从已陷入的凶恶中被解救出来。当这事完成后，就再无可怕之事，也无需惧怕任何诱惑了。然而在此生，当我们仍带着因蛇的诱惑而陷入的现世可朽之身时，不能指望这种情况发生；但我们应盼望在将来某个时候它会发生：这就是那看不见的望德，宗徒论及它时说：\BibleQuote{但看得见的望德不是望德。}不过，在此生中，天主忠信的仆人们也不应绝望于所赐的智慧。这智慧就是：我们应以最谨慎的警醒，躲避那些因主的启示而明白应当躲避的事；并以最炽热的爱慕，追求那些因主的启示而明白应当追求的事。因为这样，在死亡中卸下这现世生命的残余负担之后，在此生已开始、并时常竭尽全力去把握和确保的福乐，将在适当的时候在人的每一部分得以圆满。\par

% structure: p0047 source=h2 target=subsection reason=relative-heading-rank; base=h2

\subsection{第十章}

36. 这七项祈求之间的区别应当留意并加以称赞。因为我们现在正度过现世的生命，而永恒的生命是所盼望的；并且永恒的事物在尊贵上超越现世，尽管我们只有结束现世之后才能进入永恒。虽然前三项祈求在此生——即度过现世世界——已开始应验（因为天主之名的圣化随着谦卑之主的来临开始实现；祂的国度的来临——祂将光荣地来临——将在世界终末而非世界之后显现；祂的旨意在地上如同在天上的圆满实现——无论你以天地理解为义人和罪人，还是灵魂和肉身，或是主和教会，或是这一切的总和——都将在我们福乐圆满时完成，因此是在世界终了），然而这三项祈求将存留到永远。因为天主之名的圣化将永远继续，祂的国度无终，我们圆满的福乐被应许得永生。因此，这三项事物将在那应许给我们的生命中完全成就并彻底完成。\par

37. 但我们所求的其余四项，在我看来属于现世的生命。其中第一项是：\BibleQuote{求祢今天赐给我们日用的食粮。}因为无论是这被称为日用食粮的东西指精神食粮，还是指圣事中可见的食粮，或指我们这维持生计的食粮，它都属于现时，祂称之为\BibleQuote{今天}，并非因为精神食粮不是永恒的，而是因为圣经中称为日用食粮的东西，是通过言语的声响或任何时间的标记呈现给灵魂的；这一切当众人都受天主教导时，当然不再存在，那时他们不再用身体的动作彼此传授，而是各人凭自己心灵的纯洁汲取那不可言说的真理之光本身。或许也正因为如此，它被称为食粮而非饮料，因为食粮需要掰开和咀嚼才能转化为营养，正如圣经通过被解释和讲解来滋养灵魂；而饮料一旦预备好，便直接融入身体。因此目前，真理是食粮，当它被称为日用食粮时；但那时它将称为饮料，因为不再需要讨论和论述的劳作（如同掰开和咀嚼），只需饮用纯净透明的真理。如今我们的罪被赦免，我们也赦免他人；这是其余四项中的第二项祈求。但那时不再有罪的赦免，因为不再有罪。诱惑骚扰现世的生命；但当这些话完全应验时，诱惑将不复存在：\BibleQuote{祢必将他们隐藏在祢面前的隐密处。}我们盼望脱离的凶恶，以及脱离凶恶本身，当然属于此生，此生命因其必朽性，我们因天主的正义而应得，又因祂的仁慈而获救。\par

% structure: p0050 source=h2 target=subsection reason=relative-heading-rank; base=h2

\subsection{第十一章}

38. 这些祈求的七重数目在我看来也与前面整篇讲道所源起的七重数目相对应。因为如果贫穷的人是有福的，乃是借着对天主的敬畏，因为天国是他们的；那么，愿我们祈求天主的名能藉着那\BibleQuote{洁净而永远长存的敬畏}在人间被尊为圣。如果温良的人是有福的，乃是借着孝爱，因为他们将继承大地；那么，愿我们祈求祂的国来临，无论是临于我们自身，使我们变得温良而不抗拒祂，还是从天而降，在主来临的光荣中，那时我们将喜乐，并将受到赞美，因为祂说：\BibleQuote{来吧，我父所祝福的，继承自创世以来为你们预备的国度。}因为先知说：\BibleQuote{在上主内，我的灵魂将受赞美；温良的人听见了，必会欢喜。}如果哀恸的人是有福的，乃是借着知识，因为他们将受安慰；那么，愿我们祈求祂的旨意承行于地如同在天上，因为当身体（犹如大地）与灵魂（犹如天上）在最终完全的平安中和谐一致时，我们就不再哀恸：因为现今这个时期没有别的哀恸，除非是当这两者彼此争斗，迫使我们说：\BibleQuote{我发觉在我的肢体中另有一条法律，与我的理智法律交战}，并含泪哀叹说：\BibleQuote{我真不幸！谁能救我脱离这死亡之身？}如果饥渴慕义的人是有福的，乃是借着刚毅，因为他们将得饱足；那么，愿我们祈求今日赐给我们日用的食粮，藉此得到支持和滋养，使我们能达到那最丰盛的圆满。如果怜悯人的人是有福的，乃是借着明智，因为他们将蒙受怜悯；那么，愿我们宽恕我们的负债者，并祈求我们的债也被宽恕。如果心地纯洁的人是有福的，乃是借着明悟，因为他们将看见天主；那么，愿我们祈求不要让我们陷于诱惑，免得我们心怀二意，不追求那单一的善，好将我们的一切行为归于它，而同时贪求暂时和世俗的事物。因为那些在人看来艰难和灾祸的诱惑，如果那些藉着人视为美好和喜乐之事的诱惑不能胜过我们，它们对我们就没有力量。如果缔造和平的人是有福的，乃是借着智慧，因为他们将被称为天主的子女；那么，愿我们祈求脱离邪恶，因那自由本身将使我们成为自由的，即天主的子女，使我们能在义子的精神中呼喊：\BibleQuote{阿爸，父啊！}\par

39. 我们也不该轻率地忽略这个事实：在主教导我们祈祷的所有语句中，祂判定那与罪过赦免相关的一句是特别值得推荐的。在其中祂愿意我们慈悲为怀，因为这是逃脱悲惨的唯一智慧。因为在任何其他语句中，我们祈祷时并不像与天主立约那样：我们说：\BibleQuote{宽恕我们，如同我们也宽恕。}如果我们在这盟约中撒谎，整个祈祷就无效了。因为祂这样说：\BibleQuote{因为如果你们宽恕别人的过犯，你们的天父也必宽恕你们；但如果不宽恕别人的过犯，你们的父也必不宽恕你们的过犯。}\par

% structure: p0053 source=h2 target=subsection reason=relative-heading-rank; base=h2

\subsection{第十二章}

40. 接着是关于斋戒的诫命，涉及目前正讨论的那同样的净化心灵。因为在这项工作中，我们也必须警惕，免得某种炫耀和贪求人的称赞潜入进来，使心灵变得双重的，而不让其纯洁专一以便领悟天主。祂说：\BibleQuote{再者，你们禁食的时候，不要像伪君子那样面带愁容；因为他们把脸弄得难看，为叫人看出他们在禁食。我实在告诉你们，他们已经得到了他们的赏报。但你禁食时，要梳头洗脸，不叫人看出你在禁食，但叫你那在隐秘中的父看见；你的父在隐秘中看见，必要报答你。}从这些诫命中可以明显看出，我们的一切努力都要朝向内在的喜乐，免得我们寻求外在的赏报而效法这个世界，并失去那越是内在就越坚固的福乐的应许，天主选择在此应许中使我们成为祂圣子的肖像。\par

41. 但在此处主要应注意，炫耀不仅在于与战利品有关的辉煌与排场，也在于凄惨的肮脏本身；并且由于它假借事奉天主之名而欺骗，因此更为危险。因此，谁因对肉体过分关注以及衣着或其他事物的辉煌而格外显眼，就很容易被这些事物本身证明是世俗浮华的追随者，并不以狡猾的圣洁外表误导任何人；但至于那以基督徒身份自居，却以异常的肮脏和污秽吸引人的目光，且是出于自愿而非迫于必要者，可从其其余行为推断，他这样做是出于对多余身体关注的鄙视，还是出于某种野心：因为主曾命令我们提防披着羊皮的狼；但祂说：\BibleQuote{凭他们的果实，}祂说，\BibleQuote{你们就认出他们来。}因为当通过任何试探，那些他们在这面纱下已获得或渴望获得的东西开始被夺去或被拒绝时，那时就必然显出它是披着羊皮的狼，还是本来的羊。因为基督徒不应因此以多余的装饰取悦人的眼睛，因为伪装者也常常采取朴素而必要的衣着，以欺骗那些不警惕的人：因为如果有时狼披上羊皮，那些羊也不应该脱下自己的皮。\par

42. 因此，通常人们会问，当祂说：\BibleQuote{至于你，当你禁食时，要抹油在你的头上，并洗你的脸，好不叫人看出你在禁食。}祂是什么意思？因为任何人教导（尽管我们可按日常习惯洗脸）我们禁食时也应该在头上抹油，这是不对的。如果所有人都承认这极不恰当，那么我们必须将抹油和洗脸的这条诫命理解为涉及内在的人。因此，抹油指的是喜乐；而洗脸指的是纯洁：所以，那在内心和理智中喜乐的人，就是抹了油的头。因为我们正确理解那在灵魂中居首位者为头，并且由此明显看出人的其余部分受其统治和管理。那不从外面寻求喜乐，从而不以属肉体的方式从人的赞美中获取快乐的人，就是这样做的。因为本应服从的肉体，绝不是整个人性的头。确实，宗徒在教导关于爱妻子的诫命时说：\BibleQuote{从来没有人恨恶自己的肉体}；但男人是女人的头，基督是男人的头。因此，让他就在他的禁食中内心喜乐，因着他的禁食，他这样远离世俗的享乐，从而服从基督，基督愿意根据这条诫命而让头被抹油。因为这样他也会洗他的脸，即洁净他的心，他将以此看见天主，不再因从肮脏所染的软弱而有面纱阻隔；而是坚定稳固，因他纯洁无伪。祂说：\BibleQuote{你们要洗濯，自洁，从我眼前除掉你们的恶行。}因此，我们的脸应当被洗净，脱离那冒犯天主之眼的肮脏。因为我们也用没有帕子遮掩的脸，好像镜子一样观看主的荣光，就被变成同样的形像。\par

43. 常常，对今生必要之物的思念也伤害并玷污我们的内眼；并且它常使心怀两意，以致在那些我们似乎以邻人行事端正的事上，我们并不以主所命令我们的心去行；即，不是因为我们爱他们，而是因为我们想从他们那里获得一些好处以应付今生的需要。但我们应当为了他们永恒的救恩而行善，而非为了自己暂时的利益。因此，愿天主使我们的心倾向于祂的规诫，而非贪得无厌。因为\BibleQuote{命令的目的就是爱德，出自纯洁的心、善良的良心和无伪的信德。}但那因顾及自己今生的需要而照顾他弟兄的人，确实不是出于爱，因为他照顾的不是他应当如己爱的那个人，而是他自己；甚至也不是他自己，因为这样他就使自己的心怀两意，由此被阻碍看见天主，唯有在对天主的注视中才有确定而持久的福乐。\par

% structure: p0058 source=h2 target=subsection reason=relative-heading-rank; base=h2

\subsection{第十三章}

44. 因此，那致力于洁净我们心灵的祂，正当地以一条诫命来跟随祂所说的，祂说：\BibleQuote{不要为自己积蓄财宝在地上，那里有虫蛀和锈蚀，也有贼挖洞偷窃；但要在天上为自己积蓄财宝，那里没有虫蛀和锈蚀，也没有贼挖洞偷窃。因为你的财宝在哪里，你的心也在哪里。}因此，如果心在地上，即一人以专注于获得属世好处的心行事，那在世上打滚的心怎能洁净呢？但如果心在天上，它就会洁净，因为凡是天上的事物都是洁净的。因为任何东西与低劣的本性混合时都会受污染，尽管其本身种类上并不污染；因为黄金即使与纯银混合也会受污染：同样，我们的心灵因对属世之物的欲望而受污染，尽管地本身在其种类和秩序上是纯洁的。但在这个段落中，我们不应将天理解为任何物质的东西，因为一切物质的东西都应算作地。因为那为自己在天上积蓄财宝的人应当轻视整个世界。因此是在那个天中，正如经上所说：诸天的天是属于主的，即在属灵的穹苍中：因为我们不应当将我们的财宝和心安置在将要过去的世上，而应安置在那永远常存的；但天地都要过去。\par

45. 于此，祂显明祂赐予这一切诫命乃为洁净人心，当祂说：\Quote{身体的灯就是眼睛；所以，如果你的眼睛单纯，你全身就充满光明。但若你的眼睛邪恶，你全身就充满黑暗。因此，如果你内在的光（灯）成了黑暗，那黑暗是多么大啊！}（\BibleRef{玛 6:22-23}）这段话当如此理解：由此我们得知，当我们以纯一的心（即怀着属天的意向），以爱德为终向而行事时，我们所有的行为在天主眼中都是纯洁而中悦的；因为爱德也就是法律的满全。因此，我们应在此将\Quote{眼睛}理解为意向本身，我们以它做我们所做的一切；如果这意向纯洁、正直，并注视所当注视的，那么所有因之而行的行为就必然是善的。祂将这些行为称为\Quote{全身}；宗徒也把某些他认为不当的行为称为我们的肢体，并教导要治死它们，说：\Quote{所以要治死你们在地上的肢体：淫乱、污秽、贪欲，}（\BibleRef{哥 3:5}）以及其他诸如此类的事。\par

46. 因此，当考虑的并非人所行何事，而是他行事的意向。因为意向就是我们内在的光，因为对我们自己而言，我们知道自己在以善意向行事；因为一切显明的都是光。而那些从我们发出去到人类社会中的行为本身，其结果是不确定的；因此祂称它们为黑暗。因为当我给一个向我求乞的穷人钱时，我不知道他将用它做什么，或者他将因此遭受什么；可能他用它行了恶，或因它遭受了恶——这是我不希望发生的事，当我给他时，我也不曾抱着那样的意向。因此，如果我怀着善意去行——此事当我行时为我所知，所以称为光——那么我的行为也被光照亮，无论它有什么结果；但那个结果，因其不确定和未知，被称为黑暗。但如果我怀着恶意去行，那么这光本身就成了黑暗。因为被称为光，是因为每个人都知道自己怀着什么意向行事，即使他怀着恶意；但这光本身是黑暗，因为其目标不是单单向高处，而是转向低处，并因双重心而仿佛投下了阴影。\Quote{因此，如果你内在的光成了黑暗，那黑暗是多么大啊！}（\BibleRef{玛 6:23}）即，如果你行事时内心的意向本身（这是你所知的）因贪求属地暂时的事物而被污染、被弄瞎，那么行为本身（其结果不确定）岂不更加被污染、充满黑暗？因为，虽然你以不正不洁的意向所行的事，可能对某人有益，但算在你账上的，是你怎样行的，不是那事对那人有了怎样的结果。\par

% structure: p0062 source=h2 target=subsection reason=relative-heading-rank; base=h2

\subsection{第十四章}

47. 接着，下面的话：\Quote{没有人能服事两个主人，}也应归到同一意向，祂接着解释说：\Quote{因为他要么恨这一个而爱那一个，要么依附这一个而轻视那一个。}（\BibleRef{玛 6:24}）这些话当仔细思考；因为祂立即指出两个主人是谁，说：\Quote{你们不能服事天主又服事财宝。\OriginalTerm{财宝}{mammon}}在希伯来人中，财富被称为\Quote{玛门}。腓尼基人也用此名：因为在腓尼基语中，收益被称为玛门。但谁服事财宝，就是服事那因自己的悖逆而掌管属地之事的主——我们的主称他为今世之王。因此，人\Quote{要么恨}这一个\Quote{而爱那一个}，即天主；\Quote{要么依附这一个而轻视那一个。}因为凡服事财宝的，就是服事一个严酷而败坏的主人：因被自己的私欲缠住，他成了魔鬼的属下，他却并不爱他；因为谁爱魔鬼呢？但他却依附于他；如同在大宅中，一个人因与别人的使女有关而服于严酷的奴役，虽不爱那使女的主人，却爱那使女。\par

48. 但祂说：\Quote{他将轻视那一个，}并非说恨。因为几乎无人能在良心里恨天主；但他轻视，即不敬畏祂，仿佛因祂的良善而自觉安全。圣神藉先知呼唤我们脱离这种粗心大意毁灭性的安全，说：\Quote{我儿，不要罪上加罪，且说：天主的仁慈是大的；}（\BibleRef{德 5:5-6}）又说：\Quote{难道你不知道天主的忍耐是引导你悔改吗？}（\BibleRef{罗 2:4}）因为谁的仁慈能被提到像祂那样大呢？祂宽恕所有回头之人的一切罪过，使野橄榄枝得分享橄榄树的肥汁。谁的严厉又能像祂那样大呢？祂没有怜惜原来的枝条，因不信而折断了它们。但任何愿意爱天主并警醒不得罪祂的人，都不以为自己能服事两个主人；他当从所有双重心中解脱出正直的意向：因为这样他以善心思念上主，以纯朴的心寻求祂。\par

% structure: p0065 source=h2 target=subsection reason=relative-heading-rank; base=h2

\subsection{第十五章}

49. \Quote{因此，}祂说，\BibleQuote{我告诉你们：不要为你们的生命忧虑吃什么，也不要为你们的身体忧虑穿什么。}恐怕，虽然现在所求的并非多余之物，但人心却因必需之物的缘故而变得心怀二意，目标也因此偏斜，去追求属于我们自己的东西，当我们好像在出于同情做一些事的时候；\Emph{i.e.} 当我们想要显得是为某人的益处着想时，却在那件事上寻求自己的利益而非他的益处；而且我们似乎不因此认为自己犯罪，因为我们想要的不是多余之物，而是必需品。但主告诫我们，要记住天主，当祂造我们并使我们由身体和灵魂组成时，赐给了我们比食物和衣服多得多的东西，祂不愿我们因关心这些东西而心怀二意。祂说：\BibleQuote{生命不是比食物更贵重吗？}如此你当明白，那赐予灵魂的主，更会轻易地赐予食物。\BibleQuote{身体不比衣服更贵重吗？}即身体比衣服更贵重：\Emph{i.e.} 同样你当明白，那赐予身体的主，更会轻易地赐予衣服。\par

50. 在此段中，通常有人提出疑问：所论及的食物是否指灵魂，因为灵魂是无形的，而所论的食物是物质的。但让我们承认，在此段中灵魂代表现世的生命，其维持需要物质营养。按此意义，我们也有那句话：\BibleQuote{爱惜自己性命的，必要丧失性命；}这里，除非我们将此表达理解为现世的生命（我们应为天国而丧失它，正如殉道者显然能够做到的那样），否则这条诫命将与另一句话相矛盾：\BibleQuote{人纵然赚得了全世界，却赔上了自己的灵魂，有什么益处呢？}\par

51. 祂说：\BibleQuote{你们看天上的飞鸟：它们不播种，也不收割，也不收进仓里，你们的天父尚且养活它们；你们不比它们更贵重吗？} \Emph{i.e.} 你们更有价值。因为有理性的受造物如人在万物的本性中地位高于无理性的受造物如飞鸟。\BibleQuote{你们中谁能用思虑使自己的身量增加一肘呢？你们为什么为衣裳忧虑呢？}\Emph{i.e.} 那凭其能力和统御使你们身体长到现今身高的天主之眷顾，也能使你们得到衣服；但你们身体的增长并非由你们的关心所致，这可以从以下情形得知：如果你们思虑并希望增加一肘身量，却不能做到。因此，将保护身体的关心交托给祂，你们看见正是因祂的关心，你们才有了如此身量的身体。\par

52. 关于衣服也需要给出一个例子，正如关于食物给出的例子一样。因此祂接着说：\BibleQuote{你们观察田间的百合花怎样生长：它们不劳苦，也不纺织；但我告诉你们：连撒罗满在他极盛的荣华时所穿戴的，也不如这些花中的一朵。田野的草今天还在，明天就投在炉中，天主尚且这样装饰它，何况你们呢？小信德的人啊！}但这些例子不应视为寓意，以致我们追问天上的飞鸟或田间的百合花有什么含义；因为它们在此出现，是为了通过较小的事使我们相信更大的事；就如那个既不敬畏天主也不尊重人的判官，却因寡妇不断烦扰他而允准了她的案件，不是因为虔诚或人道，而是为了免受烦扰。那个不义的判官绝不象征天主的位格；但主愿意我们从这个情形推断出，善良公义的天主多么关心那些恳求祂的人，因为连不义的人也不能轻视那些用不断恳求烦扰他的人，即使仅仅为了避免烦恼。\par

% structure: p0070 source=h2 target=subsection reason=relative-heading-rank; base=h2

\subsection{第十六章}

53. 祂说：\BibleQuote{所以，不要忧虑说：我们吃什么？喝什么？穿什么？——因为这一切都是外邦人所寻求的；你们的天父知道你们需要这一切。你们先该寻求天主的国和祂的义德，这一切自会加给你们。}在此祂最清楚地表明，这些事不应作为我们的福乐去寻求，以致我们为了它们而应在所有行为中行善，但它们仍是必需的。因为福乐（应寻求）与必需品（取为使用）之间的区别，祂用这句话清楚地指出了：\BibleQuote{你们先该寻求天主的国和祂的义德，这一切自会加给你们。}因此，天主的国和义德是我们的善；这是应寻求的，并且应以此为终极目的，我们所作的一切都应为达到这个目的。但因为我们在此生服役，以便能到达那国度，而且我们的生命缺少这些必需品就无法维持，祂说：\BibleQuote{这一切自会加给你们；但你们先该寻求天主的国和祂的义德。}因为祂用\Quote{先}这个字，表明后者应被寻求，不是时间上的先后，而是重要性的先后：前者是我们的善，后者是我们所需的必需品；但必需品是为了那善。\par

54. 例如，我们不应为吃饭而宣讲福音，而应为宣讲福音而吃饭；因为如果我们为吃饭而宣讲福音，我们就将福音视作低于食物；那样我们的善就在吃饭，而非在宣讲福音这一为我们所需之事。宗徒也禁止这一点，当他说，他自己本可合法且由主准许，使宣讲福音的人靠福音而生活，即从福音中获得现世生活的必需品；但他并未使用这权柄。因为有许多人渴望获得买卖福音的机会，宗徒想要切断他们这种机会，因此他甘愿亲手谋生。关于这些人，他在另一处说：\BibleQuote{我愿断绝那找机会者的机会。}即使他像其他善的宗徒那样，因主的准许而靠福音生活，他也不会将宣讲福音的目的置于那生活之中，反而会使福音成为他生活的目的；即，如上所述，他不会为得食物及一切必需品而宣讲福音；而是为这目的而取用这些东西，以便能实现另一目的，即自愿而非被迫地宣讲福音。因为他责备这种事时说：\BibleQuote{你们岂不知道，在圣殿供职的人吃圣殿中之物，在祭台旁服务的分享祭台之物？主也这样吩咐，叫宣讲福音的人靠福音生活。但我未曾用过这些权柄。}由此他表明这是准许，而非命令；否则他将被视为违抗主的命令。接着他说：\BibleQuote{我写这些话，并非为叫你们这样待我；因为我宁可死，也不叫人使我所夸耀的落了空。}他说这话，是因为他已决定，因那些寻找机会的人，要亲手谋生。\BibleQuote{因为我若宣讲福音，}他说，\BibleQuote{我无可夸耀：}即，我若为得这些东西而宣讲福音，或我怀着这目的宣讲福音，为得那些东西，并将福音的目的置于饮食和衣服中。但他为何无可夸耀？\BibleQuote{为不得已，}他说，\BibleQuote{就加在我身上}；即，我宣讲福音是因没有谋生手段，或我从宣讲永恒之事中获取现世之果；这样，宣讲福音就成为不得已之事，而非出于自由选择。\BibleQuote{因为我有祸了，}他说，\BibleQuote{若不宣讲福音！}但他应如何宣讲福音？显然，应将赏报置于福音本身和天主的国中：这样他就能自由地、而非被迫地宣讲福音。\BibleQuote{因为我若甘心做这事，}他说，\BibleQuote{就有赏报；若不甘心，管家的职分（\OriginalTerm{管家的职分}{dispensation (stewardship)}）却已托付给我。}若因现世生活所需之物的缺乏而被迫宣讲福音，别人将因我而获得福音的赏报——他们在我宣讲时爱福音本身；但我却没有，因我所爱的不是福音本身，而是它在现世事物中的代价。这是有罪的：一个人不是以儿子的身份，而是以被托付管家的仆人的身份服务福音；他仿佛支付别人的财物，自己却一无所获，除了食物——这些食物不是因他共享王国而赐予，而是从外面供给悲惨奴役的生活所需。虽然他在另一处也称自己为管家。因为当仆人被收纳入儿子之列时，他能忠实地将与其共享的财产分配给同伴，他已获得了共同继承人的份。但在此处，他说：\BibleQuote{若不甘心，管家的职分（\OriginalTerm{管家的职分}{dispensation (stewardship)}）却已托付给我}时，他愿这样的管家被理解为分派别人的财物，自己却一无所获。\par

55. 因此，凡是为另一样事物而寻求的东西，无疑低于其所要求的那样事物；所以那为你所求之物是第一位的，而非你为另一事物而求的该事物。因此，若我们为食物而寻求福音和天主的国，则我们将食物放在首位，天主的国在最后；以致若食物不缺，我们就不会寻求天主的国：这就是先求食物，再求天主的国。但若我们为此目的求食物，好能获得天主的国，我们就行了所记的：\BibleQuote{你们先寻求天主的国和祂的义德，这些一切自会加给你们。}\par

% structure: p0074 source=h2 target=subsection reason=relative-heading-rank; base=h2

\subsection{第十七章}

56. 对于那些先寻求天国及其义德的人，\Emph{即}那些将此置于一切之上，以至于为天国而寻求其他事物的人，不该再存有担忧，怕那些为天国而维持此世生命所必需的东西会缺少。因为祂在上文已说过：\Quote{你们的圣父知道你们需要这一切。}因此，当祂说：\Quote{你们先寻求天国及其义德}，祂并没有说，然后寻求这些东西（虽然它们是必需的），而是肯定\Quote{这一切都要加给你们}，\Emph{即}，若你们先寻求前者，这些将会随之而来，对你们毫无阻碍：免得你们在寻求这些东西时，从前者转离；或者免得你们设立两个目标，既为天国本身寻求天国，又寻求这些必需品；但这些东西应为那另一目标而来；这样它们就不会缺少给你们。因为你们不能事奉两个主人。但人若既寻求天国作为至善，又寻求这些现世事物，便是试图事奉两个主人。然而，他不能有纯洁的眼目，唯独事奉主天主，除非他将其他一切必需之物，都归于这唯一目标，\Emph{即}为天国而接受。但正如所有当兵的人都获得粮饷，所有宣讲福音的人也获得饮食和衣服。但并非所有当兵的都是为国家的福利，有些人是为了报酬；同样，并非所有事奉天主的人都是为了教会的利益，有些人是为了这些现世事物，即他们以粮饷形式所获得的；或者既为此也为彼。但上文已经说过：\Quote{你们不能事奉两个主人。}因此，我们应当以纯朴的心，唯独为天国的缘故而善待众人；我们行善时，不应只考虑现世的报酬，也不应考虑它连同天国：祂将一切现世事物归入明天的类别，说：\Quote{不要为明天忧虑。}因为明天一词只用于时间，其中未来接续过去。所以，当我们行善时，不要思念那暂时的，而要思念那永恒的；这样才是一件善而完美的工作。祂说：\Quote{明天自会为明天忧虑；}\Emph{即}，这样，当你们需要时，你们便会取用食物、饮料或衣服，也就是说，当需要本身开始催促你们时。因为这些东西会唾手可得，因为我们圣父知道我们需要这一切。因为祂说：\Quote{一天的苦难足够一天受的了。}\Emph{即}，需要本身催促我们取用这些东西就足够了。我想，正因如此，它被称为苦难，因为对我们而言，它是惩罚性的：它属于我们因犯罪而招致的脆弱和必死性。因此，不要给这现世需要的惩罚增添更重的负担，以免你们不仅遭受这些物质的匮乏，而且为满足这需要而去为天主当兵。\par

57. 然而，在使用这段经文时，我们必须格外警惕，以免当我们看到天主的某位僕人为自己或他所受托照料的人预备这些必需品时，就断定他违反了主的诫命，为明天而焦虑。因为主自己虽然天使服事祂，但为了以身作则，免得后来有人因见到祂的僕人置办这类必需品而跌倒，祂屈尊带有钱囊，以便从中支出所需的费用；如经上记载，出卖祂的\Person{犹达斯}是钱囊的保管者和窃贼。同样，宗徒\Person{保禄}似乎也曾为明天操心，他说：\BibleQuote{关于为圣徒的捐献，我怎样吩咐了迦拉达的众教会，你们也当照样行：每週的第一日，你们各人要照自己的进项抽出来留着，免得我来的时候才凑齐。及至我来了，你们所认可的人，我就打发他们，带着你们的捐资，送到耶路撒冷去。若我也该去，他们可以和我同去。但我从马其顿经过时，必到你们那里去；因为我从马其顿经过。或者我也可能留在你们那里，甚至和你们一同过冬，好让你们送我往我所要去的地方。我不愿在路上见你们；但我希望如果主准许，在你们那里住些时候。但我要在厄弗所住到五旬节。}在\WorkTitle{宗徒大事录}中也记载，因即将来临的饥荒，未来所需的食物已得预备。因为我们读到：\BibleQuote{那时，有几位先知从耶路撒冷下到安提约基雅。众人聚在一起时，有一位名叫阿加波的站起来，藉着圣神指明天下将有大饥荒；这事在\Person{喀劳狄}皇帝时果然发生了。于是门徒各按自己的力量，决定捐款救济住在犹太的弟兄们，他们就托\Person{巴尔纳伯}和\Person{扫禄}的手送去。}至于赠与他的必需品，当同一宗徒\Person{保禄}启航时，他载着这些，似乎不是仅够一天的食物。同一宗徒\Person{保禄}写道：\BibleQuote{从前偷窃的，不要再偷；总要劳力，亲手作正经事，就可有余，分给那缺少的人。}这在误解他的人看来，似乎不遵守主的诫命，即：\BibleQuote{你们看那天上的飞鸟，也不种，也不收，也不积蓄在仓里；}以及\BibleQuote{你想野地里的百合花，怎么长起来；它也不劳苦，也不纺线；}而他却吩咐那些人亲手劳作，好有余可以分给别人。从他常说自己亲手作工以免累及他人，以及记载他与\Person{阿桂拉}因同业而一同工作以维持生计，他好像并未效法天上的飞鸟和野地的百合花。从这些以及圣经中类似的经句，足以看出主并非不赞同人按常理看顾这些事，但若有人为这些事而当天主的兵士，以致他所行所注目的不是天主的国，而是得着这些事。\par

58. 因此，这整个诫命可归结为以下准则：即便在看顾这些事时，我们也当思念天主的国；但在为天主的国服务时，我们不当思念这些事。因为这样，即使有时这些东西缺乏（这是天主为操练我们而常允许的），它们非但不削弱我们的心志，反而在考验与试炼中使之坚固。因为经上说：\BibleQuote{我们也在患难中夸耀，知道患难生坚忍，坚忍生老练，老练生盼望；盼望不至于羞耻，因为天主的爱藉着所赐给我们的圣神，已浇灌在我们心里。}同一宗徒在提到自己的患难与劳苦时，说他不仅忍受了监牢、船难及许多类似的困苦，还有饥渴、寒冷与赤身露体。我们读到这些时，不要以为天主的应许动摇，以致宗徒在寻求天主的国和祂的义时，竟忍饥受渴、赤身露体，尽管我们被告诉说：\BibleQuote{你们要先求天主的国和祂的义，这一切都要加给你们了。}因为那位我们已完全交托自己，并从祂领受今生及来生应许的医生，知道这些事物正如辅助，祂何时赐予、何时收回，都按祂认为对我们有益的方式；祂管理并引领我们这些既需在此生得安慰与锻炼，又要在今生之后得安息与坚固的人。因为人若时常从他的负重的牲畜前拿走饲料，并非不关心它，反而是出于关心而这样做。\par

% structure: p0078 source=h2 target=subsection reason=relative-heading-rank; base=h2

\subsection{第十八章}

59. 既然这些事，无论是为将来作准备，或是留作备用，若没有理由去花费它们，其存心如何便是不确定的，因为可能以纯一的心，也可能以双重心去行，祂适时地在经上添加了这句话：\BibleQuote{你们不要判断，免得你们受判断。因为你们用什么判断来判断，你们也要受什么判断；你们用什么量器量，也要用什么量器量给你们。}\BibleRef{玛 7:1-2} 在这段经文中，我以为我们所受的教导无非是：对于那些行踪不定、存心可疑的行为，我们应作更好的推定。因为经上记着：\BibleQuote{凭他们的果实，你们能认出他们来。}\BibleRef{玛 7:16} 这句话是指那些显然不能以善意去行的事，例如放荡、亵渎、偷窃、醉酒，以及所有这类的事。对于这些事，我们被允许按宗徒的话去判断：\BibleQuote{审判教外的人，关我什么事？教内的人，岂不是你们审判的吗？}\BibleRef{格前 5:12} 但关于食物的种类，因为每一种人类的食物都可以用善意和纯一的心、没有贪欲的恶习而毫无区别地取用，同一位宗徒禁止那些吃肉喝酒的人被戒绝这类食物的人判断。他说：\BibleQuote{吃的人不可轻视不吃的人；不吃的人也不可判断吃的人。}\BibleRef{罗 14:3} 他又在那里说：\BibleQuote{你是谁，竟敢判断别人的仆人？他或站住或跌倒，自有他的主人在。}\BibleRef{罗 14:4} 因为关于那些可以用善意、纯一和崇高动机去做的事，虽然也可能用相反的动机去做，那些人，虽然他们不过是人，却想判断内心的隐秘，而只有天主是判断者。\par

60. 属于这一类的，还有他在另一处所说的话：\BibleQuote{所以，时候未到，你们什么也不要判断，只等主来，他要照出暗中的隐情，显明人心的意念；那时，各人要从天主那里得到称赞。}\BibleRef{格前 4:5} 因此，有些暧昧的行为，我们不知道它们存心如何，因为它们可能出于善意或恶意；对于这些，轻率地加以判断，尤其是为了定人之罪，是鲁莽的。将来必有时间审判这些事，那时主要\BibleQuote{照出暗中的隐情，显明内心的计谋。}\BibleRef{格前 4:5} 同一位宗徒在另一处又说：\BibleQuote{有些人的罪过，在受审判以前，就已显露出来；有些人的罪过，是随后才显露出来的。}\BibleRef{弟前 5:24} 他称那些显然可见的罪过为\Quote{显露的}，即明显知道其存心的；这些罪过先到审判前，因为若随后有审判，就不是轻率的。但那些隐藏的罪过后才显露，因为它们到了时候也不能隐藏。对善行也是如此。他接着说：\BibleQuote{同样，善行也有显露的；即使不是显露的，也不能隐藏。}\BibleRef{弟前 5:25} 所以，让我们判断那些明显的；至于那些隐藏的，让我们将判断留给天主：因为它们也不能隐藏，无论善恶，当显明的时候来到，它们必会显露。\par

61. 此外，我们应在两件事上提防轻率的判断：一是当某事的存心不确定时；二是当一个人目前明显是善或恶，但他将来会成为什么样的人还不确定时。因此，例如，有人因胃不舒服而不愿禁食，而你不相信这一点，却将其归咎于贪吃的恶习，你便是在轻率判断。同样，如果你知道某人显然贪吃醉酒，便加以责备，好像这人永远不能改正和改变，你也是轻率判断。所以，我们不要责备那些我们不知道存心如何的事；也不要这样责备那些明显的事，以至于对恢复正确的心境感到绝望。这样，我们就能避免这里所说的判断：\BibleQuote{你们不要判断，免得你们受判断。}\BibleRef{玛 7:1}\par

62. 但主所说的或许会引起困惑：\BibleQuote{因为你们用什么判断来判断，人也用什么判断来判断你们；你们用什么量器来量，也要用什么量器来量给你们。} 那么，如果我们对某件事作了轻率的判断，天主也会对我们作轻率的判断吗？或者，如果我们用不公义的量器来衡量，天主那里也有不公义的量器，好用来量给我们吗？（因为，我想，这里的\Quote{尺度}一词所指的正是判断本身。）天主绝不会轻率地判断，也不会以不公义的量器报应任何人；但经文之所以如此表述，是因为你用来惩罚他人的那种轻率本身，必然也会惩罚你自己。除非，也许有人会想象，不义在某种程度上伤害了它所针对的人，而丝毫不伤害发出不义的人；但事实恰恰相反，它常常不伤害受伤害的人，而必然伤害施加伤害的人。迫害者的不义对殉道者造成了什么伤害呢？一点没有；但对迫害者自己却伤害极大。因为他们中的有些人虽然后来从错误中回头，但在他们行迫害的时候，他们的恶行是蒙蔽着他们的。同样，轻率的判断常常不损害被轻率判断的人；但轻率本身必然损害作出轻率判断的人。按照这一原则，我也理解那句话：\BibleQuote{凡持剑的，必死于剑下。} 因为有多少人持剑，却没有死在剑下，\Person{伯多禄}本人就是一例！但恐怕有人认为，他因罪过得赦免而逃脱了这种刑罚（尽管再没有比认为没有落在\Person{伯多禄}身上的剑刑，比实际落在他身上的十字架之刑更重，更荒谬的了）；然而，他们又会怎样说与主同钉的凶犯呢？因为那获得赦免的，是在被钉之后才获得的，而另一个根本没有获得赦免？或者，也许他们杀了所有被他们杀死的人，因此他们自己也该受同样的刑罚？这样想是很可笑的。因为\BibleQuote{凡持剑的，必死于剑下}这句话，除了指灵魂因它所犯的任何罪而死之外，还有什么意思呢？\par

% structure: p0083 source=h2 target=subsection reason=relative-heading-rank; base=h2

\subsection{第十九章}

63. 主正在这一段中告诫我们关于轻率和不当的判断——因为祂愿意我们做一切事时，都怀有一颗纯洁的心，单单注视天主；而且，对于许多事情，其动机是不确定的，对之轻易判断是轻率的；此外，那些对那些不确定之事尤其轻率判断并轻易指责的人，是那些宁愿指责和谴责，而不愿改正和改善的人，这缺点或出于骄傲，或出于嫉妒；因此祂接着说了这句话：\BibleQuote{为什么你看见你兄弟眼中的木屑，却不想到你自己眼中的大梁呢？} 这样，假如他（你的兄弟）在忿怒上犯了过犯，你却在仇恨上加以指责；两者之间的差别，就像木屑和大梁的差别一样。因为仇恨是积久的忿怒，它仿佛由于持续长久而获得了如此大的力量，以致可以合理地称之为大梁。现在，你可能对一个人生气，却仍希望他从错误中回头；但如果你恨一个人，你就不可能希望他悔改。\par

64. \BibleQuote{或者，你怎能对你的兄弟说：让我去掉你眼中的木屑；而看哪，你眼中却有一根大梁？假善人哪！先从你眼中去掉大梁，然后你才看得清楚，去掉你兄弟眼中的木屑。} \Emph{即}，先去掉你内心的仇恨，然后，而不是在此之前，你才能改正你所爱的人。祂说得很好：\Quote{假善人。} 因为抱怨恶习是善人、好人的职责；当恶人这样做时，他们是在扮演不属于他们的角色；就像假善人一样，他们用面具掩盖自己的本相，而用面具显示他们所不是的样子。因此，在\Quote{假善人}这个称呼下，你要理解为伪善者。事实上，有一类伪善者非常需要提防，而且令人厌烦，他们出于仇恨和恶意而抱怨各种过犯，却又希望显得是顾问。因此，我们必须虔诚而又谨慎地留意，以便当必要时我们不得不指责或规劝某人时，我们首先反思该过犯是否是我们从未有过的，或是我们现在已经摆脱了的；如果我们从未有过，让我们反思我们是人，本来可能有的；但如果我们有过，而现在已经摆脱了，那么共同的软弱应触动记忆，使我们在指责或施行规劝时，心里充满的不是仇恨，而是怜悯：这样，无论此举是否能使他（我们为之而这样做的人）悔改或使他更堕落（因为结果不确定），我们至少因眼睛的纯洁而无忧无虑。然而，如果我们反思后发现自己也陷于他（我们准备指责的人）所有的同样过犯中，那么我们就不要指责也不规劝；但仍要深深哀悼这事，并邀请他不要服从我们，而是与我们共同奋斗。\par

65. 关于宗徒所说的：\BibleQuote{对犹太人，我就成为犹太人，为要赢得犹太人；对在法律以下的人，我就成为在法律以下的人（其实我不在法律以下），为要赢得在法律以下的人；对没有法律的人，我就成为没有法律的人（其实我对天主没有法律，但对基督却在法律以下），为要赢得没有法律的人。对软弱的人，我就成为软弱的，为要赢得软弱的人：对一切人，我就成为一切，为要赢得一切人。}——他当然不是像某些人所想的那样，以虚伪的方式行事，以便他们那可恶的虚伪能借如此伟大榜样的权威得以加强；而是出于爱德行事，在爱德的影响下，他将所愿帮助之人的软弱视为自己的软弱。为此，他也预先立下了根基，他说：\BibleQuote{我虽是自由的，无人管辖，却使自己成为众人的奴仆，为要赢得更多的人。}为让你明白这是出于爱德，而非虚伪，在爱德的影响下，我们同情软弱的人，如同自己软弱一样，他在另一处告诫我们说：\BibleQuote{弟兄们，你们蒙召是为得自由，但不要以自由作为放纵肉欲的机会，却要以爱德彼此服事。}除非每个人都把别人的软弱当作自己的软弱，从而心平气和地忍受，直到他所关心的人摆脱软弱，否则这是无法做到的。\par

66. 因此，只有在很少情况下，且出于极大的需要，才应施以斥责；然而，即使在斥责中，我们也应竭力追求事奉天主，而非事奉我们自己。因为天主，而非任何其他事物，才是终点：因此，我们不可心怀二意，而应从自己的眼中除去嫉妒、恶意或虚伪的木梁，这样我们才能看清如何取出弟兄眼中的木屑。因为我们将以鸽子的眼睛——就是那属于基督净配的眼睛——去看，这净配就是天主为自己所拣选的荣耀的教会，没有瑕疵或皱纹，即纯洁而无诡诈的。\par

% structure: p0088 source=h2 target=subsection reason=relative-heading-rank; base=h2

\subsection{第二十章}

67. 然而，由于\Quote{诡诈}一词可能会误导那些渴望遵守天主诫命的人，使他们认为有时隐瞒真理是错误的，正如有时说谎是错误的；而且，由于这种方式——向那些无法承受的人揭示秘密——可能会造成比完全隐瞒更大的伤害，因此主非常恰当地补充说：\BibleQuote{不要把圣物给狗，也不要把你们的珍珠丢在猪前，免得它们用脚践踏，又转过来撕裂你们。}因为主自己虽然从未说谎，却表明祂隐瞒了一些真理，当祂说：\BibleQuote{我还有许多事要告诉你们，但你们现在不能担负。}宗徒保禄也说：\BibleQuote{弟兄们，我从前对你们说话，不能把你们当作属神的人，只能当作属血肉的人，当作在基督里的婴孩。我给你们喝的是奶，不是饭；因为那时你们还不能吃，就是现在还是不能，因为你们仍是属血肉的。}\par

68. 现在，在这条禁止我们将圣物给狗、将珍珠丢在猪前的诫命里，我们必须仔细探究：何为圣物？何为珍珠？何为狗？何为猪？圣物是一种不可侵犯、不可败坏的东西；企图或希望犯下这种罪行本身就被视为有罪，尽管那圣物在本性上仍然不可侵犯、不可败坏。所谓珍珠，是指一切我们应当高度重视的属灵事物，因为它们藏在隐秘之处，仿佛从深处取出，并且包裹在寓意的外衣里，如同打开的贝壳。因此，我们可以合理地理解，同一事物既可称为圣物，也可称为珍珠；但它之所以称为圣物，是因为它不应被败坏；称为珍珠，是因为它不应被轻视。然而，每个人都试图败坏他不愿保持完整的东西；但他轻视他认为无价值、并视之为不如自己的东西；因此，被轻视的东西就被说成是被践踏。因此，既然狗扑向东西是为了撕碎，并且不让被撕碎的东西保持原样，所以祂说：\BibleQuote{不要把圣物给狗。}因为尽管它不能被撕碎和败坏，仍保持完整和不可侵犯，但我们仍必须思想那些以苦毒和极其敌对的精神抵抗、并尽其所能企图毁灭真理之人的意愿。然而，猪虽然不像狗那样用牙齿扑向物体，却因鲁莽地践踏而玷污它：\BibleQuote{所以不要把你们的珍珠丢在猪前，免得它们用脚践踏，又转过来撕裂你们。}因此，我们可以恰当地理解，狗用来指称真理的攻击者，猪用来指称真理的藐视者。\par

69. 但当祂说：\Quote{他们转过来撕咬你们}时，祂并不是说他们撕咬珍珠本身。因为他们践踏珍珠，当转身想听更多时，却撕咬那刚把珍珠抛在他们面前而他们已践踏了的人。因为你很难发现一个践踏脚下珍珠的人有什么乐趣，也就是说，轻视神圣事物的人，其发现是费了很大劳力的。但至于教导这类人的人，我看不出他怎能逃脱他们的愤怒和盛怒而被撕碎。此外，两种动物都是不洁的，狗和猪都一样。因此我们必须警惕，不能向不接受的人开启任何事：因为更好的是他寻求隐藏的事，而不是攻击或轻视公开的事。事实上，找不到其他原因为何他们不接受那些明显而重要的事，只除了仇恨和轻视，前者使他们被称为狗，后者使他们被称为猪。所有这些不洁都是由贪恋暂时之物产生的，也就是贪恋此世，我们必须弃绝此世，才能成为洁净。因此，渴望拥有纯洁心灵的人，不应自以为有错，如果他向不能接受的人隐瞒了什么。也不能由此认为说谎是允许的：因为隐藏真理并不等于说出谎言。因此，首先要采取步骤，除去妨碍他接受的障碍；因为如果是不洁导致他不接受，就应用言语或行为尽可能洁净他。\par

70. 再者，当我们的主所说的一些话，在场许多人不接受，反而抗拒或轻视时，不应认为祂将圣物给了狗，或将珍珠丢在猪前：因为祂不是给那些不能接受的人，而是给那些能接受且在场的人；祂不应因他人的不洁而忽略这些人。当试探者向祂提问，祂回答他们，使他们无话可说时，虽然他们可能因自己的毒害而憔悴，而非因祂的食粮得饱足，但其他能接受祂教导的人，却因这些人制造的时机而听到许多有益之事。我说这些，是唯恐有人也许无法回答提问者时，会以为自己有借口，说他不愿将圣物给狗，或将珍珠丢在猪前。因为知道该如何回答的人，应该为了他人而这样做，否则若他们以为提出的问题无法回答，就会心生绝望；而这是指那些有益且属于救恩训导的事。因为闲散之人所探究的许多事是无益、虚妄，且常是有害的，但关于这些事，仍需说些什么；而这一点本身应被揭示和解释，即为何不该探究这些事。因此，关于有益的事，我们有时应回答所问的问题：正如主所做的，当撒杜塞人问祂关于那嫁了七个丈夫的女人，复活后她是哪一个的妻子时，祂回答说：在复活时，他们也不娶也不嫁，而是像天上的天使一样。但有时，提问者应当被反问另一个问题，通过回答它，他自己就能回答所问的事；但如果他拒绝陈述，在场的人就会觉得公平，如果他不再听到他所问的事。因为那些试探祂的人，问是否该纳税时，被问了另一个问题：他们自己拿出的钱币是谁的像；因为他们说出了所问的，即钱币是凯撒的像，他们就在某种程度上自己回答了所问主的问题；因此，主从他们的回答得出这个结论：\Quote{那么，凯撒的归凯撒，天主的归天主。}然而，当司祭长和民间长老问祂凭什权柄做这些事时，祂反问他们关于若翰的洗礼；当他们不愿说出那对他们不利的陈述，却又不敢因旁听者而对若翰说不好的话时，祂说：\Quote{我也不告诉你们我凭什么权柄做这些事。}这个拒绝在旁观者看来极为合理。因为他们声称不知道他们实际知道却不愿说的事。事实上，他们既然希望得到所问的答案，就应该先对别人做他们要求别人对他们做的事；如果他们这样做了，他们就一定会自己回答自己。因为他们自己曾派人到若翰那里，问他是谁；或者更确切地说，他们自己是司祭和肋未人，被派去，以为他就是基督，但他声称不是，并为主作了见证；如果他们愿意承认这个见证，他们就会教导自己，作为基督的祂凭什么权柄做这些事；而他们假装不知道，问祂，是为了寻找诽谤的途径。\par

% structure: p0093 source=h2 target=subsection reason=relative-heading-rank; base=h2

\subsection{第二十一章}

71. 因此，既然已命令不可把圣物给狗，也不可把珍珠丢在猪前，听者可能会基于自身的无知和软弱而提出异议，并听到给他这条命令：他不应给出他自己尚未领受的东西——（我再说）可能会提出异议说：你禁止我把什么圣物给狗，禁止我把什么珍珠丢在猪前，而我尚未看见我拥有这些呢？祂非常适时地补充了这句话：\BibleQuote{你们求，必要给你们；你们找，必要找着；你们敲，必要给你们开。因为凡求的，就必得到；找的，就必找着；敲的，就必给他开。}所求的，是指通过祈求获得心智的健康和力量，以便能履行所命令的职责；而寻找，则是指发现真理。因为福乐生命归结于行动和知识，行动需要力量的供应，默观则渴望事物被阐明：因此，前者是应求的，后者是应寻找的；以使前者被给予，后者被找着。但今生的知识更属于路途，而非拥有本身；但谁若找到了正路，必会到达那拥有本身，而拥有本身则给敲者开启。\par

72. 因此，为使这三件事——即求、找、敲——得以明了，我们举一个例子，假定有人肢体虚弱，不能行走：首先，他必须被治愈并增强力量，以便能行走；祂所用的\BibleQuote{求}一词即指此。但他如今能行走，甚至奔跑，若走上歧途，又有什么益处呢？因此第二件事是，他应找到通往他所欲抵达之处的道路；并且他沿着那条道路行走，抵达了他愿居住的地方，若发现门关闭，则他能行走、能行走并到达都无济于事，除非给他开门；因此，所用的\BibleQuote{敲}一词即指此。\par

73. 此外，那不欺骗人而应许的祂，已给予了并正在给予极大的希望：因为祂说：\BibleQuote{凡求的，就必得到；找的，就必找着；敲的，就必给他开。}因此，需要坚忍，以便我们得到所求的，找到所找的，所敲的门得开启。现在，就如祂论到天空的飞鸟和田间的百合花，使我们不致对衣食供应绝望，好让我们的希望从较小的事物上升到更大的事物；同样，在本段中，祂说：\BibleQuote{或者，你们中间有哪个人，儿子向他求饼，反而给他石头呢？或者求鱼，反而给他蛇呢？你们纵然不善，尚且知道把好的东西给你们的儿女，何况你们在天之父，岂不更把好东西赐给求祂的人吗？}不善的人如何能给予好东西呢？祂称那些尚是今世之恋者和罪人为不善。事实上，这些好东西应按照他们的感受称为好的，因为他们视这些为好的。虽然在万物的本性中，这些也是好的，但只是暂时的，属于这软弱的人生；而任何不善者给予它们，并非出于自己的所有；因为大地和其中的万物都属于上主，祂创造了天、地、海及其中所有的一切。因此，我们有极充分的理由盼望天主在我们求祂时，会将好东西赐给我们，而我们不会受骗，以致当我们求祂时得到一样而不是另一样；因为我们尽管不善，尚且知道给予所请求的东西？我们并不欺骗我们的儿女；我们给予的任何好东西，也并非出于我们自己，而是出于祂的。\par

% structure: p0097 source=h2 target=subsection reason=relative-heading-rank; base=h2

\subsection{第二十二章}

74. 此外，行走智慧之路的某种力量和活力寓于良好的品行中，这品行被扩展以达到内心的纯洁和纯一——祂现在已讲了许久这个主题，并这样总结：\BibleQuote{所以，凡你们愿意人给你们做的，你们也要照样给人做：法律和先知即在于此。}在希腊文抄本中，我们见到该段是这样写的：\BibleQuote{所以，凡你们愿意人给你们做的，你们也要照样给人做。}但我认为拉丁文译者加上了 \Quote{善} 字以使句子更清晰。因为有人会想到，如果有人希望别人对他做某件坏事，并将此句应用于此——例如，若有人希望被挑战而过度饮酒，并在杯间喝醉，而他先对那个他希望如此待他的人做了同样的事——那么认为他实现了这一句就荒谬可笑了。因此，由于他们受此考虑的影响，我认为，便加了一个词以使问题清楚；以致在 \BibleQuote{所以，凡你们愿意人给你们做的} 这一句中插入了 \Quote{善} 字。但若希腊文抄本中没有这个词，那么它们也应被修正：但谁敢这样做呢？因此，应当理解，即使没有加上这个词，这一句话也是完整且完全完善的。因为所用的 \Quote{凡你们愿意的} 这句话，应当理解为并非按日常和随便的意思，而是按严格的意义。因为只有在善的事上才有真正的意愿：在恶的和邪恶的事上，严格地说，是欲念，而不是意愿。并非圣经总是按严格意义说话；但在必要的地方，它们如此严格地保持一个词的意义，以致不允许有其他理解。\par

75. 此外，这条诫命似乎指向爱近人的诫命，而非也指向爱天主的诫命，因为在另一处祂说，有两条诫命，\BibleQuote{一切法律和先知都系于此。} 因为如果祂曾说：凡你们愿意人给你们做的，你们也要照样做；在这一句话中，祂就会包含这两条诫命：因为立刻就会有人说，每个人都希望自己既被天主爱，也被人们爱；因此，当这条诫命赐给他，即他愿意别人对他做的，他也要照样做，那当然就等同于他应当爱天主和爱人的诫命。但当更明确地对人说话时，\BibleQuote{所以，凡你们愿意人给你们做的，你们也要照样给人做，} 这似乎无非就是：\BibleQuote{你应当爱你的近人如你自己。} 但我们必须仔细注意祂在这里附加的话：\BibleQuote{因为这就是法律和先知。} 现在，在这两条诫命的情况下，祂不仅说：法律和先知系于此；而且还加了\BibleQuote{一切法律和先知}，这与全部的预言相同；而祂在这里没有作同样的附加，为另一条诫命（涉及爱天主）留了位置。那么，既然祂是在遵循关于单一之心的诫命，并且恐怕有人会对那些心不可见的人（即对人）怀有双心，就必须给出关于那件事本身的诫命。因为几乎没有人愿意任何存有双心的人与自己往来。但没有人能以单一之心将任何东西给予邻人，除非他如此给予，不期望从他那里得到任何现世的利益，并以我们在讨论单纯的眼时已充分讨论过的意向行事。\par

76. 因此，眼被清洁并成为单纯，就会适应并适合注视和默观它自身的内在之光。因为所说的眼是心灵的眼。如今，这样的人拥有这样的眼：他为了使自己的行为真正良善，不以讨人喜悦作为他善行的目标；即使他碰巧使人喜悦，他也使这一点倾向于他们的得救和天主的荣耀，而非自己的虚空夸耀；他也不为了获得维持现世生活所需之物而做任何有益于邻人得救的善事；他也不轻率地谴责一个人在那种不明显以何种意向和愿望所做的事中的意向和愿望；他向人施行的任何仁慈，都怀着希望别人向自己施行的同样意向去施行，即不期望从他那里得到任何现世的利益：这样，心就会是单纯而纯洁的，在其中寻求天主。因此，\BibleQuote{心里洁净的人是有福的，因为他们要看见天主。}\par

% structure: p0101 source=h2 target=subsection reason=relative-heading-rank; base=h2

\subsection{第二十三章}

77. 但因为这属于少数人，祂现在开始谈论寻求和拥有智慧，这是生命树；而且，在寻求和拥有（即默观）这智慧时，这样的眼通过前面所有内容被引领到一个地步，在那里现在可以看到狭窄的路和窄门。因此，祂接着说：\BibleQuote{你们要从窄门进去，因为宽门和大路导入丧亡，但有许多人从那里进去；那导入生命的门是多么窄，路是多么狭！找到它的人的确很少。} 祂这样说，并不是因为主的轭是粗糙的，或祂的担子是沉重的；而是因为很少有人愿意结束他们的劳苦，不太相信那位呼喊者的话：\BibleQuote{凡劳苦和负重担的，你们都到我跟前来，我要使你们安息。你们背起我的轭，跟我学吧！因为我是良善心谦的：这样你们必要找得你们灵魂的安息，因为我的轭是柔和的，我的担子是轻松的。}（此外，我们面前的这篇讲道正是以心谦和心善为起点的）：这个柔和的轭和轻松的担子，许多人藐视，只有少数人接受；因此那导入生命的道路变得狭窄，进入的门变得窄。\par

% structure: p0103 source=h2 target=subsection reason=relative-heading-rank; base=h2

\subsection{第二十四章}

78. 因此，在这里要特别谨防那些许诺自己所不拥有的智慧和真理知识的人，例如异端者，他们常常以人数稀少自夸。因此，当祂说找到窄门和窄路的人很少时，免得他们假借人数稀少而冒充自己，祂立刻附加说：\BibleQuote{你们要提防假先知！他们来到你们跟前，外面披着羊皮，里面却是凶残的豺狼。} 但这类人欺骗不了单纯的眼，这眼知道如何凭果实辨别树。因为祂说：\BibleQuote{你们可凭他们的果实辨别他们。} 然后祂加上比喻：\BibleQuote{人岂能从荆棘上采葡萄，或者从蒺藜上采无花果吗？同样，凡是好树都结好果子，而坏树都结坏果子。好树不能结坏果子，坏树也不能结好果子。凡不结好果子的树，都要砍倒，投入火中。所以，你们可凭他们的果实辨别他们。}\par

79. 在解释这段经文时，我们必须非常警惕那些人的错误，他们根据这两棵树判断存在两种本原，一种属于天主，另一种既不属於天主，也不出自祂。这种错误已在其他著作中有过非常详尽的讨论，如果还嫌不够，日后还会再讨论。但目前我们只需说明，眼前的这两棵树并不能支持他们的观点。首先，因为基督显然是在谈论人，凡是阅读前后文的人都会对他们的盲目感到惊讶。其次，他们专注于这句话：\BibleQuote{好树不能结坏果子，坏树也不能结好果子，}因此认为，恶灵魂不可能变成更好的，善灵魂也不可能变成更坏的；仿佛是说，好树不能变成坏树，坏树也不能变成好树。但经文说的是：\BibleQuote{好树不能结坏果子，坏树也不能结好果子。}树当然是指灵魂本身，也就是人本身，而果子是指人的行为；因此，恶人不能行善事，善人也不能行恶事。所以，如果恶人想要行善事，就让他先成为善人。主自己在另一处更明确地说：\BibleQuote{或者使树好，或者使树坏。}但如果祂是用这两棵树来比喻两种本原，就不会说使，因为世人中谁能创造一种本原呢？此外，在另一处提到这两棵树时，祂接着又说：\BibleQuote{你们这些伪善人，你们既然邪恶，怎能说出好话来呢？}因此，只要一个人是邪恶的，他就不能结出好果子；因为如果他结出好果子，就不再是邪恶的了。同样，最真实的话是：雪不能是暖的；因为当它开始变暖时，我们就不再称之为雪，而是水。因此，曾经是雪的东西可以不再如此；但雪不能变成暖的。同样，曾经邪恶的人可以不再邪恶；但恶人不能行善。虽然恶人有时也产生益处，但这并非出于他自己，而是通过他，凭天主眷顾的安排而成就的；例如，关于法利塞人，祂说：\BibleQuote{他们对你们说什么，你们就做什么；但不要效法他们的行为。}他们说了好的事情，并且他们所说的被有益地聆听和遵行，这本身并不属于他们；因为祂说：\BibleQuote{他们坐在梅瑟的座位上。}因此，当他们藉天主眷顾宣讲天主的法律时，他们对听众能有裨益，尽管对自己没有裨益。关于这样的人，先知在另一处说：\BibleQuote{他们播种了麦子，却要收割荆棘。}因为他们教导善事，却行恶事。因此，那些听他们说话并遵行他们所说之言的人，并非从荆棘中收取葡萄，而是透过荆棘从葡萄树上收取葡萄：就像有人把手伸过篱笆，或者至少从缠在篱笆上的葡萄藤上摘下一颗葡萄，这并非荆棘的果子，而是葡萄的果子。\par

80. 确实有一个非常恰当的问题：祂希望我们留意哪些果子，以便认出树来？因为许多人把某些属于羊皮的东西算作果子，并因此被狼欺骗；例如，禁食、祈祷或哀矜；但除非这些事连伪善者也能做，否则祂上面就不会说：\BibleQuote{你们要小心，不要在人前行你们的义，为叫他们看见。}在说了这句话之后，祂接着谈论这三件事：哀矜、祈祷、禁食。因为许多人慷慨哀矜穷人，并非出于怜悯，而是出于虚荣；许多人祈祷，或看似祈祷，却不注视天主，只愿取悦于人；许多人禁食，在那些认为禁食困难并由他们视之值得尊敬的人面前，做出惊人的克己表现；他们用这类诡计捕捉人，一面提出某物以欺骗，一面伸出另一物以捕食或杀害那些不能看见羊皮下之狼的人。因此，这些并非祂告诫我们认树的果子。因为当这些事以善意的真诚而行时，它们是羊的合宜外衣；但当它们以邪恶的欺骗而行时，它们所遮掩的不过是狼。然而，羊不应因此憎恨自己的外衣，因为狼常常藏匿其中。\par

81. 借着发现哪些果实，我们可以认出坏树，宗徒告诉我们说：\BibleQuote{本性的作为是显而易见的：即淫乱、不洁、放荡、崇拜偶像、行邪法、仇恨、纷争、嫉妒、忿怒、争吵、不睦、分党、异端、妒恨、凶杀、醉酒、宴乐，以及类似的事。我先前告诉过你们，现在又告诉你们：做这种事的人，绝不能承受天主的国。} 而借着哪些果实，我们可以认出好树，同一位宗徒接着告诉我们说：\BibleQuote{圣神的果实却是仁爱、喜乐、平安、忍耐、良善、温和、信德、柔和、节制。} 诚然，应当知道，此处\Quote{\OriginalTerm{喜乐}{gaudium}}一词是严格按照正确意义使用的；因为严格地说，恶人不能称为喜乐，而是夸张的欢跃：正如我们上面所说，严格意义上的\Quote{\OriginalTerm{意愿}{voluntas}}恶人没有，正是在这种意义上说：\BibleQuote{凡你们愿意别人给你们做的，你们也要照样给人做。} 按照这个词的严格意义，喜乐仅限于善人，先知也说话：\BibleQuote{喜乐不属于恶人，上主说。} 同样，\Quote{\OriginalTerm{信德}{fides}}也不是指任何一种信德，而是真信德；这里列出的其他事物，在恶人和欺骗者身上也有某些相似之处；因此，除非人有一双纯净和单一的眼睛来识别这些东西，否则就会完全被误导。因此，最恰当的顺序是，先讨论眼睛的洁净，然后提到需要防范的事情。\par

% structure: p0108 source=h2 target=subsection reason=relative-heading-rank; base=h2

\subsection{第二十五章}

82. 但是，无论一个人的眼睛多么纯净，也就是说，无论他生活的方式多么纯一和真诚，他仍不能看透别人的心：凡在行为或言语上未能彰显的事，就要通过试炼来揭示。试炼是双重的：或出于获得某种暂时利益的希望，或出于失去它的恐惧。尤其要警惕的是，当我们追求智慧时（智慧只能在基督内找到，在祂内蕴藏着智慧和知识的一切宝藏）——我说，我们一定要警惕，免得在基督的圣名之下，被异端者或任何智力有缺陷、热爱此世的人所欺骗。因此，祂加上一项警告，说：\BibleQuote{不是凡向我说\InnerQuote{主啊，主啊}的，就能进天国；而是那承行我在天之父旨意的，才能进天国。} 免得我们以为仅仅向主说\Quote{主啊，主啊}就属于那些果实，从而认为他是好树。但这些才是果实：承行在天之父的旨意，而祂亲自屈尊就卑，以此作为榜样。\par

83. 但完全可能提出一个问题：宗徒所说的话如何与这个句子协调一致，宗徒说：\BibleQuote{没有一个受天主圣神感动的人说：\InnerQuote{耶稣是可咒骂的}；除非受圣神感动，也没有一个能说：\InnerQuote{耶稣是主}。} 因为既不能说，凡有圣神的人，若坚持到底，就不能进天国；也不能说，那些说\Quote{主啊，主啊}而不进天国的人，就没有圣神。那么，为什么没有一个人说\Quote{耶稣是主，除非受圣神感动}呢？除非是宗徒在此处严格而正确地使用了\Quote{\OriginalTerm{说}{dicere}}这个词，表示说话者的意愿和理解？但主却普遍使用这个词：\BibleQuote{不是凡向我说\InnerQuote{主啊，主啊}的，就能进天国。} 因为那既不愿意也不理解自己所说的人，似乎也在说；但只有用声音表达自己的意愿和心志的人，才算是真正地说。正如刚才，在圣神的果实中被称为\Quote{\OriginalTerm{喜乐}{gaudium}}的，是在严格正确意义上称呼的，而不是像同一位宗徒在别处所说的\Quote{不以不义为乐}，好像有人会以不义为乐似的；因为那种心智混乱、喧哗欢跃的状态并不是喜乐；只有善人才拥有真正的喜乐。因此，那些既没有从理智上领会，也没有以意志的慎重同意参与他们所说出的话，仅仅用声音说出的人，似乎也在说；主在这种意义上说：\BibleQuote{不是凡向我说\InnerQuote{主啊，主啊}的，就能进天国。} 但真正且适当地说出的人，是那些其话语真正代表他们的意愿和意图的人；正是根据这种意义，宗徒说：\BibleQuote{除非受圣神感动，没有一个人能说：\InnerQuote{耶稣是主。}}\par

84. 此外，这与当前问题尤其相关：在追求真理的默观时，我们不仅不应被那些有基督之名却无基督之行的人藉基督之名所欺骗，而且也不应被某些行为和奇迹所欺骗。因为主曾为不信之人的缘故行过同类奇迹，祂警告我们不要被此类事物蒙蔽，以为眼见奇迹之处便有不可见的智慧。因此祂附上了这句话：\BibleQuote{在那一天，许多人要对我说：主啊，主啊，我们不是因你的名字说过预言，因你的名字驱过魔鬼，因你的名字行过许多异能吗？那时我要向他们声明说：我从来不认识你们；你们这些作恶的人，离开我吧！}因此，祂只承认那行义的人。因为祂也禁止祂自己的门徒因诸如此类的事欢乐，即因魔鬼屈服于他们而欢乐：\BibleQuote{却要欢乐，}祂说，\BibleQuote{因为你们的名字已经登记在天上了；}我想，是在那属天的\Place{耶路撒冷}城中，只有义人和圣者将在那里为王。\BibleQuote{你们岂不知道，}宗徒说，\BibleQuote{不义的人不得继承天主的国吗？}\par

85. 但或许有人会说，不义的人不能行那些可见的奇迹，反而宁愿相信那些将被发现说这样话的人是在撒谎：\BibleQuote{我们因你的名字说过预言，因你的名字驱过魔鬼，因你的名字行过许多异能。}因此，让他读一读埃及的术士们行了何等大事，他们曾抵抗天主的仆人梅瑟；或者，若他不愿读这个，因为那些事不是因基督的名行的，那么让他读一读主亲自论及假先知的话：\BibleQuote{那时，若有人对你们说：看，基督在这里，或在那里；你们不要相信。因为将有假基督、假先知兴起，行大奇事和异兆，以致连被选的人也要被迷惑。看，我预先告诉了你们。}\par

86. 因此，何等需要纯洁而单一的眼睛，好能寻得智慧的道路！因为这道路面对着邪恶乖戾之人如此众多欺骗与错误的喧嚣，逃避这一切，正是抵达最确定的平安和智慧不可动摇的稳固！因为十分可怕的是，若人因热衷争吵与辩论，便看不见那只有少数人能看见的：反对者的扰乱是微小的，除非人自己也扰乱自己。宗徒的话也指向这一点：\BibleQuote{主的仆人不应争辩，却应对众人温和，善于教导，忍耐，以温良开导那些有不同想法的人；或许天主赐予他们悔改，得以认识真理。}\BibleQuote{有福的，}因此，\BibleQuote{是缔造和平的人：因为他们要称为天主的子女。}\par

87. 因此，我们必须特别注意这整个讲道的结论是何等严肃庄重地被引入：\BibleQuote{所以，凡听了我这些话而实行的，就好像一个聪明人，把自己的房屋建在磐石上。}因为没有人能确认他所听到或理解的，除非藉着实行。如果基督就是磐石，正如许多圣经的见证所宣告的那样，那么实行从他那里所听到的人，就是建立在基督内。\BibleQuote{雨淋，水冲，风吹，袭击那座房屋，它却没有倒塌，因为它是建立在磐石上。}因此，这样的人不惧怕任何阴郁的迷信（因为雨若被用来比喻任何坏事，还能指别的什么呢？），也不惧怕人群的骚乱（我想这被比作风），也不惧怕此生的河流，如同在世俗的情欲中漫过地面。因为正是那个被顺境所诱惑的人，会被从这三者产生的逆境所摧毁；而那个将房屋建立在磐石上的人，一样也不惧怕，\Emph{即}他不仅听，而且实行主的命令。而那听了却不实行的人，与这一切都危险地接近，因为他没有稳固的根基；因着听而不行，他建造的是废墟。因为祂接着说：\BibleQuote{凡听了我这些话而不实行的，就好像一个愚昧人，把自己的房屋建在沙土上：雨淋，水冲，风吹，袭击那座房屋，它就倒塌了，且倒塌得很惨。耶稣讲完了这些话，群众都惊奇他的教训，因为他教训他们，像有权威的人，不像他们的经师。}这就是我在之前所说的，\WorkTitle{圣咏}中先知所指的，他说：\BibleQuote{我要因他而行事自信。上主的话语是纯净的言语，如同在泥土炉中炼过七次的银子。}从这个数字，我受提醒，也要把这些诫命追溯到他在这篇讲道开始时谈及真福时所置放的七个句子；以及追溯到依撒意亚先知所提及的圣神的七种运作；但在这其中，是目前的次序还是其他次序需要考虑，我们从主所听见的，若要建立在磐石上，就必须去实行。\par

\SourceRecord{Translated by William Findlay.；Nicene and Post-Nicene Fathers, First Series；Vol. 6.；Edited by Philip Schaff.；Buffalo, NY: Christian Literature Publishing Co.,；1888.；<http://www.newadvent.org/fathers/16012.htm>.}
